% Les types de mariages — Allemand 1ère A — Cours signé OnBuch+
% Programme officiel MINESEC. Compiler avec : tectonic ALA02-les-types-de-mariages.tex
\newcommand{\DOCMATIERE}{Allemand}
\newcommand{\DOCNIVEAU}{1ère A}
\newcommand{\DOCMODULE}{Chapitre 1 : Vie familiale et sociale}
\newcommand{\DOCLECON}{ALA02}
\newcommand{\DOCTITRE}{Les types de mariages}
% OnBuch+ — préambule « cours signés » ENRICHI (compilé avec tectonic / XeLaTeX)
% Système visuel « Pop » d'OnBuch+ : fond crème (jamais blanc), bordures encre
% épaisses, ombres « dures » décalées, titres Archivo Black en capitales.
% Version MATHÉMATIQUES : graphes (pgfplots/tikz), boîtes pédagogiques
% (définition, propriété, méthode, attention, le savais-tu, exercice résolu,
% exercice, TP, bilan), page de garde et signature OnBuch+ ancrée en bas.
%
% Macros attendues dans main.tex AVANT \input{preamble} :
%   \DOCMATIERE  \DOCNIVEAU  \DOCMODULE  \DOCLECON (numéro)  \DOCTITRE
\documentclass[11pt]{article}

\usepackage{fontspec}
\usepackage[ngerman,french]{babel} % français = langue principale ; l'allemand via \all{..} / textebox
\usepackage[a4paper,margin=2cm,top=3.4cm,bottom=2.8cm,headheight=1.4cm,footskip=1.3cm]{geometry}
\usepackage{amsmath,amssymb,amsfonts}
\usepackage{mathtools} % \xmapsto, \coloneqq... souvent utilisés par les modèles
\usepackage{cancel} % simplifications barrées (bilans de forces)
% \abs{z} (module, valeur absolue) et \norm{u} (norme), utilisés par les modèles
\providecommand{\abs}{}\let\abs\relax\DeclarePairedDelimiter{\abs}{\lvert}{\rvert}
\providecommand{\norm}{}\let\norm\relax\DeclarePairedDelimiter{\norm}{\lVert}{\rVert}
\usepackage[table]{xcolor} % [table] : fournit aussi \rowcolors (alternance de lignes)
\usepackage{pagecolor}
\usepackage{tikz}
\usetikzlibrary{arrows.meta,positioning,calc,shapes.geometric,shapes.misc,shapes.arrows,decorations.pathmorphing,decorations.pathreplacing,decorations.markings,patterns,shadows,mindmap,trees,fit,backgrounds,matrix,intersections,angles,quotes}
\usepackage{pgfplots}
\usepackage[version=4]{mhchem} % équations nucléaires \ce{^{235}_{92}U}
\usepackage[european,siunitx]{circuitikz} % schémas électriques
\pgfplotsset{compat=1.18}
\usepgfplotslibrary{fillbetween,groupplots} % aires entre courbes (calcul intégral)
\usepackage{enumitem}
\usepackage{fancyhdr}
% [most] charge toutes les bibliothèques tcolorbox (listings, minted, raster,
% documentation...) et gonfle inutilement la mémoire TeX sur un document long
% (~300 boîtes) : on ne charge que ce qui est réellement utilisé.
\usepackage[skins,breakable]{tcolorbox}
\usepackage{graphicx}
\usepackage{colortbl}
\usepackage{booktabs}
\usepackage{tabularx}
\usepackage{diagbox} % case d'en-tête diagonale des tableaux à double entrée
% Colonnes extensibles centrée / gauche / droite (C, L, R), souvent utilisées par les modèles
\newcolumntype{C}{>{\centering\arraybackslash}X}
\newcolumntype{L}{>{\raggedright\arraybackslash}X}
\newcolumntype{R}{>{\raggedleft\arraybackslash}X}
\usepackage{array}
\usepackage{multicol}
\usepackage{multirow}
\usepackage{siunitx}
% parse-numbers=false : accepte aussi \SI{2,5 \times 10^{-3}}{...}
\sisetup{locale=FR,output-decimal-marker={,},group-minimum-digits=4,parse-numbers=false}
\DeclareSIUnit\ohm{\ensuremath{\symup{\Omega}}} % U+2126 absent de la police
\DeclareSIUnit\division{div} % réglages d'oscilloscope (ms/div, V/div)
\DeclareSIUnit\tr{tr} % tours (vitesse de rotation, ex. \SI{1500}{\tr\per\minute})
\DeclareSIUnit\molar{M} % concentration molaire (ex. \SI{3}{\molar}, \SI{1}{\micro\molar})
\DeclareSIUnit\an{an} % année (le modèle confond parfois avec \year, un compteur TeX) — ex. \SI{5}{\kilo\meter\per\an}
\DeclareSIUnit\FCFA{FCFA} % franc CFA (ex. \SI{500}{\FCFA\per\kilo\gram})
\DeclareSIUnit\degreeBrix{\textdegree Brix} % teneur en sucre d'un sirop/jus (réfractométrie)
\DeclareSIUnit\month{mois} % le modèle utilise parfois \month, un compteur TeX, comme unité de durée
\usepackage{qrcode}
% \degree : commande fréquemment utilisée par les modèles pour un angle ou une
% température en dehors de \SI{}{} (non fournie par les paquets chargés).
\providecommand{\degree}{^{\circ}}
% \qed (fin de démonstration), utilisé par les modèles sans amsthm
\providecommand{\qed}{\hfill$\square$}
% \barre : double barre des valeurs interdites dans un tableau de variations (tkz-tab)
\providecommand{\barre}{\ensuremath{\|}}
% environnement proof (amsthm non chargé) utilisé par les modèles
\ifdefined\proof\else\newenvironment{proof}[1][Démonstration]{\par\noindent\textit{#1.}\ }{\qed\par}\fi
\usepackage{lastpage}
\usepackage{hyperref}
\hypersetup{hidelinks}

% ============================================================
% Palette « Pop » OnBuch+ — mêmes hex que la classe P (plus_ui.dart)
% ============================================================
\definecolor{poporange}{HTML}{F59321}
\definecolor{popdark}{HTML}{E07A0C}
\definecolor{popcream}{HTML}{FFF6E8}
\definecolor{popink}{HTML}{1C1714}
\definecolor{poppurple}{HTML}{7A5AE0}
\definecolor{popgreen}{HTML}{2E9E6A}
\definecolor{poppink}{HTML}{E0457B}
\definecolor{popblue}{HTML}{2F7DE1}
\definecolor{popgold}{HTML}{C98A12}
\definecolor{popmuted}{HTML}{8A7F76}
\definecolor{popline}{HTML}{EADFD2}
\definecolor{popgray}{HTML}{8A7F76}
\colorlet{popgrey}{popgray}
% Alias tolérants (le modèle nomme parfois les couleurs différemment)
\colorlet{popwhite}{white}
\colorlet{popblack}{popink}
\colorlet{popred}{poppink}
\colorlet{popyellow}{popgold}
% Teintes claires (fonds de tableaux/encadrés) — le modèle nomme parfois
% les couleurs avec un suffixe L (ex. popblueL) sans les avoir définies.
\colorlet{poporangeL}{poporange!20}
\colorlet{popdarkL}{popdark!20}
\colorlet{poppurpleL}{poppurple!20}
\colorlet{popgreenL}{popgreen!20}
\colorlet{poppinkL}{poppink!20}
\colorlet{popblueL}{popblue!20}
\colorlet{popgoldL}{popgold!20}
\colorlet{popredL}{popred!20}
\colorlet{popgrayL}{popgray!20}
% Teintes claires pour fonds de tableaux / zones
\colorlet{popblueL}{popblue!10!white}
\colorlet{poporangeL}{poporange!14!white}
\colorlet{popgreenL}{popgreen!12!white}
\colorlet{poppurpleL}{poppurple!12!white}
\colorlet{poppinkL}{poppink!12!white}
\colorlet{popgoldL}{popgold!14!white}

\pagecolor{popcream}

% ============================================================
% Typographie
% ============================================================
\defaultfontfeatures{Ligatures=TeX}
\setmainfont{PlusJakartaSans-Medium.ttf}[Path=../../../fonts/,BoldFont=PlusJakartaSans-Bold.ttf]
% Maths en sans-serif (Fira Math) avec chiffres et lettres droites de Plus
% Jakarta, pour que les formules se fondent dans le texte.
\usepackage[math-style=ISO,bold-style=ISO]{unicode-math}
\setmathfont{FiraMath-Regular.otf}[Path=../../../fonts/]
\setmathfont{PlusJakartaSans-Medium.ttf}[Path=../../../fonts/,range={up/num,up/{Latin,latin},\mathup}]
\usepackage{icomma} % virgule décimale sans espace en mode math
% Caractères mathématiques Unicode absents des polices de texte (Plus Jakarta,
% Archivo Black) : rendus par leur équivalent mathématique, en texte comme en
% formule — sinon ils s'affichent en carré vide (ex. « Arithmétique dans ℤ »).
\usepackage{newunicodechar}
\newunicodechar{ℝ}{\ensuremath{\mathbb{R}}}
\newunicodechar{ℤ}{\ensuremath{\mathbb{Z}}}
\newunicodechar{ℂ}{\ensuremath{\mathbb{C}}}
\newunicodechar{ℕ}{\ensuremath{\mathbb{N}}}
\newunicodechar{ℚ}{\ensuremath{\mathbb{Q}}}
\newunicodechar{𝔻}{\ensuremath{\mathbb{D}}}
\newunicodechar{α}{\ensuremath{\alpha}}
\newunicodechar{β}{\ensuremath{\beta}}
\newunicodechar{γ}{\ensuremath{\gamma}}
\newunicodechar{δ}{\ensuremath{\delta}}
\newunicodechar{ε}{\ensuremath{\varepsilon}}
\newunicodechar{θ}{\ensuremath{\theta}}
\newunicodechar{λ}{\ensuremath{\lambda}}
\newunicodechar{μ}{\ensuremath{\mu}}
\newunicodechar{σ}{\ensuremath{\sigma}}
\newunicodechar{τ}{\ensuremath{\tau}}
\newunicodechar{φ}{\ensuremath{\varphi}}
\newunicodechar{ψ}{\ensuremath{\psi}}
\newunicodechar{ω}{\ensuremath{\omega}}
\newunicodechar{Σ}{\ensuremath{\Sigma}}
% majuscules produites par \MakeUppercase dans les titres (α-aminés -> Α)
\newunicodechar{Α}{\ensuremath{\alpha}}
\newunicodechar{Β}{\ensuremath{\beta}}
\newunicodechar{Γ}{\ensuremath{\Gamma}}
\newunicodechar{Δ}{\ensuremath{\Delta}}
\newunicodechar{Ε}{\ensuremath{\varepsilon}}
\newunicodechar{Θ}{\ensuremath{\Theta}}
\newunicodechar{Λ}{\ensuremath{\Lambda}}
\newunicodechar{Μ}{\ensuremath{\mu}}
\newunicodechar{Τ}{\ensuremath{\tau}}
\newunicodechar{Φ}{\ensuremath{\Phi}}
\newunicodechar{Ψ}{\ensuremath{\Psi}}
\newunicodechar{Ω}{\ensuremath{\Omega}}
\newunicodechar{↦}{\ensuremath{\mapsto}}
\newunicodechar{∈}{\ensuremath{\in}}
\newunicodechar{∉}{\ensuremath{\notin}}
\newunicodechar{∧}{\ensuremath{\wedge}}
\newunicodechar{⇔}{\ensuremath{\Leftrightarrow}}
\newunicodechar{⟺}{\ensuremath{\Longleftrightarrow}}
\newunicodechar{⇒}{\ensuremath{\Rightarrow}}
\newunicodechar{⟹}{\ensuremath{\Longrightarrow}}
\newunicodechar{∘}{\ensuremath{\circ}}
\newunicodechar{∪}{\ensuremath{\cup}}
\newunicodechar{∩}{\ensuremath{\cap}}
\newunicodechar{≡}{\ensuremath{\equiv}}
\newunicodechar{⊂}{\ensuremath{\subset}}
\newunicodechar{⊥}{\ensuremath{\perp}}
\newunicodechar{∃}{\ensuremath{\exists}}
\newunicodechar{∀}{\ensuremath{\forall}}
\newunicodechar{∛}{\ensuremath{\sqrt[3]{\,}}}
\newunicodechar{∜}{\ensuremath{\sqrt[4]{\,}}}
\newunicodechar{⃗}{\ensuremath{{}^{\to}}}
\newunicodechar{ⁿ}{\ifmmode^{n}\else\textsuperscript{\ensuremath{n}}\fi}
\newunicodechar{ˣ}{\ifmmode^{x}\else\textsuperscript{\ensuremath{x}}\fi}
\newunicodechar{ᵇ}{\ifmmode^{b}\else\textsuperscript{\ensuremath{b}}\fi}
\newunicodechar{ʳ}{\ifmmode^{r}\else\textsuperscript{\ensuremath{r}}\fi}
\newunicodechar{ᵘ}{\ifmmode^{u}\else\textsuperscript{\ensuremath{u}}\fi}
\newunicodechar{ʸ}{\ifmmode^{y}\else\textsuperscript{\ensuremath{y}}\fi}
\newunicodechar{ᵃ}{\ifmmode^{a}\else\textsuperscript{\ensuremath{a}}\fi}
\newunicodechar{ᵉ}{\ifmmode^{e}\else\textsuperscript{\ensuremath{e}}\fi}
\newunicodechar{ᵏ}{\ifmmode^{k}\else\textsuperscript{\ensuremath{k}}\fi}
\newunicodechar{ᵖ}{\ifmmode^{p}\else\textsuperscript{\ensuremath{p}}\fi}
\newunicodechar{ᴺ}{\ifmmode^{N}\else\textsuperscript{\ensuremath{N}}\fi}
\newunicodechar{ᶜ}{\ifmmode^{c}\else\textsuperscript{\ensuremath{c}}\fi}
\newunicodechar{ʰ}{\ifmmode^{h}\else\textsuperscript{\ensuremath{h}}\fi}
\newunicodechar{ᵠ}{\ifmmode^{q}\else\textsuperscript{\ensuremath{q}}\fi}
\newunicodechar{ₙ}{\ifmmode_{n}\else\textsubscript{\ensuremath{n}}\fi}
\newunicodechar{ₐ}{\ifmmode_{a}\else\textsubscript{\ensuremath{a}}\fi}
\newunicodechar{ᵢ}{\ifmmode_{i}\else\textsubscript{\ensuremath{i}}\fi}
\newunicodechar{ᵣ}{\ifmmode_{r}\else\textsubscript{\ensuremath{r}}\fi}
\newunicodechar{ₖ}{\ifmmode_{k}\else\textsubscript{\ensuremath{k}}\fi}
\newunicodechar{ᵦ}{\ifmmode_{\beta}\else\textsubscript{\ensuremath{\beta}}\fi}
\newunicodechar{ₓ}{\ifmmode_{x}\else\textsubscript{\ensuremath{x}}\fi}
\newfontfamily\popdisplay{ArchivoBlack-Regular.ttf}[Path=../../../fonts/]
\newfontfamily\poplogo{SpaceGrotesk-Bold.ttf}[Path=../../../fonts/]
\newfontfamily\popsemi{PlusJakartaSans-SemiBold.ttf}[Path=../../../fonts/]
\newfontfamily\popxbold{PlusJakartaSans-ExtraBold.ttf}[Path=../../../fonts/]
\newfontfamily\popxboldls{PlusJakartaSans-ExtraBold.ttf}[Path=../../../fonts/,LetterSpace=3.0]

\setlist[enumerate]{leftmargin=1.7em,itemsep=5pt,topsep=4pt}
\setlist[itemize]{leftmargin=1.7em,itemsep=4pt,topsep=4pt,label={\color{poporange}\textbullet}}
\setlength{\parindent}{0pt}
\setlength{\parskip}{5pt}
\renewcommand{\arraystretch}{1.25}

% pgfplots : style Pop par défaut
\pgfplotsset{
  every axis/.append style={axis line style={popink,line width=1pt},tick style={popink},
    grid style={popline,line width=0.5pt},label style={font=\small},tick label style={font=\footnotesize}},
  popcurve/.style={poporange,line width=1.8pt,no markers,smooth},
}
% Même style disponible dans un \draw TikZ simple (hors pgfplots)
\tikzset{popcurve/.style={poporange,line width=1.8pt}}
\tikzset{popgrid/.style={popline,very thin}}
\colorlet{popcurve}{poporange} % le nom du style est parfois utilisé comme couleur

% ============================================================
% Pill / squiggle
% ============================================================
\newcommand{\poppill}[3]{%
  \tikz[baseline=-0.55ex]\node[draw=popink,line width=1.2pt,rounded corners=2.6pt,%
    fill=#2,inner xsep=6.5pt,inner ysep=3pt]{%
    \popxboldls\fontsize{7}{7}\selectfont\color{#3}\MakeUppercase{#1}};%
}
\newcommand{\popsquiggle}[1]{%
  \tikz[baseline=-0.4ex]{
    \draw[#1,line width=2.6pt,line cap=round]
      plot [smooth] coordinates {(0,0.09) (0.34,0.01) (0.68,0.21) (1.02,0.01) (1.36,0.10)};
  }%
}
% Mot-clé surligné (marqueur orange sous le texte)
\newcommand{\cle}[1]{{\popxbold\color{popdark}#1}}

% ============================================================
% En-tête / pied
% ============================================================
\pagestyle{fancy}
\fancyhf{}
\renewcommand{\headrulewidth}{0pt}
\renewcommand{\footrulewidth}{0pt}
\lhead{\raisebox{-0.15ex}{\poplogo\fontsize{13}{13}\selectfont\color{popink}OnBuch\color{poporange}+}}
\rhead{\poppill{\DOCMATIERE\ \textbullet\ \DOCNIVEAU\ \textbullet\ Leçon \DOCLECON}{white}{popink}}
\lfoot{\popxbold\fontsize{7.6}{7.6}\selectfont\color{popmuted}\MakeUppercase{Cours signé OnBuch+}\ \textbullet\ {\popsemi\DOCTITRE}}
\rfoot{\tikz[baseline=-0.6ex]\node[rounded corners=8.5pt,fill=popink,minimum height=17pt,minimum width=17pt,inner xsep=5pt,inner sep=0]{\popxbold\fontsize{8}{8}\selectfont\color{popcream}\thepage\,/\,\pageref*{LastPage}};}

% ============================================================
% Titres
% ============================================================
\newcommand{\chaptitre}[2]{%
  \begin{tcolorbox}[enhanced,colback=poporange,colframe=popink,boxrule=2.6pt,arc=11pt,
      drop shadow=popink,left=15pt,right=15pt,top=12pt,bottom=11pt,before skip=3pt,after skip=18pt]
    \poppill{#2}{popink}{popcream}\par\vspace{9pt}
    {\raggedright\hyphenpenalty=10000\exhyphenpenalty=10000\popdisplay\fontsize{22}{24}\selectfont\color{popcream}\MakeUppercase{#1}\par}
  \end{tcolorbox}
}
% Section numérotée : pastille orange avec le numéro + titre Archivo
\newcounter{popsec}
\newcommand{\coursec}[1]{%
  \stepcounter{popsec}\setcounter{popsub}{0}%
  \par\vspace{16pt}\noindent
  \begin{minipage}{\linewidth}
  \tikz[baseline=-0.7ex]\node[draw=popink,line width=1.6pt,fill=poporange,rounded corners=4pt,minimum size=22pt,inner sep=0]{\popdisplay\fontsize{12}{12}\selectfont\color{popcream}\thepopsec};%
  \hspace{8pt}{\popdisplay\fontsize{15}{17}\selectfont\color{popink}\MakeUppercase{#1}}\par
  \vspace{4pt}\noindent{\color{poporange}\rule{36pt}{3.6pt}}
  \end{minipage}\par\nopagebreak\vspace{8pt}
}
\newcounter{popsub}
\newcommand{\courssub}[1]{%
  \stepcounter{popsub}%
  \par\vspace{9pt}\noindent{\popxbold\fontsize{11.5}{13}\selectfont\color{popink}{\color{poporange}\thepopsec.\thepopsub}\hspace{6pt}#1}\par\nopagebreak\vspace{4pt}
}

% ============================================================
% Boîtes pédagogiques — même recette : fond blanc, bordure épaisse colorée,
% ombre dure encre, badge plein. Titre optionnel [#1].
% ============================================================
\tcbset{popbase/.style={breakable,enhanced,colback=white,boxrule=2.3pt,arc=9pt,
  shadow={3pt}{-3pt}{0pt}{popink},left=14pt,right=14pt,top=11pt,bottom=11pt,before skip=11pt,after skip=15pt,
  fontupper=\popsemi\fontsize{10.6}{14.5}\selectfont\color{popink},
  fontlower=\popsemi\fontsize{10.6}{14.5}\selectfont\color{popink}}}
% Coche dessinée (le glyphe ✓ n'existe pas dans les polices du document)
\newcommand{\popcheck}[1]{\tikz[baseline=-0.5ex]\draw[#1,line width=1.6pt,line cap=round,line join=round](-0.13,0.0)--(-0.04,-0.09)--(0.13,0.1);}
\providecommand{\coche}{\popcheck{popgreen}}
\providecommand{\resultat}[1]{\par\smallskip\noindent\fcolorbox{popgreen}{popgreenL}{\parbox{\dimexpr\linewidth-2\fboxsep-2\fboxrule\relax}{\textbf{Résultat.} #1}}\par\smallskip}
\newcommand{\popboxhead}[3]{\poppill{#1}{#2}{white}\ifx\relax#3\relax\else\hspace{6pt}{\popxbold\fontsize{10.6}{12}\selectfont\color{popink}#3}\fi\par\vspace{7pt}}

\newtcolorbox{aretenir}[1][]{popbase,colframe=popgreen,before upper={\popboxhead{À retenir}{popgreen}{#1}}}
% Texte en allemand (document de lecture, dialogue, extrait) : cadre
% d'étude. Un titre optionnel (source, auteur) s'affiche dans l'en-tête.
\newtcolorbox{textebox}[1][]{popbase,colframe=popblue,before upper={\popboxhead{Text}{popblue}{#1}\begin{otherlanguage}{ngerman}},after upper={\end{otherlanguage}}}
% Marques ✓ / ✗ (forme correcte / fautive) souvent utilisées par les modèles
\providecommand{\cross}{\textcolor{poppink}{\ensuremath{\times}}}
\providecommand{\xmark}{\cross}\providecommand{\crossmark}{\cross}\providecommand{\cmark}{\ensuremath{\checkmark}}
% Ligne de réponse / justification à compléter (inventée par les modèles)
\providecommand{\justification}[1]{\par\noindent\textit{Justification~:} \dotfill\par}
% \textipa (package tipa non chargé) : la police ne couvre pas l'API -> simple passage
\providecommand{\textipa}[1]{#1}
% Soulignement (ulem non chargé) : \uline, \ul
\providecommand{\uline}[1]{\underline{#1}}\providecommand{\ul}[1]{\underline{#1}}
% variantes ASCII d'ordinaux écrites en macro
\providecommand{\iere}{\textsuperscript{re}}\providecommand{\ieme}{\textsuperscript{e}}\providecommand{\ere}{\textsuperscript{re}}\providecommand{\eme}{\textsuperscript{e}}\providecommand{\er}{\textsuperscript{er}}\providecommand{\ie}{\textsuperscript{e}}
% Ordinaux français écrits en macro par les modèles : 1\ière, 3\ième
\newcommand{\ière}{\textsuperscript{re}}\newcommand{\ième}{\textsuperscript{e}}
% \exemple (inventée par les modèles) : étiquette « Exemple : »
\providecommand{\exemple}{\textbf{Exemple~:}\ }
% \square utilisable aussi hors mode math (cases à cocher)
\AtBeginDocument{\let\squareMath\square\renewcommand{\square}{\ensuremath{\squareMath}}}
% Mot ou phrase en allemand à l'intérieur d'un paragraphe français
\newcommand{\all}[1]{\ifmmode\text{\textit{#1}}\else\begingroup\selectlanguage{ngerman}\itshape #1\endgroup\fi}
\let\esp\all % alias toléré
\newtcolorbox{exemplebox}[1][]{popbase,colframe=poporange,before upper={\popboxhead{Exemple}{poporange}{#1}}}
\newtcolorbox{definition}[1][]{popbase,colframe=popblue,before upper={\popboxhead{Définition}{popblue}{#1}}}
\newtcolorbox{propriete}[1][]{popbase,colframe=poppurple,before upper={\popboxhead{Propriété}{poppurple}{#1}}}
\newtcolorbox{methode}[1][]{popbase,colframe=popgold,before upper={\popboxhead{Méthode}{popgold}{#1}}}
\newtcolorbox{attention}[1][]{popbase,colframe=poppink,before upper={\popboxhead{Attention}{poppink}{#1}}}
\newtcolorbox{experience}[1][]{popbase,colframe=popblue,colback=popblueL,before upper={\popboxhead{Expérience}{popblue}{#1}}}
% « Le savais-tu ? » : carte violette pleine (contexte, culture, Cameroun)
\newtcolorbox{savaistu}[1][]{popbase,colback=poppurple,colframe=popink,
  fontupper=\popsemi\fontsize{10.4}{14.2}\selectfont\color{white},
  before upper={\poppill{Le savais-tu\,?}{white}{poppurple}\ifx\relax#1\relax\else\hspace{6pt}{\popxbold\color{white}#1}\fi\par\vspace{7pt}}}
% Exercice résolu : énoncé en haut, \tcblower puis la solution sur fond crème
\newcounter{popexo}\newcounter{popexr}
\newtcolorbox{exoresolu}[1][]{popbase,colframe=poporange,colbacklower=popcream,
  segmentation style={popink,line width=1.2pt,dashed},
  before upper={\stepcounter{popexr}\popboxhead{Exercice résolu \thepopexr}{poporange}{#1}},
  before lower={\poppill{Solution}{popink}{popcream}\par\vspace{6pt}}}
% Exercice (non résolu en place) — la correction est dans la partie Corrigés
\newtcolorbox{exercice}[1][]{popbase,colframe=popink,
  before upper={\stepcounter{popexo}\popboxhead{Exercice \thepopexo}{popink}{#1}}}
% Corrigé
\newtcolorbox{corrige}[1][]{popbase,colframe=popgreen,colback=popgreenL,
  before upper={\popboxhead{Corrigé}{popgreen}{#1}}}
% Objectifs / prérequis / bilan
\newtcolorbox{objectifs}{popbase,colframe=popink,colback=poporangeL,
  before upper={\popboxhead{Objectifs de la leçon}{popink}{}}}
\newtcolorbox{prerequis}{popbase,colframe=popink,colback=popblueL,
  before upper={\popboxhead{Prérequis}{popblue}{}}}
\newtcolorbox{bilan}{popbase,colframe=popink,colback=white,boxrule=2.8pt,
  before upper={\popboxhead{Fiche bilan}{poporange}{L'essentiel à connaître pour l'examen}}}
% Figure encadrée avec légende
\newtcolorbox{popfigure}[1][]{enhanced,breakable=false,colback=white,colframe=popline,boxrule=1.4pt,arc=8pt,
  left=10pt,right=10pt,top=10pt,bottom=8pt,before skip=10pt,after skip=12pt,halign=center,
  lower separated=false,halign lower=center,
  fontlower=\popsemi\fontsize{9}{11.5}\selectfont\color{popmuted},
  #1}
\newcommand{\legende}[1]{\par\vspace{6pt}{\popsemi\fontsize{9}{11.5}\selectfont\color{popmuted}#1}}

% ============================================================
% Signature OnBuch+
% ============================================================
\newcommand{\poplogoinline}[1]{{\poplogo\fontsize{#1}{#1}\selectfont\color{popink}OnBuch\color{poporange}+}}
\newcommand{\signatureonbuch}{%
  \begin{tcolorbox}[enhanced,colback=popink,colframe=popink,boxrule=2.2pt,arc=10pt,
      drop shadow=poporange,left=14pt,right=14pt,top=10pt,bottom=10pt,before skip=0pt,after skip=0pt]
    \tikz[baseline=-0.7ex]\node[circle,fill=popgreen,draw=popcream,line width=1.3pt,minimum size=22pt,inner sep=0]{\popcheck{white}};%
    \hspace{9pt}\begin{minipage}[c]{0.62\linewidth}
      {\popxbold\fontsize{12}{13}\selectfont\color{popcream}Signé\ \textbullet\ {\poplogo OnBuch\color{poporange}+}}\par\vspace{2pt}
      {\raggedright\popsemi\fontsize{8}{10.5}\selectfont\color{popline}\MakeUppercase{Équipe pédagogique OnBuch+}\\\MakeUppercase{Conforme au programme officiel MINESEC}\par}
    \end{minipage}\hfill
    {\poplogo\fontsize{22}{22}\selectfont\color{popcream}OnBuch\color{poporange}+}
  \end{tcolorbox}%
}

% ============================================================
% Page de garde
% #1 titre · #2 matière · #3 niveau · #4 module · #5 n° leçon · #6 durée ·
% #7 description / objectifs (texte ou itemize)
% ============================================================
\newcommand{\pagegarde}[7]{%
  \thispagestyle{empty}
  \newgeometry{a4paper,margin=1.7cm,top=1.6cm,bottom=1.4cm}%
  \begin{tikzpicture}[remember picture,overlay]
    \fill[poporange!18] ([xshift=-3.2cm,yshift=-3.6cm]current page.north east) circle (4.6cm);
    \fill[poppurple!14] ([xshift=2.4cm,yshift=7.6cm]current page.south west) circle (3.8cm);
  \end{tikzpicture}%
  {\poplogo\fontsize{30}{30}\selectfont\color{popink}OnBuch\color{poporange}+}\hfill
  \raisebox{0.6ex}{\poppill{Cours signé \textbullet\ Premium}{popink}{popcream}}\par
  \vspace{1pt}\popsquiggle{poppurple}\par
  \vspace{8pt}
  \begin{tcolorbox}[enhanced,colback=poporange,colframe=popink,boxrule=2.8pt,arc=13pt,
      drop shadow=popink,left=18pt,right=18pt,top=12pt,bottom=13pt,after skip=10pt]
    \poppill{#2 \textbullet\ #3}{popink}{popcream}\hspace{5pt}\poppill{#4}{white}{popink}\par\vspace{8pt}
    {\popdisplay\fontsize{14}{14}\selectfont\color{popink}LEÇON #5}\par\vspace{4pt}
    {\raggedright\hyphenpenalty=10000\exhyphenpenalty=10000\popdisplay\fontsize{25}{27}\selectfont\color{popcream}\MakeUppercase{#1}\par}
    \vspace{8pt}
    {\popxbold\fontsize{9.5}{9.5}\selectfont\color{popink}\MakeUppercase{Durée conseillée : #6}}
  \end{tcolorbox}
  \begin{tcolorbox}[enhanced,colback=white,colframe=popink,boxrule=2.2pt,arc=10pt,
      drop shadow=popink,left=16pt,right=16pt,top=10pt,bottom=10pt,after skip=8pt]
    \poppill{Dans cette leçon}{poppurple}{white}\par\vspace{5pt}
    \popsemi\fontsize{9.3}{12.6}\selectfont\color{popink}#7
  \end{tcolorbox}
  \noindent\begin{minipage}[t]{0.487\linewidth}
    \begin{tcolorbox}[enhanced,colback=popgreen,colframe=popink,boxrule=2.2pt,arc=10pt,
        drop shadow=popink,left=11pt,right=11pt,top=10pt,bottom=10pt,height=3.1cm,valign=center]
      \tcbox[colback=white,colframe=popink,boxrule=1.3pt,arc=5pt,boxsep=0pt,nobeforeafter,
        left=3pt,right=3pt,top=3pt,bottom=3pt,valign=center]{\qrcode[height=1.9cm]{https://play.google.com/store/apps/details?id=com.luvvix.onbuch&hl=fr}}%
      \hspace{8pt}\begin{minipage}[c]{0.5\linewidth}
        {\popdisplay\fontsize{11}{12.5}\selectfont\color{white}\MakeUppercase{Télécharge OnBuch}}\par\vspace{4pt}
        {\raggedright\popsemi\fontsize{8.4}{11}\selectfont\color{white}Gratuit \textbullet\ Google Play\\Tuteur IA, annales, concours, résultats\par}
      \end{minipage}
    \end{tcolorbox}
  \end{minipage}\hfill
  \begin{minipage}[t]{0.487\linewidth}
    \begin{tcolorbox}[enhanced,colback=poppurple,colframe=popink,boxrule=2.2pt,arc=10pt,
        drop shadow=popink,left=11pt,right=11pt,top=10pt,bottom=10pt,height=3.1cm,valign=center]
      \tikz[baseline=-0.7ex]\node[circle,fill=white,draw=popink,line width=1.3pt,minimum size=1.6cm,inner sep=0]{\popdisplay\fontsize{24}{24}\selectfont\color{poppurple}+};%
      \hspace{8pt}\begin{minipage}[c]{0.6\linewidth}
        {\popdisplay\fontsize{11}{12.5}\selectfont\color{white}\MakeUppercase{Passe à OnBuch+}}\par\vspace{4pt}
        {\raggedright\popsemi\fontsize{8.4}{11}\selectfont\color{white}Tous les cours signés, Tuteur IA illimité, fiches PDF — onglet OnBuch+ de l'app\par}
      \end{minipage}
    \end{tcolorbox}
  \end{minipage}
  \vspace{8pt}
  \popcontenu
  \vfill
  \signatureonbuch
  \restoregeometry
  \newpage
}

% Rangée « Ce cours contient » (page de garde)
\newcommand{\poptile}[3]{% #1 couleur #2 gros texte #3 légende
  \begin{tcolorbox}[enhanced,colback=#1,colframe=popink,boxrule=2pt,arc=9pt,shadow={2.5pt}{-2.5pt}{0pt}{popink},
      left=6pt,right=6pt,top=9pt,bottom=9pt,halign=center,valign=center,height=1.75cm,width=\linewidth,nobeforeafter]
    {\popdisplay\fontsize{12.5}{13}\selectfont\color{white}\MakeUppercase{#2}}\par\vspace{3pt}
    {\centering\popsemi\fontsize{7.8}{9.5}\selectfont\color{white}#3\par}
  \end{tcolorbox}}
\newcommand{\popcontenu}{%
  \par\noindent{\popxboldls\fontsize{7.5}{7.5}\selectfont\color{popmuted}CE COURS SIGNÉ CONTIENT}\par\vspace{7pt}
  \noindent\begin{minipage}[t]{0.235\linewidth}\poptile{popblue}{Cours}{complet, illustré, pas à pas}\end{minipage}\hfill
  \begin{minipage}[t]{0.235\linewidth}\poptile{poporange}{Méthodes}{et exercices résolus}\end{minipage}\hfill
  \begin{minipage}[t]{0.235\linewidth}\poptile{poppink}{Exercices}{type examen + corrigés}\end{minipage}\hfill
  \begin{minipage}[t]{0.235\linewidth}\poptile{popgreen}{Fiche bilan}{l'essentiel à retenir}\end{minipage}\par}

% Sommaire
\newtcolorbox{sommaire}{enhanced,colback=white,colframe=popink,boxrule=2.2pt,arc=10pt,
  drop shadow=popink,left=18pt,right=18pt,top=13pt,bottom=13pt,before skip=6pt,after skip=20pt,
  before upper={\poppill{Sommaire}{poppurple}{white}\par\vspace{9pt}\popsemi\fontsize{10.6}{14.5}\selectfont\color{popink}}}

% Signature de fin de cours — poussée tout en bas de la dernière page
\newcommand{\findecours}{%
  \par\vfill\nopagebreak
  \begin{center}\popsquiggle{poporange}\end{center}\vspace{4pt}
  \signatureonbuch
}
\AtBeginDocument{\def\llbracket{\mathopen{[\mkern-3.5mu[}}\def\rrbracket{\mathclose{]\mkern-3.5mu]}}} % intervalles d'entiers (glyphe ⟦ absent de la police)
% Glyphes absents de Fira Math : redéfinis à partir de glyphes disponibles
\AtBeginDocument{%
  \def\setminus{\mathbin{\backslash}}\let\smallsetminus\setminus
  \def\vdots{\mathord{\vbox{\baselineskip3.2pt\lineskiplimit0pt\kern4pt\hbox{.}\hbox{.}\hbox{.}}}}%
  \def\ddots{\mathinner{\mkern1mu\raise6.5pt\hbox{.}\mkern2mu\raise3.6pt\hbox{.}\mkern2mu\raise0.7pt\hbox{.}\mkern1mu}}%
  \def\triangle{\mathord{\scalebox{0.9}{$\symup{\Delta}$}}}%
  \def\nabla{\mathord{\raisebox{\depth}{\scalebox{1}[-1]{$\symup{\Delta}$}}}}%
  \def\checkmark{\popcheck{popdark}}%
  \def\star{\mathbin{\ast}}\def\bigstar{\mathord{\ast}}%
  \def\diamond{\mathbin{\scalebox{0.6}{$\Diamond$}}}%
  \def\bigcup{\mathop{\mathchoice{\raisebox{-0.3em}{\scalebox{1.7}{$\cup$}}}{\scalebox{1.25}{$\cup$}}{\cup}{\cup}}\displaylimits}%
  \def\bigcap{\mathop{\mathchoice{\raisebox{-0.3em}{\scalebox{1.7}{$\cap$}}}{\scalebox{1.25}{$\cap$}}{\cap}{\cap}}\displaylimits}%
  \def\lesseqqgtr{\mathrel{\lessgtr}}\def\bigoplus{\mathop{\oplus}\displaylimits}%
}
\providecommand{\ssi}{si et seulement si} % abréviation fréquente
\AtBeginDocument{\providecommand{\wideparen}{\overparen}} % arc de cercle

%% FIGFIX-BEGIN v5 — correctifs globaux des figures (docs/onbuch-generation/tools/figfix)
\usepackage{adjustbox}\usepackage{etoolbox}\tcbuselibrary{raster}
% toute figure TikZ/pgfplots plus large que la ligne est réduite à la largeur disponible
\BeforeBeginEnvironment{tikzpicture}{\begin{adjustbox}{max width=\linewidth}}
\AfterEndEnvironment{tikzpicture}{\end{adjustbox}}
% légendes pgfplots lisibles par-dessus les courbes
\pgfplotsset{every axis/.append style={legend style={fill=white,fill opacity=0.92,text opacity=1,draw=black!15,inner sep=3pt}}}
% étiquettes d'arêtes (syntaxe edge["..."]) avec fond blanc
\tikzset{every edge quotes/.append style={fill=white,inner sep=1.2pt}}
%% FIGFIX-END
\begin{document}

\pagegarde{Les types de mariages}{Allemand}{1ère A}{Chapitre 1 : Vie familiale et sociale}{ALA02}{2 h}{Cette leçon de type « texte » propose la lecture et le commentaire d'un texte original sur les types de mariages en Allemagne et au Cameroun. Elle permet d'acquérir le vocabulaire thématique du mariage et de maîtriser l'accusatif, le datif et les verbes de modalité au présent.
\vspace{4pt}
\begin{multicols}{2}\small
\begin{itemize}
\item À la fin de cette leçon, je sais comprendre globalement puis en détail un texte sur les types de mariages.
\item À la fin de cette leçon, je sais employer correctement le vocabulaire du mariage (die Hochzeit, die Verlobung, die Mitgift, standesamtlich...) avec articles et pluriels.
\item À la fin de cette leçon, je sais distinguer et utiliser l'accusatif et le datif après les verbes et les prépositions courantes.
\item À la fin de cette leçon, je sais conjuguer les verbes de modalité au présent (dürfen, können, müssen, sollen, wollen, mögen/möchten).
\item À la fin de cette leçon, je sais construire une phrase avec un verbe modal et l'infinitif en fin de phrase.
\item À la fin de cette leçon, je sais répondre à des questions de compréhension en phrases complètes.
\end{itemize}\end{multicols}}

\begin{sommaire}
\begin{multicols}{2}
\begin{enumerate}
\item Texte support : Heiraten in Deutschland und in Kamerun
\item Compréhension du texte
\item Lexique thématique : die Hochzeit und die Ehe
\item Grammaire 1 : l'accusatif et le datif
\item Grammaire 2 : les verbes de modalité au présent
\item Méthode du commentaire et production guidée
\item Activité d'intégration
\item Méthodes et exercices résolus
\item Exercices
\item Corrigés des exercices
\item Fiche bilan
\end{enumerate}
\end{multicols}
\end{sommaire}

\begin{objectifs}
\begin{itemize}
\item À la fin de cette leçon, je sais comprendre globalement puis en détail un texte sur les types de mariages.
\item À la fin de cette leçon, je sais employer correctement le vocabulaire du mariage (die Hochzeit, die Verlobung, die Mitgift, standesamtlich...) avec articles et pluriels.
\item À la fin de cette leçon, je sais distinguer et utiliser l'accusatif et le datif après les verbes et les prépositions courantes.
\item À la fin de cette leçon, je sais conjuguer les verbes de modalité au présent (dürfen, können, müssen, sollen, wollen, mögen/möchten).
\item À la fin de cette leçon, je sais construire une phrase avec un verbe modal et l'infinitif en fin de phrase.
\item À la fin de cette leçon, je sais répondre à des questions de compréhension en phrases complètes.
\item À la fin de cette leçon, je sais rédiger un résumé et un court paragraphe de comparaison entre le Cameroun et l'Allemagne.
\item À la fin de cette leçon, je sais commenter un texte selon la méthode : présentation, compréhension, analyse, opinion personnelle.
\end{itemize}
\end{objectifs}

\begin{prerequis}
\begin{itemize}
\item Les articles définis et indéfinis au nominatif (der, die, das / ein, eine) vus en Seconde.
\item La conjugaison des verbes réguliers et irréguliers au présent.
\item La structure de base de la phrase allemande (verbe conjugué en 2e position).
\item Le vocabulaire de la famille (der Vater, die Mutter, die Eltern...) vu en Seconde.
\item Les pronoms personnels (ich, du, er/sie/es, wir, ihr, sie, Sie).
\end{itemize}
\end{prerequis}

\newpage

\coursec{Texte support : Heiraten in Deutschland und in Kamerun}

\courssub{Situation d'accroche et problématique}

À Yaoundé comme à Douala, tout le monde connaît cette scène : le samedi matin, la famille se réunit pour le mariage traditionnel, avec la remise de la dot ; le lundi, le couple passe à la mairie, en costume et en robe blanche ; et parfois, le dimanche suivant, une messe de mariage est célébrée à l'église. Trois cérémonies pour un seul couple ! En Allemagne, la situation est différente : seul le mariage civil, \all{die standesamtliche Trauung}, est juridiquement valable. La cérémonie religieuse est facultative et n'a aucune valeur légale.

Comment parler de ces réalités en allemand ? Comment comparer les traditions camerounaises et celles des pays germanophones ? Cette leçon nous donne le vocabulaire et les outils grammaticaux pour le faire.

\textbf{Problématique :} quels types de mariages existe-t-il en Allemagne et au Cameroun, et comment peut-on les décrire et les comparer en allemand ?

\courssub{Lecture globale du texte}

\begin{methode}[Lire un texte en trois étapes]
Pour bien comprendre un texte allemand, procède toujours en trois étapes :
\begin{enumerate}
\item \textbf{Lecture globale sans dictionnaire :} lis le texte une première fois (ou écoute la lecture de l'enseignant) pour saisir le sens général. Ne t'arrête pas sur les mots inconnus.
\item \textbf{Lecture détaillée avec le glossaire :} relis paragraphe par paragraphe, souligne les mots-clés du champ lexical et vérifie les mots nouveaux dans le glossaire.
\item \textbf{Reformulation :} raconte le contenu du texte avec tes propres mots, d'abord oralement, puis par écrit en quelques phrases simples.
\end{enumerate}
\end{methode}

\textbf{Première étape :} l'enseignant lit le texte à voix haute, lentement. Les élèves écoutent, livre fermé, puis lisent individuellement en silence. Question de compréhension globale : \emph{De quoi parle ce texte ?} (Il compare les mariages en Allemagne et au Cameroun.)

\begin{textebox}[Heiraten in Deutschland und in Kamerun]
In Deutschland heiraten viele Paare zuerst standesamtlich, denn nur die standesamtliche Trauung ist vor dem Gesetz gültig. Das Standesamt befindet sich meistens im Rathaus. Dort unterschreiben der Bräutigam und die Braut die Dokumente, und ein Beamter erklärt das Paar zu Ehemann und Ehefrau. Wer möchte, kann danach auch in der Kirche heiraten. Die kirchliche Trauung ist aber nur eine schöne Zeremonie ohne juristischen Wert. Vor der Hochzeit gibt es oft eine Verlobung: Der Mann schenkt der Frau einen Ring, und die beiden Familien lernen sich kennen.

In Kamerun gibt es meistens drei Hochzeiten: die traditionelle Hochzeit mit der Mitgift für die Familie der Braut, die standesamtliche Trauung im Rathaus und oft auch die kirchliche Trauung in der Kirche. Bei der traditionellen Hochzeit bringt die Familie des Bräutigams Geschenke und die Mitgift zur Familie der Braut. Diese Zeremonie ist sehr wichtig, denn sie verbindet zwei Familien und nicht nur zwei Personen. Viele Gäste essen, trinken, singen und tanzen den ganzen Tag.

In beiden Ländern müssen der Ehemann und die Ehefrau nach der Hochzeit zusammenarbeiten, denn eine Ehe will gelernt sein! Man muss Geduld haben, man darf nicht egoistisch sein, und man soll den Partner respektieren. Die Hochzeit dauert einen Tag, aber die Ehe dauert das ganze Leben.

\emph{Glossaire :} gültig = valable ; das Gesetz, -e = la loi ; das Rathaus, -äuser = la mairie ; der Bräutigam, -e = le fiancé/le marié ; die Braut, -en = la fiancée/la mariée ; unterschreiben = signer ; der Beamte, -n = le fonctionnaire ; der Wert, -e = la valeur ; schenken = offrir ; sich kennenlernen = faire connaissance ; die Mitgift, -en = la dot ; das Geschenk, -e = le cadeau ; verbinden = unir ; die Geduld = la patience ; respektieren = respecter ; dauern = durer.
\end{textebox}

\begin{attention}[Die Hochzeit ou die Ehe ?]
Ne confonds pas \all{die Hochzeit} (la cérémonie, la fête du mariage : un jour précis) et \all{die Ehe} (l'institution, la vie conjugale : toute la vie). De même, \all{die Verlobung} désigne les fiançailles (avant le mariage), tandis que \all{die Trauung} désigne la cérémonie de mariage elle-même. Compare : \all{Die Verlobung war im Januar, die Trauung im Juni.} « Les fiançailles ont eu lieu en janvier, le mariage en juin. »
\end{attention}

\courssub{Identification du thème et du type de texte}

Après la lecture globale, on identifie les caractéristiques du texte :

\begin{center}
\begin{tabularx}{\linewidth}{|X|X|}
\hline
\rowcolor{popblueL} \textbf{Élément} & \textbf{Réponse} \\
\hline
Thème & Les types de mariages (\all{die Hochzeitsformen}) \\
\hline
Type de texte & Texte informatif et descriptif \\
\hline
Structure & 3 paragraphes : Allemagne / Cameroun / conseils pour la vie conjugale \\
\hline
Idée principale & En Allemagne, seul le mariage civil compte juridiquement ; au Cameroun, on célèbre souvent trois mariages. \\
\hline
\end{tabularx}
\end{center}

\begin{definition}[Les idées principales du texte]
\begin{itemize}
\item \textbf{Paragraphe 1 :} en Allemagne, la \all{standesamtliche Trauung} est obligatoire et seule valable devant la loi ; la \all{kirchliche Trauung} est facultative.
\item \textbf{Paragraphe 2 :} au Cameroun, on célèbre souvent trois mariages : traditionnel (avec la \all{Mitgift}), civil et religieux.
\item \textbf{Paragraphe 3 :} dans les deux pays, le couple doit collaborer, car le mariage s'apprend.
\end{itemize}
\end{definition}

\courssub{Lecture détaillée paragraphe par paragraphe}

\textbf{Paragraphe 1 — L'Allemagne.} Souligne les mots-clés du champ lexical du mariage : \all{heiraten}, \all{die standesamtliche Trauung}, \all{das Standesamt}, \all{der Bräutigam}, \all{die Braut}, \all{Ehemann und Ehefrau}, \all{die kirchliche Trauung}, \all{die Verlobung}, \all{der Ring}. Observe la phrase clé : \all{Nur die standesamtliche Trauung ist vor dem Gesetz gültig.} « Seul le mariage civil est valable devant la loi. »

\textbf{Paragraphe 2 — Le Cameroun.} Souligne : \all{die traditionelle Hochzeit}, \all{die Mitgift}, \all{die Familie der Braut}, \all{die Geschenke}, \all{die Gäste}. Observe : \all{Die Familie des Bräutigams bringt Geschenke und die Mitgift zur Familie der Braut.} « La famille du marié apporte des cadeaux et la dot à la famille de la mariée. »

\textbf{Paragraphe 3 — Les conseils.} Ce paragraphe contient plusieurs \cle{verbes de modalité} : \all{müssen} (devoir), \all{dürfen} (avoir le droit), \all{sollen} (devoir, conseil). Observe : \all{Man muss Geduld haben, man darf nicht egoistisch sein, und man soll den Partner respektieren.} « Il faut avoir de la patience, on ne doit pas être égoïste et l'on doit respecter son partenaire. »

\textbf{Repérage grammatical dans le texte :}

\begin{center}
\begin{tabularx}{\linewidth}{|X|X|X|}
\hline
\rowcolor{popblueL} \textbf{Exemple du texte} & \textbf{Cas} & \textbf{Fonction} \\
\hline
\all{Der Mann schenkt der Frau einen Ring.} & einen Ring = accusatif & COD (la chose offerte) \\
\hline
\all{Der Mann schenkt der Frau einen Ring.} & der Frau = datif & COI (la personne) \\
\hline
\all{... mit der Mitgift für die Familie der Braut} & die Familie = accusatif & après la préposition \all{für} \\
\hline
\all{... mit der Mitgift} & der Mitgift = datif & après la préposition \all{mit} \\
\hline
\all{Man soll den Partner respektieren.} & den Partner = accusatif & COD du verbe \all{respektieren} \\
\hline
\end{tabularx}
\end{center}

\begin{aretenir}[Verbes de modalité repérés dans le texte]
\begin{itemize}
\item \all{Wer möchte, kann danach auch in der Kirche heiraten.} — \all{können} (pouvoir) : possibilité.
\item \all{Der Ehemann und die Ehefrau müssen zusammenarbeiten.} — \all{müssen} (devoir) : obligation.
\item \all{Man darf nicht egoistisch sein.} — \all{dürfen} : interdiction avec la négation.
\item \all{Man soll den Partner respektieren.} — \all{sollen} : conseil, recommandation.
\end{itemize}
Règle d'ordre des mots : le modal est conjugué en \cle{deuxième position}, l'infinitif va à la \cle{fin} de la phrase : \all{Man \textbf{muss} Geduld \textbf{haben}.}
\end{aretenir}

\courssub{Carte mentale du champ lexical}

\begin{popfigure}
\begin{center}\tcbox[colback=popink,colframe=popink,coltext=white,arc=9pt,boxsep=4pt,left=6pt,right=6pt]{\bfseries\small die Ehe}\end{center}
\vspace{-2pt}
\begin{tcbraster}[raster columns=2,raster equal height=rows,raster column skip=6pt,raster row skip=6pt,enhanced,blankest]
\begin{tcolorbox}[enhanced,colback=poporange,colframe=poporange,coltext=white,boxrule=0.9pt,arc=3pt,left=4pt,right=4pt,top=3pt,bottom=3pt,boxsep=1pt,halign=left]\bfseries\scriptsize die Verlobung \footnotesize der Ring\end{tcolorbox}
\begin{tcolorbox}[enhanced,colback=popgreen,colframe=popgreen,coltext=white,boxrule=0.9pt,arc=3pt,left=4pt,right=4pt,top=3pt,bottom=3pt,boxsep=1pt,halign=left]\bfseries\scriptsize die standesamtliche Trauung \footnotesize das Rathaus\end{tcolorbox}
\begin{tcolorbox}[enhanced,colback=poppurple,colframe=poppurple,coltext=white,boxrule=0.9pt,arc=3pt,left=4pt,right=4pt,top=3pt,bottom=3pt,boxsep=1pt,halign=left]\bfseries\scriptsize die kirchliche Trauung \footnotesize die Kirche\end{tcolorbox}
\begin{tcolorbox}[enhanced,colback=poppink,colframe=poppink,coltext=white,boxrule=0.9pt,arc=3pt,left=4pt,right=4pt,top=3pt,bottom=3pt,boxsep=1pt,halign=left]\bfseries\scriptsize die traditionelle Hochzeit \footnotesize die Mitgift\end{tcolorbox}
\end{tcbraster}

\legende{Carte mentale : le champ lexical du mariage autour de \all{die Ehe}.}
\end{popfigure}

\courssub{Reformulation orale du texte}

\begin{methode}[Reformuler un texte avec ses propres mots]
\begin{enumerate}
\item Repère les trois idées principales (une par paragraphe).
\item Remplace les structures du texte par des phrases simples : sujet + verbe + compléments.
\item Utilise des connecteurs : \all{zuerst} (d'abord), \all{dann} (ensuite), \all{auch} (aussi), \all{aber} (mais), \all{denn} (car).
\end{enumerate}
\textbf{Modèle de reformulation :} \all{In Deutschland muss jedes Paar standesamtlich heiraten, denn nur diese Trauung ist gültig. Die Kirche ist nicht obligatorisch. In Kamerun feiert man oft drei Hochzeiten: traditionell, standesamtlich und kirchlich. Nach der Hochzeit sollen Ehemann und Ehefrau zusammenarbeiten.} « En Allemagne, chaque couple doit se marier civilement, car seule cette cérémonie est valable. L'église n'est pas obligatoire. Au Cameroun, on fête souvent trois mariages : traditionnel, civil et religieux. Après le mariage, mari et femme doivent collaborer. »
\end{methode}

\begin{exoresolu}[Questions de compréhension]
Réponds en phrases complètes : 1) \all{Wo heiratet man in Deutschland standesamtlich?} 2) \all{Was schenkt der Mann der Frau bei der Verlobung?} 3) \all{Wie viele Hochzeiten gibt es oft in Kamerun?}
\tcblower
1) \all{In Deutschland heiratet man standesamtlich im Rathaus (auf dem Standesamt).} « En Allemagne, on se marie civilement à la mairie. » — 2) \all{Bei der Verlobung schenkt der Mann der Frau einen Ring.} « Aux fiançailles, l'homme offre une bague à la femme. » — 3) \all{In Kamerun gibt es oft drei Hochzeiten: die traditionelle, die standesamtliche und die kirchliche Hochzeit.} « Au Cameroun, il y a souvent trois mariages : traditionnel, civil et religieux. »
\end{exoresolu}

\begin{savaistu}[Le mariage civil obligatoire en Allemagne]
En Allemagne, le mariage religieux seul n'a \textbf{aucune valeur juridique} : le passage au \all{Standesamt} (l'état civil) est obligatoire pour être reconnu comme marié. C'est presque l'inverse de certaines pratiques camerounaises, où le mariage traditionnel avec la dot est souvent considéré comme le « vrai » mariage par les familles, avant même le passage à la mairie. En Autriche et en Suisse, le principe est le même qu'en Allemagne : seul l'acte civil compte devant la loi.
\end{savaistu}

\coursec{Compréhension du texte}

Dans la section précédente, nous avons lu et écouté le texte \all{Heiraten in Deutschland und in Kamerun}. Nous avons découvert les premiers mots du thème et dégagé l'idée générale. Il est temps maintenant de passer à l'étape décisive de toute leçon de texte : la \cle{compréhension}. Comprendre un texte allemand, ce n'est pas traduire mot à mot : c'est savoir répondre à des questions, repérer les informations importantes et justifier ses réponses par le texte. C'est exactement ce qui sera demandé à l'évaluation et au baccalauréat. Nous allons procéder en trois temps : compréhension globale, compréhension détaillée, puis questions de vocabulaire et d'opinion.

Pour travailler, nous nous appuyons sur le texte support, que voici rappelé dans sa version intégrale :

\begin{textebox}[Heiraten in Deutschland und in Kamerun]
In Deutschland und in Kamerun heiraten die Menschen, aber die Traditionen sind sehr verschieden.

In Deutschland gibt es nur eine gültige Hochzeit: die standesamtliche Trauung. Das Paar geht zum Standesamt und unterschreibt dort ein Dokument. Erst dann ist die Ehe offiziell. Viele Paare machen danach auch eine kirchliche Hochzeit in der Kirche, aber das ist nur eine Tradition, kein Gesetz. Vor der Hochzeit gibt es oft eine Verlobung: Der Mann schenkt der Frau einen Ring, und die Familien feiern zusammen.

In Kamerun ist das anders. Hier gibt es drei Arten von Hochzeiten: die traditionelle Hochzeit, die standesamtliche Hochzeit und die kirchliche Hochzeit. Bei der traditionellen Hochzeit besucht die Familie des Mannes die Familie der Frau. Der Mann muss eine Mitgift bezahlen, zum Beispiel Geld, Getränke oder Stoffe. Die beiden Familien essen, trinken und tanzen zusammen. Viele junge Leute in Yaoundé oder Douala machen heute alle drei Hochzeiten: zuerst die traditionelle, dann die standesamtliche und am Ende die kirchliche.

In beiden Ländern ist die Hochzeit ein großes Fest. Die Ehefrau und der Ehemann bekommen viele Geschenke, und die Gäste feiern bis spät in die Nacht. Aber die Wege zur Ehe sind verschieden — und das macht die Kulturen so interessant.
\end{textebox}

\textbf{Glossaire :} \all{heiraten} (se marier) ; \all{die Trauung, die Trauungen} (la cérémonie de mariage) ; \all{das Standesamt, die Standesämter} (l'état civil, la mairie) ; \all{gültig} (valable) ; \all{das Gesetz, die Gesetze} (la loi) ; \all{die Verlobung, die Verlobungen} (les fiançailles) ; \all{schenken} (offrir) ; \all{die Mitgift, die Mitgiften} (la dot) ; \all{der Stoff, die Stoffe} (le tissu, le pagne) ; \all{das Geschenk, die Geschenke} (le cadeau) ; \all{der Gast, die Gäste} (l'invité) ; \all{verschieden} (différent).

\courssub{Compréhension globale : de quoi parle le texte ?}

Avant de chercher des détails, il faut toujours répondre à la question la plus simple : \emph{de quoi parle le texte ?} En allemand, on demande : \all{Worum geht es im Text?} (« De quoi s'agit-il dans le texte ? »).

\begin{exemplebox}[Questions globales et réponses]
\textbf{Frage 1 :} \all{Worum geht es im Text?}\\
\textbf{Antwort :} \all{Im Text geht es um die Hochzeiten in Deutschland und in Kamerun.}\\
« Question : De quoi parle le texte ? — Réponse : Le texte parle des mariages en Allemagne et au Cameroun. »

\textbf{Frage 2 :} \all{Welche Länder werden verglichen?}\\
\textbf{Antwort :} \all{Deutschland und Kamerun werden verglichen.}\\
« Quels pays sont comparés ? — L'Allemagne et le Cameroun sont comparés. »

\textbf{Frage 3 :} \all{Was ist in beiden Ländern gleich?}\\
\textbf{Antwort :} \all{In beiden Ländern ist die Hochzeit ein großes Fest.}\\
« Qu'est-ce qui est identique dans les deux pays ? — Dans les deux pays, le mariage est une grande fête. »
\end{exemplebox}

Remarquez la méthode : la réponse \textbf{reprend les mots de la question} et forme une \textbf{phrase complète}, avec un verbe conjugué. On ne répond jamais par un seul mot comme \all{Hochzeiten} ou \all{Deutschland}.

\courssub{La technique de réponse : de la question à la phrase complète}

\begin{methode}[Répondre à une question de compréhension en quatre étapes]
\begin{enumerate}
\item \textbf{Lire la question} et souligner le mot interrogatif : \all{wer} (qui), \all{was} (quoi), \all{wo} (où), \all{wann} (quand), \all{warum} (pourquoi), \all{welche} (quel).
\item \textbf{Repérer la phrase-clé} dans le texte : chercher les mots de la question ou leurs synonymes dans le texte.
\item \textbf{Construire une phrase complète} qui commence par les mots de la question ou par l'élément demandé — ne jamais recopier tout le paragraphe.
\item \textbf{Justifier} si on le demande : citer la phrase du texte entre guillemets allemands „…“.
\end{enumerate}
\end{methode}

\begin{popfigure}
\begin{tikzpicture}[
  node distance=0.55cm,
  box/.style={draw=popblue, fill=popblueL, rounded corners=3pt, align=center, text width=3.1cm, minimum height=1.1cm, font=\small},
  arr/.style={-{Stealth[length=3mm]}, thick, poporange}
]
\node[box, align=center] (q){\textbf{1. Question}\\ \all{Welche Trauung ist gültig?}};
\node[box, right=of q, align=center] (r){\textbf{2. Repérage}\\ phrase-clé dans le texte};
\node[box, right=of r, align=center] (a){\textbf{3. Réponse}\\ phrase complète};
\node[box, right=of a, align=center] (j){\textbf{4. Justification}\\ citation „…“};
\draw[arr] (q) -- (r);
\draw[arr] (r) -- (a);
\draw[arr] (a) -- (j);
\end{tikzpicture}
\legende{La démarche de compréhension : de la question à la réponse justifiée.}
\end{popfigure}

\courssub{Compréhension détaillée : les informations précises}

Passons maintenant aux questions de détail. Pour chaque question, nous donnons la réponse modèle et la phrase du texte qui la justifie.

\begin{exemplebox}[Questions détaillées et réponses modèles]
\textbf{Frage 1 :} \all{Welche Trauung ist in Deutschland gültig?}\\
\textbf{Antwort :} \all{Die standesamtliche Trauung ist in Deutschland gültig.}\\
Justification : „\all{In Deutschland gibt es nur eine gültige Hochzeit: die standesamtliche Trauung.}“\\
« Quelle cérémonie est valable en Allemagne ? — La cérémonie civile est valable en Allemagne. »

\textbf{Frage 2 :} \all{Was passiert bei der Verlobung?}\\
\textbf{Antwort :} \all{Bei der Verlobung schenkt der Mann der Frau einen Ring, und die Familien feiern zusammen.}\\
« Que se passe-t-il aux fiançailles ? — L'homme offre une bague à la femme et les familles fêtent cela ensemble. »

\textbf{Frage 3 :} \all{Welche drei Arten von Hochzeiten gibt es in Kamerun?}\\
\textbf{Antwort :} \all{In Kamerun gibt es die traditionelle Hochzeit, die standesamtliche Hochzeit und die kirchliche Hochzeit.}\\
« Quels sont les trois types de mariage au Cameroun ? — Le mariage traditionnel, le mariage civil et le mariage religieux. »

\textbf{Frage 4 :} \all{Was muss der Mann bei der traditionellen Hochzeit bezahlen?}\\
\textbf{Antwort :} \all{Bei der traditionellen Hochzeit muss der Mann eine Mitgift bezahlen.}\\
« Que doit payer l'homme au mariage traditionnel ? — Il doit payer une dot. »
\end{exemplebox}

Observez bien la réponse à la question 3 : elle commence par le complément de lieu \all{In Kamerun}, et le verbe \all{gibt} vient immédiatement après, \emph{avant} le sujet \all{es}. C'est la fameuse \cle{inversion} allemande.

\begin{propriete}[L'inversion après un complément en tête de phrase]
En allemand, le \textbf{verbe conjugué occupe toujours la deuxième position} de la phrase. Si la phrase commence par un complément de lieu (\all{In Kamerun}), de temps (\all{Heute}, \all{Vor der Hochzeit}) ou par l'élément demandé dans la question, alors le \textbf{sujet se place après le verbe} : c'est l'inversion.

\all{In Kamerun gibt es drei Hochzeiten.} — « Au Cameroun, il y a trois mariages. »

\all{Bei der Verlobung schenkt der Mann der Frau einen Ring.} — « Aux fiançailles, l'homme offre une bague à la femme. »

\all{Vor der Hochzeit gibt es oft eine Verlobung.} — « Avant le mariage, il y a souvent des fiançailles. »
\end{propriete}

\begin{center}
\begin{tabularx}{\linewidth}{|X|X|X|}
\hline
\rowcolor{popblueL} Français & Deutsch & Remarque \\
\hline
Au Cameroun, il y a trois mariages. & \all{In Kamerun gibt es drei Hochzeiten.} & verbe en 2\textsuperscript{e} position \\
\hline
En Allemagne, seul le mariage civil est valable. & \all{In Deutschland ist nur die standesamtliche Hochzeit gültig.} & sujet après le verbe \\
\hline
L'homme offre une bague à la femme. & \all{Der Mann schenkt der Frau einen Ring.} & ordre normal : sujet d'abord \\
\hline
Aujourd'hui, les jeunes font les trois mariages. & \all{Heute machen die jungen Leute alle drei Hochzeiten.} & inversion après \all{heute} \\
\hline
\end{tabularx}
\end{center}

\begin{attention}[Le piège de l'ordre des mots]
En français on dit : « Au Cameroun, \emph{il y a} trois mariages » — le sujet reste avant le verbe. En allemand, c'est impossible : après un complément en tête, le verbe doit passer devant le sujet. Ne dites jamais \all{In Kamerun es gibt drei Hochzeiten} : c'est la faute la plus fréquente des francophones. Dites : \all{In Kamerun gibt es drei Hochzeiten.} Astuce : comptez les positions — position 1 : le complément ; position 2 : le verbe ; position 3 : le sujet.
\end{attention}

\courssub{Vrai ou faux : justifier par une citation}

L'exercice « vrai ou faux » (\all{Richtig oder falsch?}) exige une justification : il faut \textbf{citer la phrase du texte} qui prouve votre réponse.

\begin{exoresolu}[Richtig oder falsch?]
\textbf{Énoncé.} Dites si les phrases sont vraies (\all{richtig}) ou fausses (\all{falsch}) et justifiez avec une citation du texte.

1. \all{In Deutschland ist die kirchliche Hochzeit ein Gesetz.}\\
2. \all{In Kamerun gibt es nur eine Hochzeit.}\\
3. \all{Die Hochzeit ist in beiden Ländern ein großes Fest.}
\tcblower
\textbf{Solution détaillée.}

1. \all{Falsch.} Le texte dit : „\all{Viele Paare machen danach auch eine kirchliche Hochzeit in der Kirche, aber das ist nur eine Tradition, kein Gesetz.}“ — Le mariage religieux est une tradition, pas une loi.

2. \all{Falsch.} Le texte dit : „\all{Hier gibt es drei Arten von Hochzeiten.}“ — Il y a trois types de mariage au Cameroun, pas un seul.

3. \all{Richtig.} Le texte dit : „\all{In beiden Ländern ist die Hochzeit ein großes Fest.}“ — C'est la phrase exacte du dernier paragraphe.
\end{exoresolu}

\courssub{Questions de vocabulaire en contexte}

On vous demandera souvent de retrouver un mot du texte à partir d'un synonyme, d'un contraire ou d'une définition. Voici comment procéder.

\begin{exemplebox}[Trouver le mot dans le texte]
\textbf{Synonyme :} Trouve un synonyme de \all{die Eheschließung} (le mariage).\\
Réponse : \all{die Hochzeit} ou \all{die Trauung}.

\textbf{Contraire :} Trouve le contraire de \all{ungültig} (non valable).\\
Réponse : \all{gültig} — „\all{nur eine gültige Hochzeit}“.

\textbf{Définition :} Quel mot signifie « l'argent ou les biens que la famille de l'homme donne à la famille de la femme » ?\\
Réponse : \all{die Mitgift} — „\all{Der Mann muss eine Mitgift bezahlen.}“

\textbf{Définition :} Quel mot signifie « la cérémonie avant le mariage, avec un anneau » ?\\
Réponse : \all{die Verlobung}.
\end{exemplebox}

\begin{methode}[Retrouver un mot dans le texte]
\begin{enumerate}
\item Traduire mentalement le mot français ou la définition donnée.
\item Balayer le texte paragraphe par paragraphe en cherchant un mot de la même famille.
\item Vérifier que le mot trouvé convient en le remettant dans sa phrase.
\item Recopier le mot \textbf{avec son article} : on écrit \all{die Mitgift}, pas seulement \all{Mitgift}.
\end{enumerate}
\end{methode}

\courssub{Questions d'opinion personnelle}

Enfin, l'expression personnelle : on vous demande votre avis. Il n'y a pas de « bonne » réponse, mais il y a une \textbf{bonne forme} : donner son opinion avec \all{ich finde...} ou \all{ich denke...} et justifier avec \all{weil} (parce que). Attention : après \all{weil}, le verbe conjugué va \textbf{à la fin} de la phrase.

\begin{exemplebox}[Donner son opinion]
\textbf{Frage :} \all{Welche Hochzeit findest du besser? Warum?}\\
« Quel mariage préfères-tu ? Pourquoi ? »

\textbf{Antwort 1 :} \all{Ich finde die traditionelle Hochzeit besser, weil die beiden Familien zusammen feiern.}\\
« Je préfère le mariage traditionnel parce que les deux familles fêtent ensemble. »

\textbf{Antwort 2 :} \all{Ich finde die standesamtliche Hochzeit wichtig, weil sie offiziell ist.}\\
« Je trouve le mariage civil important parce qu'il est officiel. »

\textbf{Antwort 3 :} \all{Ich denke, alle drei Hochzeiten sind schön, weil jede Hochzeit eine andere Tradition hat.}\\
« Je pense que les trois mariages sont beaux parce que chaque mariage a une tradition différente. »
\end{exemplebox}

\begin{aretenir}[L'essentiel de la compréhension de texte]
\begin{itemize}
\item Répondre toujours par une \textbf{phrase complète} qui reprend les mots de la question.
\item Ne jamais recopier tout le paragraphe : repérer la \textbf{phrase-clé} et la reformuler.
\item Justifier les réponses « vrai/faux » par une \textbf{citation} entre guillemets „…“.
\item Respecter la \textbf{place du verbe} : 2\textsuperscript{e} position, avec inversion si la réponse commence par un complément (\all{In Kamerun gibt es...}).
\item Pour l'opinion : \all{ich finde / ich denke} + justification avec \all{weil} + verbe à la fin.
\end{itemize}
\end{aretenir}

\begin{exercice}[À toi de jouer]
Réponds par des phrases complètes aux questions suivantes :
\begin{enumerate}
\item \all{Wo unterschreibt das Paar in Deutschland ein Dokument?}
\item \all{Was bekommen die Ehefrau und der Ehemann?}
\item \all{Welche Hochzeit machen die jungen Leute in Yaoundé zuerst?}
\item Richtig oder falsch ? \all{Bei der traditionellen Hochzeit tanzen die Familien nicht.} Justifie avec une citation.
\item \all{Welche Hochzeit findest du besser? Warum?}
\end{enumerate}
\end{exercice}

\begin{corrige}[Exercice « À toi de jouer »]
\begin{enumerate}
\item \all{In Deutschland unterschreibt das Paar im Standesamt ein Dokument.} (inversion après le complément de lieu)
\item \all{Die Ehefrau und der Ehemann bekommen viele Geschenke.}
\item \all{Die jungen Leute in Yaoundé machen zuerst die traditionelle Hochzeit.}
\item \all{Falsch.} Citation : „\all{Die beiden Familien essen, trinken und tanzen zusammen.}“
\item Réponse libre, par exemple : \all{Ich finde die kirchliche Hochzeit schön, weil die Kirche eine feierliche Atmosphäre hat.}
\end{enumerate}
\end{corrige}

La compréhension du texte est maintenant assurée : vous savez dégager le sens global, retrouver les détails, justifier par des citations et donner votre avis en allemand correct. Dans la section suivante, nous approfondirons le lexique thématique du mariage et de la famille, mot par mot, avec articles et pluriels.

\coursec{Lexique thématique : die Hochzeit und die Ehe}

Dans la section précédente, nous avons lu et compris le texte \all{Heiraten in Deutschland und in Kamerun}. Nous y avons rencontré des mots comme \all{die Verlobung}, \all{das Standesamt} ou \all{die Mitgift}. Il est temps maintenant de nous arrêter sur ce vocabulaire : le connaître précisément, avec l'article et le pluriel de chaque nom, c'est la condition pour comprendre n'importe quel texte sur la famille et pour pouvoir parler soi-même du mariage en allemand.

\courssub{Observer avant d'apprendre : un mini-texte}

Lisons d'abord ce court texte original. Soulignez mentalement tous les mots qui appartiennent au thème du mariage.

\begin{textebox}[Eine Hochzeit in Yaoundé]
Am Samstag heiratet meine Cousine Sandrine ihren Freund Thomas. Das Brautpaar hat sich vor einem Jahr verlobt: bei der Verlobung hat Thomas seiner Braut einen Ring geschenkt, und die Familie des Bräutigams hat die Mitgift gezahlt. Das ist bei uns in Kamerun eine wichtige Tradition.

Zuerst findet die standesamtliche Trauung im Rathaus statt, denn nur diese Ehe ist vor dem Gesetz gültig. Danach geht das Paar in die Kirche: die kirchliche Trauung beginnt um elf Uhr. Die Braut trägt ein weißes Kleid, der Bräutigam einen schwarzen Anzug.

Am Abend gibt es eine große Feier mit vielen Gästen. Die Familien bringen Geschenke mit, es gibt Musik und Tanz. Am nächsten Tag sagt man nicht mehr „das Brautpaar“, sondern „das Ehepaar“: Sandrine ist jetzt die Ehefrau von Thomas, und Thomas ist ihr Ehemann. Alle wünschen dem Ehepaar eine glückliche Ehe!
\end{textebox}

\textbf{Glossaire :} \all{die Cousine, -n} : la cousine ; \all{gültig} : valable ; \all{tragen} (il porte : \all{er trägt}) : porter (vêtement) ; \all{der Anzug, -e} : le costume ; \all{wünschen} : souhaiter ; \all{glücklich} : heureux.

\begin{exercice}[Questions de compréhension rapide]
\begin{enumerate}
\item \all{Wann findet die standesamtliche Trauung statt?}
\item \all{Wer zahlt die Mitgift?}
\item \all{Was trägt die Braut?}
\end{enumerate}
\end{exercice}

\begin{corrige}[Exercice « Questions de compréhension rapide »]
\begin{enumerate}
\item \all{Die standesamtliche Trauung findet zuerst im Rathaus statt.} (Elle a lieu d'abord à la mairie.)
\item \all{Die Familie des Bräutigams zahlt die Mitgift.} (La famille du marié paie la dot.)
\item \all{Die Braut trägt ein weißes Kleid.} (La mariée porte une robe blanche.)
\end{enumerate}
\end{corrige}

\courssub{Champ lexical 1 : les personnes}

\begin{center}
\begin{tabularx}{\linewidth}{|X|X|X|}\hline
\rowcolor{popblueL} Deutsch (article + pluriel) & Français & Remarque \\ \hline
der Ehemann, die Ehemänner & le mari (époux) &  \\ \hline
die Ehefrau, die Ehefrauen & la femme (épouse) &  \\ \hline
der Bräutigam, die Bräutigame & le fiancé, le marié & le jour du mariage \\ \hline
die Braut, die Bräute & la fiancée, la mariée &  \\ \hline
das Brautpaar, die Brautpaare & les jeunes mariés & le jour de la noce \\ \hline
das Ehepaar, die Ehepaare & le couple marié & dans la vie quotidienne \\ \hline
der Gast, die Gäste & l'invité &  \\ \hline
die Familie, die Familien & la famille &  \\ \hline
\end{tabularx}
\end{center}

\begin{exemplebox}[Les personnes en contexte]
\all{Der Bräutigam und die Braut stehen vor dem Standesamt.} « Le marié et la mariée sont devant l'état civil. »

\all{Die Gäste gratulieren dem Ehepaar.} « Les invités félicitent le couple marié. » (\cle{gratulieren} + datif)

\all{Ihre Familie kommt aus Douala.} « Sa famille vient de Douala. »
\end{exemplebox}

\courssub{Champ lexical 2 : les étapes du mariage}

\begin{center}
\begin{tabularx}{\linewidth}{|X|X|X|}\hline
\rowcolor{popblueL} Deutsch (article + pluriel) & Français & Remarque / Construction \\ \hline
die Verlobung, -en & les fiançailles & \cle{sich verloben} (+ \cle{mit} + datif) \\ \hline
die Trauung, -en & la cérémonie de mariage & civile ou religieuse \\ \hline
die Hochzeit, -en & le mariage (la noce, l'événement) & le jour de la fête \\ \hline
heiraten & épouser & \cle{heiraten} + \cle{accusatif} (sans préposition) \\ \hline
die Ehe, -n & le mariage (l'union, l'état) & l'état durable \\ \hline
die Scheidung, -en & le divorce & \cle{sich scheiden lassen} \\ \hline
\end{tabularx}
\end{center}

\begin{propriete}[Deux mots pour « mariage » : \cle{Hochzeit} ou \cle{Ehe} ?]
\all{die Hochzeit} désigne l'\textbf{événement}, la fête du jour J : \all{Die Hochzeit war schön.} « La noce était belle. »

\all{die Ehe} désigne l'\textbf{état durable}, l'union dans le temps : \all{Ihre Ehe ist glücklich.} « Leur mariage est heureux. »

Comparaison : le français « mariage » couvre les deux sens ; l'allemand distingue nettement le jour (\all{Hochzeit}) et la durée (\all{Ehe}).
\end{propriete}

\begin{attention}[Piège : \cle{heiraten} se construit SANS préposition]
En français on dit « épouser quelqu'un » ou « se marier \emph{avec} quelqu'un ». En allemand, \all{heiraten} prend directement l'\textbf{accusatif}, jamais la préposition \all{mit} :

\all{Er heiratet seine Freundin.} « Il épouse sa petite amie. »

\all{Sie heiratet einen Lehrer.} « Elle épouse un enseignant. »

Faux : \emph{Er heiratet mit seiner Freundin.} — La préposition \all{mit} existe seulement avec \all{sich verloben} : \all{Er verlobt sich mit seiner Freundin.} « Il se fiance avec sa petite amie. »
\end{attention}

\begin{popfigure}
\begin{tikzpicture}[node distance=0.7cm]
\node[draw=popblue, fill=popblueL, rounded corners, minimum width=3.2cm, minimum height=1.1cm, align=center] (a) {\textbf{die Verlobung}\\les fiançailles};
\node[draw=popgreen, fill=popgreenL, rounded corners, minimum width=3.2cm, minimum height=1.1cm, align=center, right=of a] (b) {\textbf{die standesamtliche}\\\textbf{Trauung}\\mariage civil};
\node[draw=poporange, fill=poporangeL, rounded corners, minimum width=3.2cm, minimum height=1.1cm, align=center, right=of b] (c) {\textbf{die kirchliche}\\\textbf{Trauung}\\mariage religieux};
\node[draw=poppurple, fill=poppurpleL, rounded corners, minimum width=3.2cm, minimum height=1.1cm, align=center, right=of c] (d) {\textbf{die Hochzeitsfeier}\\fête de mariage};
\draw[-{Stealth}, very thick, popdark] (a) -- (b);
\draw[-{Stealth}, very thick, popdark] (b) -- (c);
\draw[-{Stealth}, very thick, popdark] (c) -- (d);
\end{tikzpicture}
\legende{Frise chronologique : les étapes du mariage, des fiançailles à la fête.}
\end{popfigure}

\courssub{Champ lexical 3 : les lieux et les institutions}

\begin{center}
\begin{tabularx}{\linewidth}{|X|X|X|}\hline
\rowcolor{popblueL} Deutsch (article + pluriel) & Français & Remarque \\ \hline
das Standesamt, die Standesämter & l'état civil, le bureau des mariages &  \\ \hline
das Rathaus, die Rathäuser & la mairie &  \\ \hline
die Kirche, -n & l'église &  \\ \hline
standesamtlich & civil (de l'état civil) & adjectif / adverbe \\ \hline
kirchlich & religieux (de l'église) & adjectif / adverbe \\ \hline
traditionell & traditionnel & adjectif / adverbe \\ \hline
\end{tabularx}
\end{center}

\begin{exemplebox}[Les lieux en contexte]
\all{Das Paar heiratet standesamtlich im Rathaus.} « Le couple se marie civilement à la mairie. »

\all{Die kirchliche Trauung findet in der Kirche statt.} « La cérémonie religieuse a lieu à l'église. »

\all{In Kamerun gibt es auch die traditionelle Hochzeit.} « Au Cameroun, il y a aussi le mariage traditionnel. »
\end{exemplebox}

\begin{propriete}[Former des adverbes : très facile en allemand !]
Contrairement au français (civil \textrightarrow{} civile\textbf{ment}), l'allemand utilise l'adjectif tel quel comme adverbe, sans terminaison :

\all{Er heiratet standesamtlich.} « Il se marie civilement. »

\all{Sie feiern traditionell.} « Ils fêtent traditionnellement. »
\end{propriete}

\courssub{Champ lexical 4 : les objets et les coutumes}

\begin{center}
\begin{tabularx}{\linewidth}{|X|X|X|}\hline
\rowcolor{popblueL} Deutsch (article + pluriel) & Français & Remarque \\ \hline
der Ring, -e & la bague, l'anneau &  \\ \hline
das Geschenk, -e & le cadeau & verbe : \cle{schenken} (qqc à qqn : dat + akk) \\ \hline
die Mitgift, -en & la dot & coutume camerounaise \\ \hline
das Kleid, -er & la robe & pluriel : \all{die Kleider} \\ \hline
die Feier, -n & la fête, la cérémonie & verbe : \cle{feiern} \\ \hline
das Fest, -e & la fête & synonyme de \all{Feier} \\ \hline
\end{tabularx}
\end{center}

\begin{definition}[die Mitgift]
\all{die Mitgift, -en} : la dot, c'est-à-dire les biens ou l'argent versés par la famille du marié à la famille de la mariée dans le mariage traditionnel camerounais.

\all{Die Familie des Bräutigams zahlt die Mitgift.} « La famille du marié paie la dot. »
\end{definition}

\begin{exemplebox}[Expressions utiles à réemployer]
\all{ein Paar} : un couple — \all{Das Paar wohnt in Yaoundé.} « Le couple habite à Yaoundé. »

\all{vor dem Gesetz gültig} : valable devant la loi — \all{Nur die standesamtliche Ehe ist vor dem Gesetz gültig.} « Seul le mariage civil est valable devant la loi. »

\all{einen Ring schenken} : offrir une bague — \all{Er schenkt seiner Braut einen Ring.} « Il offre une bague à sa fiancée. » (Dat + Akk)

\all{eine Mitgift zahlen} : payer une dot — \all{Seine Familie zahlt eine Mitgift.} « Sa famille paie une dot. »
\end{exemplebox}

\courssub{Familles de mots : un mot en cache trois !}

Apprendre les familles de mots multiplie votre vocabulaire sans effort. Observez la dérivation nominale et adjectivale :

\begin{center}
\begin{tabularx}{\linewidth}{|X|X|X|}\hline
\rowcolor{popblueL} Verbe & Nom (article + pluriel) & Autres dérivés \\ \hline
heiraten & die Heirat, -en (le mariage, action) & der Heiratsantrag, -e (demande en mariage) \\ \hline
--- & die Ehe, -n (le mariage, état) & ehelich (conjugal), der Ehemann, die Ehefrau \\ \hline
feiern & die Feier, -n (la fête) & das Fest, -e (la fête) \\ \hline
schenken & das Geschenk, -e (le cadeau) & --- \\ \hline
\end{tabularx}
\end{center}

\begin{exemplebox}[La famille de « heiraten » en action]
\all{Er macht seiner Freundin einen Heiratsantrag.} « Il fait une demande en mariage à sa petite amie. » (\cle{machen} + dat + akk)

\all{Die Heirat findet im Juni statt.} « Le mariage a lieu en juin. »

\all{Die Verlobung findet vor der Hochzeit statt.} « Les fiançailles ont lieu avant le mariage. » (\cle{vor} + datif)
\end{exemplebox}

\begin{savaistu}[\cle{Brautpaar} ou \cle{Ehepaar} ?]
En allemand, \all{das Brautpaar} désigne les jeunes mariés \textbf{le jour de la noce} : la femme porte encore la robe, l'homme le costume. Dès le lendemain, on dit \all{das Ehepaar} pour parler du couple marié dans la vie quotidienne. Le français « jeunes mariés » ne fait pas cette distinction de vocabulaire !
\end{savaistu}

\begin{methode}[Comment mémoriser le vocabulaire allemand]
\begin{enumerate}
\item Apprends toujours chaque nom \textbf{avec son article et son pluriel} : pas seulement \all{Ring}, mais \all{der Ring, -e}.
\item Range les mots dans des \textbf{champs lexicaux} (personnes, étapes, lieux, objets) : le cerveau retient mieux les familles que les listes désordonnées.
\item Fabrique une \textbf{phrase personnelle} avec chaque mot : \all{Mein Onkel ist mein Lieblingsgast bei jeder Hochzeit.}
\item Teste-toi oralement : un camarade dit le mot français, tu donnes l'allemand avec article et pluriel.
\end{enumerate}
\end{methode}

\begin{exoresolu}[Réemploi du vocabulaire]
Énoncé : complétez les phrases avec le mot qui convient (article compris) : \all{die Mitgift}, \all{das Standesamt}, \all{der Bräutigam}, \all{die Verlobung}, \all{die Gäste}.

\begin{enumerate}
\item \all{Die Trauung findet auf \_\_\_\_\_ statt.}
\item \all{\_\_\_\_\_ schenkt der Braut einen Ring.}
\item \all{Vor der Hochzeit feiert die Familie \_\_\_\_\_.}
\item \all{In Kamerun zahlt die Familie des Mannes oft \_\_\_\_\_.}
\item \all{\_\_\_\_\_ bringen viele Geschenke zur Feier mit.}
\end{enumerate}

\tcblower
Solution détaillée :
\begin{enumerate}
\item \all{Die Trauung findet auf dem Standesamt statt.} — \all{auf dem Standesamt} : datif après la préposition \all{auf} (lieu). « La cérémonie a lieu à l'état civil. »
\item \all{Der Bräutigam schenkt der Braut einen Ring.} — sujet au nominatif ; \all{der Braut} est au datif (on offre quelque chose \emph{à} quelqu'un), \all{einen Ring} à l'accusatif. « Le marié offre une bague à la mariée. »
\item \all{Vor der Hochzeit feiert die Familie die Verlobung.} — \all{die Verlobung} à l'accusatif, complément direct de \all{feiern}. « Avant le mariage, la famille fête les fiançailles. »
\item \all{In Kamerun zahlt die Familie des Mannes oft die Mitgift.} — \all{die Mitgift} à l'accusatif, complément direct de \all{zahlen}. « Au Cameroun, la famille du mari paie souvent la dot. »
\item \all{Die Gäste bringen viele Geschenke zur Feier mit.} — \all{Die Gäste} nominatif pluriel ; \all{mitbringen} verbe à particule séparable : \all{mit} va à la fin ; \all{zur Feier} = \all{zu der Feier} (datif). « Les invités apportent beaucoup de cadeaux à la fête. »
\end{enumerate}
\end{exoresolu}

\begin{attention}[Erreurs fréquentes des francophones]
\begin{itemize}
\item Ne dites pas \emph{die Ehemann} : c'est \all{der Ehemann} (masculin), même si « mari » finit par une voyelle en français.
\item \all{das Kleid} signifie « la robe », pas « le vêtement » en général (faux ami avec l'anglais \emph{clothes}).
\item N'oubliez jamais la majuscule aux noms : \emph{die hochzeit} est faux, il faut \all{die Hochzeit}.
\item \all{die Feier} et \all{das Fest} sont proches, mais l'expression figée est \all{die Hochzeitsfeier} pour la fête de mariage.
\item \all{schenken} (offrir) construit le destinataire au \textbf{datif} et l'objet à l'\textbf{accusatif} : \all{schenken + Dativ + Akkusativ}.
\end{itemize}
\end{attention}

\begin{experience}[Quiz oral rapide et mini-production]
Par groupes de deux : l'élève A dit un mot français du champ lexical (par exemple « la dot », « la mairie », « épouser »), l'élève B répond en allemand avec article et pluriel (\all{die Mitgift, -en}). Cinq mots chacun, puis on inverse les rôles.

Ensuite, chaque élève écrit trois phrases personnelles avec trois mots du jour, par exemple : \all{Meine Tante hat im Rathaus geheiratet.} « Ma tante s'est mariée à la mairie. »
\end{experience}

\begin{aretenir}[L'essentiel du lexique]
\begin{itemize}
\item Les personnes : \all{der Ehemann}, \all{die Ehefrau}, \all{der Bräutigam}, \all{die Braut}, \all{das Ehepaar}, \all{der Gast}.
\item Les étapes : \all{die Verlobung} \textrightarrow{} \all{die Trauung} (standesamtlich / kirchlich / traditionell) \textrightarrow{} \all{die Hochzeit} \textrightarrow{} \all{die Ehe} (\textrightarrow{} \all{die Scheidung}).
\item Les lieux : \all{das Standesamt}, \all{das Rathaus}, \all{die Kirche}.
\item Les objets et coutumes : \all{der Ring}, \all{das Geschenk}, \all{die Mitgift}, \all{das Kleid}, \all{die Feier}.
\item Règle d'or : \cle{heiraten} + \cle{accusatif}, sans préposition : \all{Er heiratet seine Freundin.}
\item Construction \cle{schenken} : \all{schenken} + \cle{datif} (à qui) + \cle{accusatif} (quoi).
\end{itemize}
\end{aretenir}

\coursec{Grammaire 1 : l'accusatif et le datif}

Dans la section précédente, nous avons appris le vocabulaire du mariage : \all{die Hochzeit}, \all{der Bräutigam}, \all{die Braut}, \all{die Mitgift}, \all{das Geschenk}. Mais pour employer ces mots dans des phrases correctes, il faut savoir les \cle{décliner}. En français, on dit « le mari offre une bague à la mariée » : les articles ne changent jamais. En allemand, au contraire, l'article change de forme selon la \cle{fonction} du nom dans la phrase. C'est ce qu'on appelle les \cle{cas} : le nominatif (sujet), l'accusatif (COD) et le datif (COI). C'est la grammaire la plus importante de l'année : observons-la d'abord dans un texte.

\courssub{Observation : un texte et une phrase à analyser}

\begin{textebox}[Die Hochzeit von Stefan und Amina — Texte d'observation]
Stefan und Amina heiraten heute in Yaoundé. Am Morgen geht Stefan zum Standesamt. Er gibt dem Beamten die Dokumente, und der Beamte zeigt dem Paar das Formular. Um elf Uhr beginnt die Feier. Der Bräutigam schenkt der Braut einen Ring, und die Braut schenkt dem Bräutigam eine Uhr. Die Gäste gratulieren dem Ehepaar und bringen den Eltern Geschenke. Aminas Mutter dankt den Gästen für die vielen Geschenke. Am Nachmittag fährt die Familie mit dem Bus zum Restaurant. Dort sitzt das Ehepaar neben den Eltern. Ein Freund aus Deutschland schickt dem Paar eine Karte: „Wir wünschen euch viel Glück!“ Am Abend hilft Stefan seiner Frau in der Küche. Amina sagt: „Das Geschenk meiner Schwester gefällt mir sehr. Der Ring gehört mir jetzt!“ Alle lachen. Die Hochzeit ist ein großes Fest für die ganze Familie.
\end{textebox}

\textbf{Glossaire}

\begin{center}
\begin{tabularx}{\linewidth}{|X|X|}\hline
\rowcolor{popblueL} \textbf{Deutsch} & \textbf{Français} \\ \hline
das Standesamt, die Standesämter & l'état civil, le bureau de l'état civil \\ \hline
der Beamte, -n & le fonctionnaire \\ \hline
das Formular, -e & le formulaire \\ \hline
die Uhr, -en & la montre \\ \hline
gratulieren (+ Dat.) & féliciter \\ \hline
die Karte, -n & la carte (de vœux) \\ \hline
gefallen (+ Dat.) & plaire \\ \hline
gehören (+ Dat.) & appartenir \\ \hline
\end{tabularx}
\end{center}

\textbf{Questions d'observation}

\begin{enumerate}
\item \all{Wer schenkt einen Ring?} — Qui offre une bague ?
\item \all{Was schenkt der Bräutigam?} — Qu'offre le marié ?
\item \all{Wem schenkt der Bräutigam den Ring?} — À qui le marié offre-t-il la bague ?
\end{enumerate}

Réponses : c'est \all{der Bräutigam} (le sujet, nominatif) qui offre ; il offre \all{einen Ring} (quoi ? → accusatif) \all{der Braut} (à qui ? → datif). Remarquez : \all{die Braut} est devenue \all{der Braut}, et \all{ein Ring} est devenu \all{einen Ring}. Les articles ont changé de forme : c'est la déclinaison.

\courssub{La fonction des deux cas : COD et COI}

\begin{definition}[L'accusatif et le datif]
L'\cle{accusatif} (Akkusativ) est le cas du \cle{complément d'objet direct} (COD) : il répond aux questions \all{wen?} (qui ?) et \all{was?} (quoi ?). Le \cle{datif} (Dativ) est le cas du \cle{complément d'objet indirect} (COI) : il répond à la question \all{wem?} (à qui ?).
\end{definition}

Analysons la phrase clé du texte :

\all{Der Mann schenkt der Frau einen Ring.} « L'homme offre une bague à la femme. »

\begin{itemize}
\item \all{Der Mann} : qui offre ? → \textbf{sujet}, nominatif.
\item \all{einen Ring} : il offre \emph{quoi} ? → \textbf{COD}, accusatif.
\item \all{der Frau} : il offre \emph{à qui} ? → \textbf{COI}, datif.
\end{itemize}

\begin{popfigure}
\begin{center}
\begin{tikzpicture}[node distance=0.5cm and 0.6cm]
\node[draw=popblue, fill=popblueL, rounded corners, align=center, minimum width=2.6cm] (subj) {\textbf{Subjekt}\\ \all{Der Mann}\\ Nom. : wer?};
\node[draw=popdark, fill=popcream, rounded corners, align=center, minimum width=2.2cm, right=of subj] (verb) {\textbf{Verb}\\ \all{schenkt}};
\node[draw=popgreen, fill=popgreenL, rounded corners, align=center, minimum width=2.6cm, right=of verb] (dat) {\textbf{Dativ}\\ \all{der Frau}\\ à qui ? wem?};
\node[draw=poporange, fill=poporangeL, rounded corners, align=center, minimum width=2.6cm, right=of dat] (akk) {\textbf{Akkusativ}\\ \all{einen Ring}\\ quoi ? was?};
\draw[-{Stealth}, thick, popblue] (subj) -- (verb);
\draw[-{Stealth}, thick, popdark] (verb) -- (dat);
\draw[-{Stealth}, thick, popdark] (verb) -- (akk);
\end{tikzpicture}
\end{center}
\legende{Structure de la phrase à deux compléments : sujet (nominatif), verbe, COI (datif), COD (accusatif).}
\end{popfigure}

\begin{attention}[La grande différence avec le français]
En français, le COI est introduit par la préposition « à » : « Il offre une bague \textbf{à la} mariée ». En allemand, il n'y a \textbf{aucune préposition} : seule la forme de l'article indique la fonction. \all{der Braut} signifie à lui seul « à la mariée ». Ne traduisez jamais « à » par une préposition devant un COI !
\end{attention}

\courssub{L'accusatif : le cas du complément d'objet direct}

\begin{propriete}[Règle de l'accusatif]
Le COD répond à la question \all{wen? was?} (qui ? quoi ?). Bonne nouvelle : \textbf{seul le masculin singulier change de forme} : \all{der} $\rightarrow$ \all{den}, \all{ein} $\rightarrow$ \all{einen}. Le féminin, le neutre et le pluriel restent identiques au nominatif.
\end{propriete}

\begin{exemplebox}[Exemples à l'accusatif]
\all{Ich sehe den Ehemann.} « Je vois le mari. » (masculin : \all{der} $\rightarrow$ \all{den})

\all{Der Bräutigam kauft einen Ring.} « Le marié achète une bague. » (\all{ein} $\rightarrow$ \all{einen})

\all{Wir besuchen die Familie.} « Nous rendons visite à la famille. » (féminin : inchangé)

\all{Sie bringt das Geschenk.} « Elle apporte le cadeau. » (neutre : inchangé)

\all{Er lädt die Gäste ein.} « Il invite les invités. » (pluriel : inchangé ; verbe à particule \all{einladen} : la particule \all{ein} va à la fin)
\end{exemplebox}

\courssub{Le datif : le cas du complément d'objet indirect}

\begin{propriete}[Règle du datif]
Le COI répond à la question \all{wem?} (à qui ?). Les formes du datif sont : \all{dem} (masculin), \all{der} (féminin), \all{dem} (neutre), \all{den} au pluriel \textbf{avec un -n ajouté au nom} quand c'est possible : \all{den Gästen}, \all{den Kindern}, \all{den Frauen}. Exception : les noms dont le pluriel se termine déjà par -n ou par -s ne prennent pas de -n supplémentaire : \all{den Frauen}, \all{den Autos}.
\end{propriete}

\begin{exemplebox}[Exemples au datif]
\all{Er dankt dem Gast.} « Il remercie l'invité. » (masculin : \all{dem})

\all{Er dankt den Gästen.} « Il remercie les invités. » (pluriel : \all{den} + -n)

\all{Wir gratulieren der Braut.} « Nous félicitons la mariée. » (féminin : \all{der})

\all{Sie hilft dem Kind.} « Elle aide l'enfant. » (neutre : \all{dem})
\end{exemplebox}

\courssub{Tableaux de déclinaison}

\begin{popfigure}
\begin{center}
\begin{tabular}{|l|l|l|l|}\hline
\rowcolor{popblue!30}  & \textbf{Nominativ} & \textbf{Akkusativ} & \textbf{Dativ} \\ \hline
Maskulin & der / ein & \cellcolor{poporangeL}den / einen & \cellcolor{popgreenL}dem / einem \\ \hline
Feminin & die / eine & die / eine & \cellcolor{popgreenL}der / einer \\ \hline
Neutrum & das / ein & das / ein & \cellcolor{popgreenL}dem / einem \\ \hline
Plural & die / -- & die / -- & \cellcolor{popgreenL}den + n / -- \\ \hline
\end{tabular}
\end{center}
\legende{Articles définis et indéfinis. En orange : la seule forme d'accusatif qui change (masculin). En vert : les formes du datif, toutes différentes du nominatif.}
\end{popfigure}

\begin{center}
\begin{tabularx}{\linewidth}{|X|X|X|X|}\hline
\rowcolor{popblueL} \textbf{Cas} & \textbf{Masculin} & \textbf{Féminin} & \textbf{Neutre / Pluriel} \\ \hline
Nominatif & der Vater / ein Vater & die Mutter / eine Mutter & das Kind / die Kinder \\ \hline
Akkusativ & \textbf{den} Vater / \textbf{einen} Vater & die Mutter / eine Mutter & das Kind / die Kinder \\ \hline
Dativ & \textbf{dem} Vater / \textbf{einem} Vater & \textbf{der} Mutter / \textbf{einer} Mutter & \textbf{dem} Kind / \textbf{den} Kinder\textbf{n} \\ \hline
\end{tabularx}
\end{center}

\begin{propriete}[Les possessifs se déclinent comme \all{ein}]
Les adjectifs possessifs (\all{mein}, \all{dein}, \all{sein}, \all{ihr}, \all{unser}...) suivent la même déclinaison que l'article indéfini : \all{mein Bruder} $\rightarrow$ accusatif \all{meinen Bruder}, datif \all{meinem Bruder} ; \all{meine Schwester} $\rightarrow$ datif \all{meiner Schwester}.

\all{Stefan hilft seiner Frau.} « Stefan aide sa femme. » (datif féminin : \all{sein} $\rightarrow$ \all{seiner})

\all{Das Geschenk meiner Schwester gefällt mir.} « Le cadeau de ma sœur me plaît. »
\end{propriete}

\begin{center}
\begin{tabularx}{\linewidth}{|X|X|X|}\hline
\rowcolor{popblueL} \textbf{Nominatif} & \textbf{Akkusativ} & \textbf{Dativ} \\ \hline
ich (je) & mich (me) & mir (à moi) \\ \hline
du (tu) & dich (te) & dir (à toi) \\ \hline
er (il) & ihn (le) & ihm (à lui) \\ \hline
sie (elle) & sie (la) & ihr (à elle) \\ \hline
es (il, neutre) & es (le) & ihm (à lui) \\ \hline
wir (nous) & uns (nous) & uns (à nous) \\ \hline
ihr (vous) & euch (vous) & euch (à vous) \\ \hline
sie (ils/elles) & sie (les) & ihnen (à eux) \\ \hline
Sie (vous, politesse) & Sie (vous) & Ihnen (à vous) \\ \hline
\end{tabularx}
\end{center}

\begin{exemplebox}[Pronoms en contexte]
\all{Ich danke dir für das Geschenk.} « Je te remercie pour le cadeau. »

\all{Die Braut sieht ihn und lächelt.} « La mariée le voit et sourit. »

\all{Wir gratulieren Ihnen, Herr Mbarga!} « Nous vous félicitons, Monsieur Mbarga ! »

\all{Die Mutter gibt ihm eine Uhr.} « La mère lui donne une montre. »
\end{exemplebox}

\courssub{Quels verbes, quel cas ?}

\textbf{1. Les verbes à double complément} (quelque chose \emph{à} quelqu'un) : la \emph{chose} est à l'accusatif, la \emph{personne} au datif.

\begin{center}
\begin{tabularx}{\linewidth}{|X|X|X|}\hline
\rowcolor{popblueL} \textbf{Verbe} & \textbf{Exemple} & \textbf{Traduction} \\ \hline
geben (donner) & \all{Er gibt dem Beamten die Dokumente.} & Il donne les documents au fonctionnaire. \\ \hline
schenken (offrir) & \all{Der Bräutigam schenkt der Braut einen Ring.} & Le marié offre une bague à la mariée. \\ \hline
zeigen (montrer) & \all{Sie zeigt den Gästen das Foto.} & Elle montre la photo aux invités. \\ \hline
bringen (apporter) & \all{Die Gäste bringen den Eltern Geschenke.} & Les invités apportent des cadeaux aux parents. \\ \hline
schicken (envoyer) & \all{Ein Freund schickt dem Paar eine Karte.} & Un ami envoie une carte au couple. \\ \hline
\end{tabularx}
\end{center}

\begin{propriete}[Ordre des deux compléments]
Quand les deux compléments sont des \textbf{noms}, le datif (personne) précède l'accusatif (chose) : \all{Der Bräutigam schenkt \textbf{der Braut} \textbf{einen Ring}.} Quand l'accusatif est un \textbf{pronom}, il passe devant : \all{Der Bräutigam schenkt \textbf{ihn} der Braut.} « Le marié le offre à la mariée » (= la bague).
\end{propriete}

\textbf{2. Les verbes suivis du datif seul} (à apprendre par cœur : ils n'ont pas de COD) :

\begin{center}
\begin{tabularx}{\linewidth}{|X|X|X|}\hline
\rowcolor{popblueL} \textbf{Verbe} & \textbf{Exemple} & \textbf{Traduction} \\ \hline
danken (remercier) & \all{Aminas Mutter dankt den Gästen.} & La mère d'Amina remercie les invités. \\ \hline
helfen (aider) & \all{Stefan hilft seiner Frau.} & Stefan aide sa femme. \\ \hline
gratulieren (féliciter) & \all{Wir gratulieren dem Ehepaar.} & Nous félicitons le couple marié. \\ \hline
gefallen (plaire) & \all{Das Geschenk gefällt mir.} & Le cadeau me plaît. \\ \hline
gehören (appartenir) & \all{Der Ring gehört mir jetzt.} & La bague m'appartient maintenant. \\ \hline
\end{tabularx}
\end{center}

\begin{attention}[Piège : \all{helfen} et \all{danken}]
En français, on « aide quelqu'un » et on « remercie quelqu'un » : ce sont des COD. En allemand, \all{helfen} et \all{danken} exigent le \textbf{datif} : \all{Ich helfe \textbf{dem Mann}} (et non \all{den Mann}). Même logique pour \all{gratulieren}, \all{gefallen}, \all{gehören}. Attention aussi au sens de \all{gefallen} : \all{Das Geschenk gefällt mir} signifie « le cadeau me plaît » — le sujet allemand est la chose, la personne est au datif.
\end{attention}

\courssub{Les prépositions et les cas}

\begin{propriete}[Prépositions à cas fixe]
Certaines prépositions exigent toujours l'\textbf{accusatif} : \all{für} (pour), \all{ohne} (sans), \all{gegen} (contre), \all{durch} (à travers), \all{um} (autour de). D'autres exigent toujours le \textbf{datif} : \all{mit} (avec), \all{von} (de), \all{zu} (à, chez), \all{bei} (chez, près de), \all{nach} (à, après), \all{aus} (de, hors de).
\end{propriete}

\begin{propriete}[Prépositions à double cas : \all{in}, \all{an}, \all{auf}, \all{neben}]
Avec les prépositions de lieu, le cas dépend du sens : \textbf{accusatif = mouvement} (où va-t-on ? \all{wohin?}), \textbf{datif = position} (où est-on ? \all{wo?}).

\all{Stefan geht \textbf{zum} Standesamt.} (\all{zu dem}, datif) « Stefan va à l'état civil. »

\all{Die Familie fährt in\textbf{s} Restaurant.} (\all{in das}, mouvement → accusatif) « La famille va au restaurant. »

\all{Das Ehepaar sitzt \textbf{in dem} Restaurant.} (position → datif) « Le couple est assis dans le restaurant. »

\all{Das Ehepaar sitzt neben \textbf{den} Eltern.} (position → datif pluriel) « Le couple est assis à côté des parents. »
\end{propriete}

\begin{exemplebox}[Prépositions en contexte]
\all{Das Geschenk ist für die Ehefrau.} « Le cadeau est pour l'épouse. » (\all{für} + accusatif)

\all{Er kommt mit der Familie.} « Il vient avec la famille. » (\all{mit} + datif)

\all{Ein Freund aus Deutschland schickt eine Karte.} « Un ami d'Allemagne envoie une carte. » (\all{aus} + datif)

\all{Ohne den Ring ist die Feier nicht komplett.} « Sans la bague, la fête n'est pas complète. » (\all{ohne} + accusatif)
\end{exemplebox}

\begin{aretenir}[Contractions utiles]
Certaines prépositions fusionnent avec l'article : \all{zu dem} $\rightarrow$ \all{zum}, \all{zu der} $\rightarrow$ \all{zur}, \all{an dem} $\rightarrow$ \all{am}, \all{in dem} $\rightarrow$ \all{im}, \all{in das} $\rightarrow$ \all{ins}, \all{an das} $\rightarrow$ \all{ans}, \all{für das} $\rightarrow$ \all{fürs}. Exemple du texte : \all{Am Morgen} = \all{an dem Morgen} « le matin ».
\end{aretenir}

\courssub{Comparaison français / allemand}

\begin{center}
\begin{tabularx}{\linewidth}{|X|X|X|}\hline
\rowcolor{popblueL} \textbf{Français} & \textbf{Deutsch} & \textbf{Remarque} \\ \hline
Je vois le mari. & \all{Ich sehe den Mann.} & COD : masculin \all{der} $\rightarrow$ \all{den} \\ \hline
Il offre une bague à la mariée. & \all{Er schenkt der Braut einen Ring.} & « à la » = datif seul, sans préposition \\ \hline
Il remercie les invités. & \all{Er dankt den Gästen.} & \all{danken} + datif ; -n au pluriel \\ \hline
Le cadeau est pour la femme. & \all{Das Geschenk ist für die Frau.} & \all{für} + accusatif (féminin inchangé) \\ \hline
Elle vient avec les enfants. & \all{Sie kommt mit den Kindern.} & \all{mit} + datif pluriel + -n \\ \hline
\end{tabularx}
\end{center}

\begin{methode}[Comment choisir entre accusatif et datif ?]
\begin{enumerate}
\item Trouve le \textbf{verbe} de la phrase.
\item Pose la question au verbe : « qui ? quoi ? » (\all{wen? was?}) → \textbf{accusatif} ; « à qui ? » (\all{wem?}) → \textbf{datif}.
\item Si le verbe fait partie de la liste des verbes à datif (\all{danken}, \all{helfen}, \all{gratulieren}, \all{gefallen}, \all{gehören}), mets le complément au datif.
\item S'il y a une \textbf{préposition}, vérifie son cas dans la liste apprise (\all{für}, \all{ohne} + accusatif ; \all{mit}, \all{von}, \all{zu} + datif).
\item Au \textbf{pluriel du datif}, n'oublie pas le -n final : \all{den Gästen}, \all{mit den Kindern}.
\end{enumerate}
\end{methode}

\begin{exoresolu}[Complète les terminaisons (Ergänzen Sie die Endungen)]
Énoncé. Complète avec la bonne forme de l'article :

1. \all{Der Bräutigam schenkt d--- Braut ein--- Ring.}

2. \all{Wir gratulieren d--- Ehepaar.}

3. \all{Das Geschenk ist für d--- Ehemann.}

4. \all{Sie kommt mit d--- Kinder---.}

5. \all{Ich helfe d--- Mutter.}
\tcblower
Solution détaillée.

1. « Qui offre ? » \all{der Bräutigam} (sujet). « Il offre à qui ? » \all{d\textbf{er} Braut} → datif féminin. « Il offre quoi ? » \all{ein\textbf{en} Ring} → accusatif masculin. Phrase : \all{Der Bräutigam schenkt d\textbf{er} Braut ein\textbf{en} Ring.}

2. \all{gratulieren} exige le datif ; \all{das Ehepaar} est neutre → \all{d\textbf{em} Ehepaar}.

3. \all{für} exige l'accusatif ; \all{der Ehemann} est masculin → \all{für d\textbf{en} Ehemann}.

4. \all{mit} exige le datif ; pluriel → \all{mit d\textbf{en} Kinder\textbf{n}} (ne pas oublier le -n du datif pluriel).

5. \all{helfen} exige le datif ; \all{die Mutter} est féminin → \all{Ich helfe d\textbf{er} Mutter.}
\end{exoresolu}

\begin{attention}[Erreurs fréquentes des francophones]
\begin{itemize}
\item Oublier le \textbf{-n du datif pluriel} : on dit \all{mit den Kindern} et non \all{mit den Kinder}.
\item Dire \all{den Frau} au lieu de \all{der Frau} au datif féminin : \all{Er dankt \textbf{der} Frau} (il remercie la femme).
\item Traduire « à » par une préposition devant un COI : « Il donne le livre à son frère » se dit \all{Er gibt \textbf{seinem Bruder} das Buch}, sans \all{zu}.
\item Croire que tout change à l'accusatif : seul le masculin singulier change (\all{der} $\rightarrow$ \all{den}) ; \all{die}, \all{das} et le pluriel ne changent pas.
\item Confondre \all{der} datif féminin et \all{der} nominatif masculin : dans \all{Er hilft der Frau}, \all{der} n'est pas un masculin mais bien le datif de \all{die Frau}.
\end{itemize}
\end{attention}

\begin{exercice}[À ton tour]
1. Mets à l'accusatif : \all{der Ring}, \all{die Uhr}, \all{das Geschenk}, \all{ein Freund}.

2. Mets au datif : \all{die Braut}, \all{das Kind}, \all{die Gäste}, \all{der Ehemann}.

3. Traduis en allemand : « Le marié offre une montre à la mariée. » / « Nous aidons les parents. » / « La carte est pour le couple. »

4. Choisis le bon cas : \all{Er kommt mit (die / der / den) Familie.} — \all{Ich sehe (der / den / dem) Bräutigam.}
\end{exercice}

\begin{corrige}[Exercice « À ton tour »]
1. \all{den Ring}, \all{die Uhr} (inchangé), \all{das Geschenk} (inchangé), \all{einen Freund}.

2. \all{der Braut}, \all{dem Kind}, \all{den Gästen} (-n au pluriel), \all{dem Ehemann}.

3. \all{Der Bräutigam schenkt der Braut eine Uhr.} / \all{Wir helfen den Eltern.} (\all{helfen} + datif) / \all{Die Karte ist für das Paar.} (\all{für} + accusatif, neutre inchangé).

4. \all{Er kommt mit \textbf{der} Familie.} (\all{mit} + datif féminin) — \all{Ich sehe \textbf{den} Bräutigam.} (accusatif masculin).
\end{corrige}

\begin{savaistu}[L'accusatif et le datif au Bac]
Au baccalauréat, la question « \all{Ergänzen Sie die Endungen} » (complétez les terminaisons) porte presque toujours sur l'accusatif et le datif. Maîtriser le tableau de déclinaison, c'est gagner des points faciles ! Astuce de correcteur : apprends les formes du datif en chantant \all{dem – der – dem – den + n}. En Allemagne, les enfants apprennent eux aussi ces cas à l'école primaire avec les questions \all{wer? wen? wem?} : si tu sais poser la bonne question, tu trouves toujours le bon cas.
\end{savaistu}

\begin{aretenir}[L'essentiel de la section]
\begin{itemize}
\item L'\cle{accusatif} = COD (\all{wen? was?}) ; seul le masculin change : \all{der} $\rightarrow$ \all{den}, \all{ein} $\rightarrow$ \all{einen}.
\item Le \cle{datif} = COI (\all{wem?}) ; formes : \all{dem – der – dem – den + n} au pluriel.
\item Verbes à double complément (\all{geben}, \all{schenken}, \all{zeigen}, \all{bringen}, \all{schicken}) : la personne au datif, la chose à l'accusatif ; le datif précède l'accusatif quand ce sont deux noms.
\item Verbes à datif seul : \all{danken}, \all{helfen}, \all{gratulieren}, \all{gefallen}, \all{gehören}.
\item Prépositions + accusatif : \all{für}, \all{ohne}, \all{gegen}, \all{durch}, \all{um} ; + datif : \all{mit}, \all{von}, \all{zu}, \all{bei}, \all{nach}, \all{aus} ; \all{in}, \all{an}, \all{auf}, \all{neben} : accusatif = mouvement, datif = position.
\item Le français dit « à quelqu'un », l'allemand décline sans préposition.
\end{itemize}
\end{aretenir}

\coursec{Grammaire 2 : les verbes de modalité au présent}

Dans la section précédente, nous avons appris à distinguer l'accusatif et le datif, c'est-à-dire à identifier \emph{qui} fait l'action (sujet, nominatif), \emph{qui} la subit directement (objet direct, accusatif) et \emph{à qui} elle est destinée (objet indirect, datif). Mais dans le texte \all{Heiraten in Deutschland und in Kamerun}, un autre phénomène grammatical saute aux yeux : presque chaque phrase contient un petit verbe qui exprime une \emph{obligation}, une \emph{possibilité}, une \emph{permission} ou une \emph{volonté}. Ce sont les \cle{verbes de modalité} (\all{die Modalverben}), indispensables pour parler des règles du mariage — et de toutes les règles de la vie sociale.

\courssub{Observation : que remarque-t-on dans le texte ?}

Reprenons deux phrases du texte support :

\begin{exemplebox}[Phrases du texte à observer]
\all{Wer möchte, kann auch in der Kirche heiraten.}\\
« Qui veut peut aussi se marier à l'église. »

\medskip

\all{Der Ehemann und die Ehefrau müssen zusammenarbeiten.}\\
« Le mari et la femme doivent travailler ensemble. »
\end{exemplebox}

Observons attentivement :

\begin{itemize}
\item Dans chaque phrase, il y a \textbf{deux verbes} : un verbe modal conjugué (\all{kann}, \all{müssen}) en deuxième position et un verbe principal à l'infinitif (\all{heiraten}, \all{zusammenarbeiten}) en position finale.
\item Le verbe modal conjugué occupe la \textbf{deuxième position} de la phrase (place normale du verbe conjugué en allemand).
\item L'infinitif du verbe principal est renvoyé \textbf{à la toute fin} de la phrase, \emph{après} tous les compléments (lieu, temps, objet) : \all{kann auch in der Kirche \textbf{heiraten}}.
\item Remarque de forme (singulier) : on dit \all{er kann} (et non \emph{er kannt}) ; \all{ich muss} (et non \emph{ich müsse}). La voyelle du radical change souvent (\all{ö} $\rightarrow$ \all{a}, \all{ü} $\rightarrow$ \all{u}) et la 3e personne du singulier \textbf{n'a pas de -t final}.
\end{itemize}

\courssub{Les six verbes de modalité et leur sens}

\begin{definition}[Les verbes de modalité]
Les \cle{verbes de modalité} (\all{die Modalverben}) sont six verbes qui n'expriment pas une action concrète en eux-mêmes, mais modifient le sens du verbe principal : ils indiquent si l'action est \emph{permise}, \emph{possible}, \emph{obligatoire}, \emph{conseillée} (par autrui) ou \emph{voulue}. Ils se combinent obligatoirement avec un infinitif (sauf \all{möchten} qui peut s'employer seul).
\end{definition}

\begin{center}
\begin{tabularx}{\linewidth}{|l|X|X|}
\hline
\rowcolor{popblueL} \textbf{Modal} & \textbf{Sens principal} & \textbf{Exemple (contexte du mariage)} \\
\hline
\all{dürfen} & permission, droit, autorisation & \all{Du darfst tanzen.} (Tu as le droit de danser.) \\
\hline
\all{können} & capacité, possibilité, savoir-faire & \all{Man kann kirchlich heiraten.} (On peut se marier à l'église.) \\
\hline
\all{müssen} & obligation, nécessité (interne ou loi) & \all{Man muss standesamtlich heiraten.} (On doit se marier civilement.) \\
\hline
\all{sollen} & devoir conseillé, ordre/attente venant d'autrui & \all{Der Bräutigam soll eine Mitgift zahlen.} (Le marié est censé payer une dot.) \\
\hline
\all{wollen} & volonté, intention ferme & \all{Das Paar will im Juni heiraten.} (Le couple veut se marier en juin.) \\
\hline
\all{mögen} / \all{möchten} & aimer / désir poli (forme de politesse) & \all{Die Braut möchte ein weißes Kleid.} (La mariée voudrait une robe blanche.) \\
\hline
\end{tabularx}
\end{center}

\begin{propriete}[Règle fondamentale : la parenthèse verbale (\all{die Satzklammer})]
Le verbe modal est \textbf{conjugué en deuxième position} de la phrase et renvoie l'\textbf{infinitif du verbe principal à la toute fin}. Le modal et l'infinitif forment une « parenthèse » (\all{Satzklammer}) qui encadre tous les compléments (temps, lieu, objet, manière).

\medskip

\all{Das Paar will im Juni in der Kirche heiraten.}\\
« Le couple veut se marier en juin à l'église. »
\end{propriete}

\begin{popfigure}
\begin{tikzpicture}[font=\small]
\node[draw=popblue, fill=popblueL, rounded corners, align=center] (subj) at (0,0) {Das Paar\\ (sujet)};
\node[draw=poporange, fill=poporangeL, rounded corners, align=center] (mod) at (3,0) {will\\ (modal conjugué,\\ 2e position)};
\node[draw=popmuted, fill=popcream, rounded corners, align=center] (compl) at (6.5,0) {im Juni in der Kirche\\ (compléments)};
\node[draw=popgreen, fill=popgreenL, rounded corners, align=center] (inf) at (10.2,0) {heiraten.\\ (infinitif,\\ dernière position)};
\draw[-{Stealth}, very thick, popgreen] (mod.south) .. controls (5,-1.8) and (8.5,-1.8) .. node[below, align=center, text=popdark]{l'infinitif est renvoyé à la fin} (inf.south);
\draw[-{Stealth}, thick, popdark] (subj) -- (mod);
\draw[-{Stealth}, thick, popdark] (mod) -- (compl);
\draw[-{Stealth}, thick, popdark] (compl) -- (inf);
\end{tikzpicture}
\legende{La parenthèse verbale (Satzklammer) : le modal conjugué en 2e position encadre les compléments et renvoie l'infinitif en fin de phrase.}
\end{popfigure}

\courssub{Conjugaison des six modaux au présent de l'indicatif}

\begin{propriete}[Formation au présent]
Au présent de l'indicatif, les modaux (sauf \all{sollen}) présentent une irrégularité au singulier : ils perdent la voyelle de leur infinitif (\all{können} $\rightarrow$ \all{ich kann}, \all{müssen} $\rightarrow$ \all{ich muss}, \all{dürfen} $\rightarrow$ \all{ich darf}, \all{wollen} $\rightarrow$ \all{ich will}, \all{mögen} $\rightarrow$ \all{ich mag}). Les 1re et 3e personnes du singulier sont \textbf{identiques} et \textbf{n'ont pas de -t final}. Au pluriel, on retrouve la forme de l'infinitif (radical + \all{-en}). \all{Sollen} est régulier : \all{ich soll, er soll}.
\end{propriete}

\begin{popfigure}
\begin{tikzpicture}[font=\footnotesize]
\matrix[matrix of nodes, 
nodes={minimum height=0.62cm, anchor=center, align=center},
column 1/.style={nodes={fill=popblueL, font=\bfseries\footnotesize}},
row 1/.style={nodes={fill=poporangeL, font=\bfseries\footnotesize}},
nodes in empty cells] (m) {
Personne & dürfen & können & müssen & sollen & wollen & mögen \\
ich & darf & kann & muss & soll & will & mag \\
du & darfst & kannst & musst & sollst & willst & magst \\
er/sie/es & darf & kann & muss & soll & will & mag \\
wir & dürfen & können & müssen & sollen & wollen & mögen \\
ihr & dürft & könnt & müsst & sollt & wollt & mögt \\
sie/Sie & dürfen & können & müssen & sollen & wollen & mögen \\
};
\foreach \i in {1,...,7} {\draw[popline] (m-1-\i.north west) -- (m-7-\i.south west);}
\draw[popline] (m-1-7.north east) -- (m-7-7.south east);
\foreach \j in {1,...,7} {\draw[popline] (m-\j-1.north west) -- (m-\j-7.north east);}
\draw[popline] (m-7-1.south west) -- (m-7-7.south east);
\end{tikzpicture}
\legende{Tableau de conjugaison des six modaux au présent de l'indicatif. Singulier : radical modifié, 1re = 3e personne sans -t. Pluriel : forme de l'infinitif.}
\end{popfigure}

\begin{aretenir}[Astuce de mémorisation : le trio « 1 = 3 sans -t »]
Retiens la 1re personne du singulier : \all{ich darf, ich kann, ich muss, ich soll, ich will, ich mag}. La 3e personne est \textbf{identique} : \all{er darf, er kann, er muss, er soll, er will, er mag}. Au pluriel, tout redevient régulier : \all{wir dürfen, wir können, wir müssen, wir sollen, wir wollen, wir mögen}.
\end{aretenir}

\begin{aretenir}[Le cas particulier de \all{möchten} (politesse)]
\all{Möchten} n'est pas un infinitif présent, mais le \textbf{Konjunktiv II} de \all{mögen}. Il se conjugue \textbf{régulièrement} au présent et s'emploit comme un modal de politesse (« voudrais / aimerait »).
\medskip
\begin{center}
\begin{tabular}{|l|l|l|l|}\hline
\rowcolor{popblueL} \textbf{Personne} & \textbf{möchten} & \textbf{Personne} & \textbf{möchten} \\ \hline
ich & möchte & wir & möchten \\ \hline
du & möchtest & ihr & möchtet \\ \hline
er/sie/es & möchte & sie/Sie & möchten \\ \hline
\end{tabular}
\end{center}
\emph{Attention :} \all{ich möchte} (pas \emph{ich möcht}), \all{du möchtest} (pas \emph{du möchtst}).
\end{aretenir}

\courssub{Emplois de chaque modal dans le contexte du mariage}

\begin{exemplebox}[Un exemple par modal, traduit]
\all{In Deutschland muss man standesamtlich heiraten.} « En Allemagne, on doit se marier civilement. » (obligation légale, \all{müssen})

\medskip

\all{Man kann danach auch kirchlich heiraten.} « On peut ensuite aussi se marier à l'église. » (possibilité, \all{können})

\medskip

\all{Der Bräutigam soll eine Mitgift zahlen.} « Le marié doit [est censé] payer une dot. » (attente familiale, tradition, \all{sollen})

\medskip

\all{Die Braut möchte ein weißes Kleid.} « La mariée voudrait une robe blanche. » (désir poli, \all{möchten} sans infinitif)

\medskip

\all{Du darfst auf der Hochzeit tanzen.} « Tu as le droit de danser au mariage. » (permission, \all{dürfen})

\medskip

\all{Wir wollen im Dezember heiraten.} « Nous voulons nous marier en décembre. » (volonté ferme, \all{wollen})
\end{exemplebox}

\courssub{Comparaison français / allemand : l'ordre des mots}

\begin{center}
\begin{tabularx}{\linewidth}{|X|X|X|}
\hline
\rowcolor{popblueL} Français & Deutsch & Remarque \\
\hline
Il peut se marier. & \all{Er kann heiraten.} & Même ordre : modal (2e) + infinitif (fin). \\
\hline
Il peut se marier à l'église. & \all{Er kann in der Kirche heiraten.} & En allemand, l'infinitif reste \textbf{dernier}, après le complément de lieu. \\
\hline
Elle veut payer la dot. & \all{Sie will die Mitgift zahlen.} & L'objet (\all{die Mitgift}) reste \textbf{avant} l'infinitif. Jamais \emph{Sie will zahlen die Mitgift}. \\
\hline
Je voudrais un café. & \all{Ich möchte einen Kaffee.} & \all{möchten} s'emploie souvent seul (verbe complet), sans infinitif. \\
\hline
\end{tabularx}
\end{center}

\begin{attention}[Piège n° 1 : la place de l'infinitif (erreur n° 1 des francophones)]
En français, l'infinitif suit directement le modal : « il peut \emph{se marier} à l'église ». En allemand, quand la phrase contient des compléments, l'infinitif reste \textbf{en dernière position absolue} : on dit \all{Er kann in der Kirche heiraten.} et \textbf{jamais} \emph{Er kann heiraten in der Kirche.}
\medskip
\textbf{Règle d'or :} Le verbe conjugué (modal) est en position 2, l'infinitif est en position finale. Entre les deux : tout le reste.
\end{attention}

\begin{attention}[Piège n° 2 : ne jamais conjuguer l'infinitif !]
L'infinitif ne change \textbf{jamais de forme}. On n'écrit jamais \emph{er kann heiratet} ou \emph{wir müssen gehen}. Toujours : \all{er kann heiraten}, \all{wir müssen gehen}.
\end{attention}

\begin{attention}[Piège n° 3 : müssen ou sollen ? will ou möchte ?]
\all{müssen} exprime une obligation personnelle, une contrainte interne ou une loi : \all{Ich muss lernen.} (Je dois apprendre, c'est nécessaire pour moi / pour l'examen).
\medskip
\all{sollen} exprime un ordre, un conseil ou une attente venant \textbf{d'autrui} (parents, prof, société) : \all{Du sollst mehr lernen!} (Tu es censé apprendre plus — dit par tes parents).
\medskip
De même, \all{ich will} = « je veux » (volonté forte, parfois perçue comme impolie ou enfantine), tandis que \all{ich möchte} = « je voudrais » (forme polie, standard pour demander). Au restaurant, à l'administration, en société : \textbf{toujours \all{möchte}} : \all{Ich möchte einen Kaffee, bitte.}
\end{attention}

\begin{methode}[Conjuguer et construire une phrase avec un modal en 4 étapes]
\begin{enumerate}
\item \textbf{Choisis le modal} selon le sens visé : permission (\all{dürfen}), possibilité (\all{können}), obligation (\all{müssen}), conseil/attente (\all{sollen}), volonté (\all{wollen}), désir poli (\all{möchten}).
\item \textbf{Conjugue-le} selon le sujet : singulier $\rightarrow$ radical modifié, 1re/3e pers. identiques sans -t (\all{ich/er kann}) ; pluriel $\rightarrow$ forme de l'infinitif (\all{wir können}).
\item \textbf{Place le modal conjugué en 2e position} (après le sujet ou un complément en tête de phrase).
\item \textbf{Mets le verbe principal à l'infinitif, à la toute fin} de la phrase (après tous les compléments).
\end{enumerate}
Exemple : « Nous devons inviter la famille. » $\rightarrow$ Modal = \all{müssen} (obligation). Sujet = \all{wir} $\rightarrow$ \all{wir müssen}. Complément = \all{die Familie}. Infinitif = \all{einladen}. $\rightarrow$ \all{Wir müssen die Familie einladen.}
\end{methode}

\begin{textebox}[Die Hochzeit von Amina und Thomas]
Amina aus Yaoundé und Thomas aus München wollen im Juli heiraten. Es gibt viel zu tun! Zuerst müssen sie zum Standesamt gehen, denn in Deutschland kann man nur standesamtlich heiraten. Danach dürfen sie auch in der Kirche heiraten, wenn sie möchten. Aminas Mutter sagt: „Thomas soll eine Mitgift zahlen, das ist bei uns Tradition!“ Thomas will das verstehen und fragt viele Fragen über Kamerun. Amina möchte ein weißes Kleid und auch ein traditionelles Kleid tragen — sie kann beides anziehen! Die Gäste sollen um 14 Uhr kommen. Am Abend dürfen alle tanzen und singen. Thomas sagt: „Wir müssen zusammenarbeiten, dann wird unsere Ehe stark.“ Amina lächelt: „Ich will dich immer unterstützen.“

\medskip

\textbf{Glossaire :} das Standesamt, die Standesämter (l'état civil) ; die Mitgift, -en (la dot) ; der Gast, die Gäste (l'invité) ; tragen (porter un vêtement) ; anziehen (s'habiller / mettre un vêtement, séparable) ; unterstützen (soutenir) ; lächeln (sourire) ; die Ehe, -n (le mariage, l'union).

\medskip

\textbf{Questions de compréhension :}
\begin{enumerate}
\item Warum müssen Amina und Thomas zum Standesamt gehen?
\item Was soll Thomas nach der Tradition zahlen?
\item Was möchte Amina tragen?
\item Was dürfen die Gäste am Abend machen?
\end{enumerate}
\end{textebox}

\begin{exoresolu}[Exercice résolu : conjugue et construis]
Énoncé — Conjugue le modal entre parenthèses et place correctement l'infinitif (ou construis la phrase si pas d'infinitif) :
\begin{enumerate}
\item Man (müssen) --- standesamtlich --- heiraten.
\item Die Gäste (dürfen) --- auf der Hochzeit --- tanzen.
\item Ich (möchten) --- ein weißes Kleid.
\item Wir (wollen) --- im Juli --- heiraten.
\item Der Bräutigam (sollen) --- die Mitgift --- zahlen.
\end{enumerate}
\tcblower
\textbf{Solution détaillée :}
\begin{enumerate}
\item \all{Man muss standesamtlich heiraten.} (Sujet \all{man} = 3e pers. sg. $\rightarrow$ \all{muss} sans -t ; infinitif \all{heiraten} à la fin.)
\item \all{Die Gäste dürfen auf der Hochzeit tanzen.} (Sujet pluriel \all{die Gäste} $\rightarrow$ \all{dürfen} comme l'infinitif ; infinitif \all{tanzen} à la fin.)
\item \all{Ich möchte ein weißes Kleid.} (\all{möchten} 1e pers. sg. $\rightarrow$ \all{ich möchte} ; s'emploie sans infinitif ici, \all{ein weißes Kleid} est l'objet accusatif.)
\item \all{Wir wollen im Juli heiraten.} (Sujet \all{wir} $\rightarrow$ \all{wollen} ; complément \all{im Juli} inside la parenthèse ; infinitif \all{heiraten} final.)
\item \all{Der Bräutigam soll die Mitgift zahlen.} (Sujet \all{der Bräutigam} = 3e pers. sg. $\rightarrow$ \all{soll} ; objet \all{die Mitgift} ; infinitif \all{zahlen} final.)
\end{enumerate}
\end{exoresolu}

\begin{savaistu}[Un proverbe pour le Bac]
\all{Eine Ehe will gelernt sein.} — « Le mariage, ça s'apprend ! » (Litt. : « Un mariage veut être appris. ») Ce proverbe allemand utilise le modal \all{wollen} au sens de « exiger / nécessiter » (impersonnel). Citer un proverbe authentique dans une production écrite au Baccalauréat est très apprécié des correcteurs : retiens-en un par thème !
\end{savaistu}

\coursec{Méthode du commentaire et production guidée}

Dans la section précédente, nous avons appris à conjuguer les verbes de modalité au présent (\all{können, müssen, dürfen, wollen, sollen, mögen}) et à les placer correctement dans la phrase : le modal en position 2, l'infinitif à la fin. Nous allons maintenant mettre toutes ces connaissances en pratique : commenter le texte \all{Heiraten in Deutschland und in Kamerun}, le résumer, puis produire un court paragraphe comparatif et un dialogue oral. C'est exactement la compétence exigée au baccalauréat : comprendre, résumer, donner son avis.

\courssub{La méthode du commentaire de texte en quatre étapes}

\begin{methode}[Commenter un texte en allemand : les 4 étapes]
Un commentaire de texte (\all{die Textbesprechung}) suit toujours le même plan, en allemand comme en français :
\begin{enumerate}
\item \textbf{Présentation} (\all{die Vorstellung}) : indiquer le type de texte, le thème et la source. Commencer par : \all{Der Text handelt von...} (Le texte traite de...), \all{Der Text informiert über...} (Le texte informe sur...).
\item \textbf{Résumé} (\all{die Zusammenfassung}) : donner les idées principales, une phrase par paragraphe, avec les connecteurs chronologiques : \all{Zuerst..., dann..., danach..., schließlich...} (D'abord..., puis..., ensuite..., enfin...).
\item \textbf{Analyse} (\all{die Analyse}) : relever le vocabulaire du thème, les structures grammaticales importantes (modaux, datif, accusatif) et comparer les deux pays.
\item \textbf{Opinion personnelle} (\all{die eigene Meinung}) : terminer par son avis argumenté : \all{Meiner Meinung nach...} (À mon avis...), \all{Ich finde..., dass...} (Je trouve que...), \all{Ich finde es interessant, dass...} (Je trouve intéressant que...).
\end{enumerate}
\end{methode}

\begin{popfigure}
\begin{tikzpicture}[
  node distance=0.45cm,
  etape/.style={rectangle, rounded corners, draw=popblue, fill=popblueL, text width=4.1cm, align=center, minimum height=1.5cm, font=\small},
  fleche/.style={-{Stealth[length=2.5mm]}, thick, poporange}
]
\node[etape, align=center] (e1){\textbf{1. Présentation}\\ \all{Der Text handelt von...}\\ \all{Der Text informiert über...}};
\node[etape, below=of e1, align=center] (e2){\textbf{2. Résumé}\\ \all{Zuerst..., dann...,}\\ \all{danach..., schließlich...}};
\node[etape, below=of e2, align=center] (e3){\textbf{3. Analyse}\\ Vocabulaire du thème,\\ modaux, datif/accusatif,\\ comparaison};
\node[etape, below=of e3, align=center] (e4){\textbf{4. Meine Meinung}\\ \all{Meiner Meinung nach...}\\ \all{Ich finde..., dass...}};
\draw[fleche] (e1) -- (e2);
\draw[fleche] (e2) -- (e3);
\draw[fleche] (e3) -- (e4);
\end{tikzpicture}
\legende{Organigramme du commentaire de texte : quatre étapes avec les expressions utiles.}
\end{popfigure}

\begin{aretenir}[Les phrases-cadres à mémoriser]
\begin{itemize}
\item Présentation : \all{Der Text handelt von den Hochzeiten in Deutschland und in Kamerun.} « Le texte traite des mariages en Allemagne et au Cameroun. »
\item Résumé : \all{Zuerst erklärt er die standesamtliche Trauung.} « D'abord, il explique le mariage civil. »
\item Analyse : \all{Der Text vergleicht die Traditionen der beiden Länder.} « Le texte compare les traditions des deux pays. »
\item Opinion : \all{Meiner Meinung nach ist die Mitgift eine schöne Tradition.} « À mon avis, la dot est une belle tradition. »
\end{itemize}
\end{aretenir}

\courssub{Commentaire modèle rédigé}

Lisons ensemble ce commentaire modèle du texte support. Il respecte les quatre étapes et réutilise le vocabulaire et les structures étudiés dans la leçon.

\begin{textebox}[Kommentar zum Text „Heiraten in Deutschland und in Kamerun“ — commentaire modèle]
Der Text informiert über die Hochzeit in Deutschland und in Kamerun. Er zeigt die Unterschiede und die Gemeinsamkeiten der beiden Länder.

Zuerst erklärt er die standesamtliche Trauung in Deutschland: Das Paar muss zum Standesamt gehen, und ein Beamter gibt dem Brautpaar die Urkunde. Dann beschreibt er die drei Hochzeiten in Kamerun: die traditionelle Hochzeit mit der Mitgift, die standesamtliche Hochzeit und die kirchliche Hochzeit. Danach erklärt er, dass die Familie in Kamerun eine große Rolle spielt. Schließlich sagt er, dass viele Paare in beiden Ländern auch eine Party mit Musik und Essen organisieren.

Der Text benutzt viele Modalverben: „muss“, „kann“, „darf“. Das Vokabular des Themas ist reich: die Verlobung, der Bräutigam, die Braut, die Mitgift, die Ehefrau, der Ehemann.

Ich finde es interessant, dass man in Kamerun dreimal feiern kann. Meiner Meinung nach ist die Mitgift eine schöne Tradition, aber sie darf nicht zu teuer sein. Ich finde auch, dass die Familie bei einer Hochzeit sehr wichtig ist.
\end{textebox}

\textbf{Glossaire :} der Unterschied, -e (la différence) ; die Gemeinsamkeit, -en (le point commun) ; das Standesamt, -ämter (l'état civil) ; der Beamte, -n (le fonctionnaire) ; die Urkunde, -n (l'acte, le certificat) ; die Rolle, -n (le rôle) ; reich (riche).

\textbf{Analyse collective du modèle.} Repérons les quatre étapes : le paragraphe 1 présente le texte (\all{Der Text informiert über...}) ; le paragraphe 2 résume avec \all{zuerst, dann, danach, schließlich} ; le paragraphe 3 analyse (modaux, vocabulaire) ; le paragraphe 4 donne l'opinion (\all{Ich finde es interessant, dass... / Meiner Meinung nach...}). Remarquons aussi les cas : \all{gibt dem Brautpaar die Urkunde} — \all{dem Brautpaar} est au datif (personne qui reçoit), \all{die Urkunde} à l'accusatif (chose donnée).

\courssub{Le plan de résumé : une phrase par paragraphe}

\begin{propriete}[Résumer avec les connecteurs chronologiques]
Pour résumer un texte, on rédige \textbf{une phrase par paragraphe}, reliées par des connecteurs placés en position 1 (le verbe conjugué reste en position 2) :
\begin{center}
\begin{tabularx}{\linewidth}{|l|X|X|}\hline
\rowcolor{popblueL} \textbf{Connecteur} & \textbf{Français} & \textbf{Exemple} \\ \hline
\all{zuerst} & d'abord & \all{Zuerst erklärt er die Trauung.} \\ \hline
\all{dann} & puis & \all{Dann beschreibt er die Mitgift.} \\ \hline
\all{danach} & ensuite & \all{Danach spricht er von der Familie.} \\ \hline
\all{schließlich} & enfin & \all{Schließlich nennt er die Party.} \\ \hline
\end{tabularx}
\end{center}
\end{propriete}

\begin{attention}[La place du verbe après le connecteur]
En français, on dit « D'abord, il explique... ». En allemand, le connecteur occupe la position 1, donc le sujet passe après le verbe : \all{Zuerst \textbf{erklärt er}...} et jamais \emph{Zuerst er erklärt...}. C'est la règle d'or de l'inversion : le verbe conjugué est toujours en deuxième position.
\end{attention}

\courssub{Production écrite guidée : un paragraphe comparatif}

\begin{methode}[Rédiger « Die Hochzeit in Kamerun und in Deutschland » (6 à 8 phrases)]
\begin{enumerate}
\item Commencer par une phrase d'introduction : \all{Die Hochzeit ist in Kamerun und in Deutschland ein großes Fest.}
\item Comparer un point précis avec \all{aber} ou \all{in Deutschland nicht} : utiliser au moins \textbf{trois modaux} (\all{muss, kann, darf...}).
\item Inclure au moins \textbf{deux compléments au datif} : \all{der Bräutigam gibt der Familie der Braut eine Mitgift} ; \all{das Paar dankt den Gästen}.
\item Terminer par son opinion : \all{Ich finde..., dass...} ou \all{Meiner Meinung nach...}.
\item Relire : verbe en position 2, infinitif à la fin avec les modaux, majuscules aux noms.
\end{enumerate}
\end{methode}

\begin{exemplebox}[Phrases-cadres traduites pour la production]
\all{In Kamerun muss der Bräutigam eine Mitgift zahlen, in Deutschland nicht.} « Au Cameroun, le marié doit payer une dot, pas en Allemagne. »

\all{Der Bräutigam gibt der Familie der Braut eine Mitgift.} « Le marié donne une dot à la famille de la mariée. » (datif : \all{der Familie der Braut})

\all{In Deutschland kann das Paar nur standesamtlich heiraten.} « En Allemagne, le couple peut se marier seulement à l'état civil. »

\all{Ich finde die traditionelle Hochzeit schön, weil die ganze Familie feiert.} « Je trouve le mariage traditionnel beau parce que toute la famille fait la fête. » (avec \all{weil}, le verbe \all{feiert} va à la fin)
\end{exemplebox}

\begin{exoresolu}[Production guidée : rédiger le paragraphe comparatif]
Énoncé : rédigez un paragraphe de 6 à 8 phrases intitulé \all{Die Hochzeit in Kamerun und in Deutschland}, avec au moins 3 verbes de modalité et 2 compléments au datif.
\tcblower
\textbf{Solution détaillée (corrigé possible) :}

\all{Die Hochzeit ist in Kamerun und in Deutschland ein wichtiges Fest. In Kamerun muss der Bräutigam der Familie der Braut eine Mitgift geben, in Deutschland nicht. Dort muss das Paar nur zum Standesamt gehen. In Kamerun kann man dreimal feiern: traditionell, standesamtlich und kirchlich. Nach der Trauung dankt das Paar den Gästen für die Geschenke. In beiden Ländern wollen alle Gäste essen, trinken und tanzen. Ich finde die traditionelle Hochzeit schön, weil die ganze Familie feiert. Meiner Meinung nach darf die Mitgift aber nicht zu teuer sein.}

Vérification : 8 phrases ; 5 modaux (\all{muss, muss, kann, wollen, darf}) avec infinitif à la fin ; 2 datifs (\all{der Familie der Braut}, \all{den Gästen}) ; opinion finale avec \all{weil} + verbe en fin de phrase.
\end{exoresolu}

\begin{attention}[La subordonnée avec \all{dass} : le verbe à la fin]
Dans \all{Ich finde, dass...} (Je trouve que...), la subordonnée introduite par \all{dass} renvoie le verbe conjugué \textbf{à la fin} : \all{Ich finde, dass die Mitgift eine schöne Tradition \textbf{ist}.} Ne dites jamais \emph{Ich finde, dass die Mitgift ist eine schöne Tradition}. Même règle avec \all{weil} : \all{..., weil die ganze Familie \textbf{feiert}.} C'est un point capital pour le Bac.
\end{attention}

\courssub{Expression orale guidée : dialogue en binôme}

\begin{experience}[Dialogue : un élève camerounais et un élève allemand]
Par binôme, jouez le dialogue suivant : un élève camerounais (A) et un élève allemand (B) s'expliquent leurs traditions de mariage. Chacun pose au moins deux questions et utilise au moins deux modaux.

\textbf{A :} \all{Wie heiratet man in Deutschland?}

\textbf{B :} \all{Wir müssen zuerst zum Standesamt gehen. Dort bekommt das Paar die Urkunde. Und in Kamerun?}

\textbf{A :} \all{Bei uns gibt es drei Hochzeiten! Der Bräutigam muss der Familie der Braut eine Mitgift geben.}

\textbf{B :} \all{Das ist interessant! Kann man auch in der Kirche heiraten?}

\textbf{A :} \all{Ja, natürlich. Nach der Trauung feiern wir mit der ganzen Familie. Darf man in Deutschland eine große Party machen?}

\textbf{B :} \all{Ja, klar! Wir wollen auch tanzen und essen. Ich finde eure Traditionen sehr schön.}

Puis inversez les rôles et ajoutez deux questions personnelles : \all{Was findest du schön? Was darf man nicht vergessen?}
\end{experience}

\courssub{Grille d'auto-évaluation}

Après votre production écrite ou orale, évaluez-vous avec cette grille. Cochez chaque critère atteint.

\begin{center}
\begin{tabularx}{\linewidth}{|X|l|l|}\hline
\rowcolor{popblueL} \textbf{Critère} & \textbf{Oui} & \textbf{Non} \\ \hline
J'emploie au moins 5 mots du vocabulaire du thème (\all{die Hochzeit, die Mitgift, die Verlobung...}) & $\square$ & $\square$ \\ \hline
Mes cas sont corrects : accusatif pour la chose, datif pour la personne (\all{der Familie, den Gästen}) & $\square$ & $\square$ \\ \hline
Mes modaux sont bien placés : modal en position 2, infinitif à la fin & $\square$ & $\square$ \\ \hline
Après \all{dass} et \all{weil}, mon verbe conjugué est à la fin & $\square$ & $\square$ \\ \hline
Mes phrases sont complètes : sujet, verbe, majuscules aux noms, point final & $\square$ & $\square$ \\ \hline
Je donne mon opinion avec \all{Meiner Meinung nach...} ou \all{Ich finde..., dass...} & $\square$ & $\square$ \\ \hline
\end{tabularx}
\end{center}

\begin{savaistu}[La stratégie gagnante au Bac]
Au baccalauréat, l'expression écrite est notée sur trois critères : le respect de la consigne (nombre de phrases, éléments imposés), la richesse du vocabulaire et la correction grammaticale. Réutiliser les structures et le vocabulaire du texte support — les modaux, les datifs, les connecteurs \all{zuerst, dann, schließlich} — est une stratégie gagnante : vous montrez que vous avez compris le texte et que vous savez vous en servir. Entraînez-vous à rédiger un paragraphe de ce type en 15 minutes, comme le jour de l'examen.
\end{savaistu}

\begin{exercice}[À vous de jouer]
\begin{enumerate}
\item Résumez le texte support en quatre phrases avec \all{zuerst, dann, danach, schließlich}.
\item Rédigez votre propre paragraphe \all{Die Hochzeit in Kamerun und in Deutschland} (6 à 8 phrases, 3 modaux, 2 datifs).
\item Traduisez : « Je trouve que le mariage traditionnel est une belle fête, parce que toute la famille danse. »
\end{enumerate}
\end{exercice}

\begin{corrige}[Exercice : À vous de jouer]
\begin{enumerate}
\item \all{Zuerst erklärt der Text die standesamtliche Trauung in Deutschland. Dann beschreibt er die drei Hochzeiten in Kamerun. Danach spricht er von der Mitgift und der Familie. Schließlich sagt er, dass alle Gäste feiern wollen.}
\item Corrigé individuel : vérifier avec la grille d'auto-évaluation (5 mots du thème, 3 modaux avec infinitif final, 2 datifs, opinion finale).
\item \all{Ich finde, dass die traditionelle Hochzeit ein schönes Fest ist, weil die ganze Familie tanzt.} (verbes \all{ist} et \all{tanzt} à la fin des subordonnées).
\end{enumerate}
\end{corrige}

\newpage

\coursec{Activité d'intégration}

\courssub{Objectifs de l'activité}

À la fin de cette activité d'intégration, tu seras capable de :
\begin{itemize}
\item réinvestir le vocabulaire du mariage (\all{die Hochzeit}, \all{die Verlobung}, \all{die Mitgift}, \all{der Ehemann}, \all{die Ehefrau}, \all{standesamtlich}, \all{kirchlich}, \all{der Bräutigam}, \all{die Braut}) dans une communication orale et écrite authentique ;
\item formuler des questions et des affirmations avec les verbes de modalité au présent (\all{müssen}, \all{wollen}, \all{können}, \all{dürfen}, \all{sollen}) en respectant la place du verbe (position 2) et de l'infinitif (fin de phrase) ;
\item employer correctement l'accusatif (complément d'objet direct) et le datif (complément d'objet indirect / bénéficiaire) avec les articles, les pronoms et les adjectifs ;
\item rédiger un court article journalistique de 8 à 10 lignes pour le journal de l'école en respectant des contraintes linguistiques précises ;
\item travailler en binôme : écouter, réagir, corriger et améliorer une production commune.
\end{itemize}

\courssub{Situation}

Le journal de ton lycée, \emph{La Plume du Lycée}, prépare un numéro spécial sur les traditions familiales. Un journaliste allemand, Thomas Weber, est en reportage à Yaoundé pour un échange scolaire entre un \all{Gymnasium} de Munich et ton établissement. Il veut comprendre comment on se marie au Cameroun et quelles différences existent avec l'Allemagne. Il rencontre Junior, un jeune Camerounais de 24 ans, fiancé à Aïcha, et lui accorde une interview. Ensuite, chaque binôme de la classe rédige un court article pour le journal de l'école.

\begin{textebox}[Extrait du carnet de notes du journaliste]
Thomas Weber prépare son interview. Voici ses notes :\par
\medskip
\all{Notizen für das Interview (Yaoundé):}\par
\all{1. Verlobung: Wie lange? Traditionelle Zeremonie?}\par
\all{2. Mitgift: Muss der Mann eine Mitgift zahlen? Was genau?}\par
\all{3. Hochzeit: standesamtlich? kirchlich? traditionell?}\par
\all{4. Gäste: Wie viele? Wer hilft der Familie?}\par
\all{5. Unterschiede zu Deutschland?}\par
\medskip
Glossaire : \emph{die Notiz, -en} : la note ; \emph{die Zeremonie, -n} : la cérémonie ; \emph{der Gast, die Gäste} : l'invité ; \emph{der Unterschied, -e} : la différence ; \emph{traditionell} : traditionnel(le).
\end{textebox}

\courssub{Consignes}

\textbf{Étape 1 — Préparer les questions (journaliste).} En binôme, l'élève A joue le journaliste allemand. À l'aide des notes de Thomas Weber, il rédige \textbf{5 questions complètes} contenant chacune un verbe de modalité au présent. Modèles : \all{Müssen Sie eine Mitgift zahlen?} ; \all{Wollen Sie auch kirchlich heiraten?} ; \all{Können die Gäste der Familie helfen?} Attention à la place du modal en position 2 (ou position 1 pour question fermée) et de l'infinitif à la fin de la phrase.

\textbf{Étape 2 — Préparer les réponses (fiancé).} L'élève B joue Junior, le jeune fiancé camerounais. Il prépare des réponses en \textbf{phrases complètes}, avec au moins \textbf{2 compléments à l'accusatif} et \textbf{2 compléments au datif} repérables (articles ou adjectifs soulignés). Exemple : \all{Ich muss meiner Familie und der Familie meiner Verlobten viele Geschenke geben.}

\textbf{Étape 3 — Jouer l'interview.} Jouez la scène à voix haute devant la classe ou en ateliers. Le journaliste pose ses questions, Junior répond, puis le journaliste peut poser une question supplémentaire spontanée. Les autres binômes écoutent et notent les verbes de modalité entendus.

\textbf{Étape 4 — Rédiger l'article.} Ensemble, rédigez un court article de \textbf{8 à 10 lignes} intitulé \all{„Eine Hochzeit in Kamerun“} pour le journal de l'école. Contraintes obligatoires :
\begin{enumerate}
\item au moins \textbf{4 verbes de modalité} différents au présent ;
\item au moins \textbf{2 compléments au datif} (article ou adjectif au datif souligné) ;
\item au moins \textbf{5 mots du lexique de la leçon} (par exemple \all{die Verlobung}, \all{die Mitgift}, \all{die Hochzeit}, \all{standesamtlich}, \all{die Ehefrau}) ;
\item une phrase de conclusion qui compare avec l'Allemagne.
\end{enumerate}

\textbf{Étape 5 — Relecture et affichage.} Échangez votre article avec un autre binôme : vérifiez la place des verbes, les cas (accusatif/datif), les terminaisons des adjectifs et les majuscules des noms. Corrigez, puis remettez la version finale. Les meilleurs articles sont lus en classe et affichés au mur de la salle d'allemand.

\courssub{Ressources à mobiliser}

\begin{aretenir}[Boîte à outils de l'interview : Lexique et structures]
\begin{itemize}
\item \textbf{Lexique de base (nom, article, pluriel) :}
\begin{center}
\begin{tabular}{|l|l|l|}
\hline
\rowcolor{popblueL} \textbf{Singulier} & \textbf{Pluriel} & \textbf{Français} \\ \hline
\all{die Hochzeit} & \all{die Hochzeiten} & le mariage, la noce \\ \hline
\all{die Verlobung} & \all{die Verlobungen} & les fiançailles \\ \hline
\all{die Mitgift} & \all{die Mitgiften} & la dot \\ \hline
\all{der Ehemann} & \all{die Ehemänner} & l'époux, le mari \\ \hline
\all{die Ehefrau} & \all{die Ehefrauen} & l'épouse, la femme \\ \hline
\all{der Bräutigam} & \all{die Bräutigame} & le fiancé, le marié \\ \hline
\all{die Braut} & \all{die Bräute} & la fiancée, la mariée \\ \hline
\all{die Feier} & \all{die Feiern} & la fête, la célébration \\ \hline
\all{der Gast} & \all{die Gäste} & l'invité \\ \hline
\all{das Standesamt} & \all{die Standesämter} & l'état civil, la mairie \\ \hline
\end{tabular}
\end{center}
\item \textbf{Adverbes / Prépositions utiles :} \all{standesamtlich heiraten} (se marier à la mairie), \all{kirchlich heiraten} (se marier à l'église), \all{traditionell} (traditionnel), \all{zur Hochzeit} (pour le mariage / au mariage), \all{bei der Feier} (lors de la fête).
\end{itemize}
\end{aretenir}

\begin{propriete}[Conjugaison des verbes de modalité au présent (Indikativ)]
Les verbes de modalité (\all{müssen, wollen, können, dürfen, sollen}) sont irréguliers aux 3 personnes du singulier (changement de voyelle pour \all{müssen, wollen, können, dürfen}, pas pour \all{sollen}). L'infinitif se place \textbf{toujours en fin de phrase}.
\begin{center}
\begin{tabular}{|l|c|c|c|c|c|}
\hline
\rowcolor{popblueL} \textbf{Personne} & \textbf{müssen} & \textbf{wollen} & \textbf{können} & \textbf{dürfen} & \textbf{sollen} \\ \hline
\all{ich} & \all{muss} & \all{will} & \all{kann} & \all{darf} & \all{soll} \\ \hline
\all{du} & \all{musst} & \all{willst} & \all{kannst} & \all{darfst} & \all{sollst} \\ \hline
\all{er / sie / es} & \all{muss} & \all{will} & \all{kann} & \all{darf} & \all{soll} \\ \hline
\all{wir} & \all{müssen} & \all{wollen} & \all{können} & \all{dürfen} & \all{sollen} \\ \hline
\all{ihr} & \all{müsst} & \all{wollt} & \all{könnt} & \all{dürft} & \all{sollt} \\ \hline
\all{sie / Sie} & \all{müssen} & \all{wollen} & \all{können} & \all{dürfen} & \all{sollen} \\ \hline
\end{tabular}
\end{center}
\emph{Exemples :} \all{Ich muss die Mitgift zahlen.} (Je dois payer la dot.) — \all{Wir wollen standesamtlich heiraten.} (Nous voulons nous marier civilement.) — \all{Kannst du mir helfen?} (Peux-tu m'aider ?)
\end{propriete}

\begin{propriete}[L'accusatif et le datif : articles, pronoms et adjectifs]
\textbf{Rappel :} L'accusatif marque le \textbf{complément d'objet direct (COD)} (qui/quoi ?). Le datif marque le \textbf{complément d'objet indirect / bénéficiaire (COI)} (à qui ?) et suit certaines prépositions (\all{mit, zu, bei, von, nach, seit, aus, gegenüber}).
\begin{center}
\begin{tabular}{|l|c|c|c|}
\hline
\rowcolor{popblueL} & \textbf{Nominatif} & \textbf{Accusatif} & \textbf{Datif} \\ \hline
\textbf{Masculin} & \all{der / ein} & \all{den / einen} & \all{dem / einem} \\ \hline
\textbf{Féminin} & \all{die / eine} & \all{die / eine} & \all{der / einer} \\ \hline
\textbf{Neutre} & \all{das / ein} & \all{das / ein} & \all{dem / einem} \\ \hline
\textbf{Pluriel} & \all{die / —} & \all{die / —} & \all{den / —} \\ \hline
\end{tabular}
\end{center}
\textbf{Déclinaison des adjectifs (après article défini / indéfini) :}
\begin{center}
\begin{tabular}{|l|c|c|}
\hline
\rowcolor{popblueL} \textbf{Cas} & \textbf{Article défini (\all{der/die/das})} & \textbf{Article indéfini (\all{ein/eine})} \\ \hline
\all{Acc. masc.} & \all{den großen Ring} & \all{einen großen Ring} \\ \hline
\all{Dat. masc./neut.} & \all{dem großen Ring / dem kleinen Kind} & \all{einem großen Ring / einem kleinen Kind} \\ \hline
\all{Dat. fém.} & \all{der großen Familie} & \all{einer großen Familie} \\ \hline
\all{Dat. pluriel} & \all{den großen Gästen} & \all{— großen Gästen} \\ \hline
\end{tabular}
\end{center}
\emph{Verbes clés au datif :} \all{helfen} (aider), \all{geben} (donner qqch \textbf{à qqn} : \all{geben + Dat + Acc}), \all{gehören}, \all{gefallen}, \all{danken}, \all{folgen}, \all{glauben}.
\end{propriete}

\begin{attention}[Pièges à éviter pour les francophones]
\begin{itemize}
\item \textbf{Ordre des mots (Modal + Infinitif) :} Ne dis pas \all{*Ich muss zahlen die Mitgift}. L'infinitif va \textbf{à la fin} : \all{Ich muss die Mitgift zahlen.} / Question : \all{Müssen Sie die Mitgift zahlen?}
\item \textbf{\all{helfen} + DATIF :} \all{Die Gäste helfen der Familie.} (Jamais \all{*die Familie} accusatif).
\item \textbf{Prépositions + DATIF :} Après \all{mit, zu, von, bei, nach, seit, aus} : toujours le datif. \all{zur Feier} = \all{zu der Feier}, \all{bei der Hochzeit}.
\item \textbf{MaJUSCULES :} Tous les noms prennent une majuscule : \all{die Hochzeit}, \all{das Geschenk}, \all{die Mitgift}, \all{der Bräutigam}.
\item \textbf{\all{kein} vs \all{nicht} :} \all{keine Mitgift} (accusatif, nie \all{*kein Mitgift}), \all{nicht standesamtlich}.
\item \textbf{Faux ami :} \all{das Geschenk} = le cadeau (pas \all{le poison} = \all{das Gift}). \all{die Mitgift} = la dot (pas \all{le don}).
\end{itemize}
\end{attention}

\courssub{Entraînement grammatical : Cas et Modaux}

\begin{exercice}[Mise en application immédiate]
Complétez les phrases avec la forme correcte (article + adjectif si nécessaire) au \textbf{cas accusatif} ou \textbf{datif}, et conjuguez le \textbf{verbe de modalité} entre parenthèses.
\begin{enumerate}
\item Junior \_\_\_\_\_\_\_\_\_\_\_\_ (müssen) \_\_\_\_\_\_\_\_\_\_\_\_ (eine große Mitgift) \_\_\_\_\_\_\_\_\_\_\_\_ (geben).
\item \_\_\_\_\_\_\_\_\_\_\_\_ (Können) \_\_\_\_\_\_\_\_\_\_\_\_ (die Gäste) \_\_\_\_\_\_\_\_\_\_\_\_ (der Familie) \_\_\_\_\_\_\_\_\_\_\_\_ (helfen)?
\item Wir \_\_\_\_\_\_\_\_\_\_\_\_ (wollen) \_\_\_\_\_\_\_\_\_\_\_\_ (zuerst) \_\_\_\_\_\_\_\_\_\_\_\_ (standesamtlich) \_\_\_\_\_\_\_\_\_\_\_\_ (heiraten).
\item \_\_\_\_\_\_\_\_\_\_\_\_ (Sollen) \_\_\_\_\_\_\_\_\_\_\_\_ (wir) \_\_\_\_\_\_\_\_\_\_\_\_ (viele Geschenke) \_\_\_\_\_\_\_\_\_\_\_\_ (mitbringen)?
\item Die Braut \_\_\_\_\_\_\_\_\_\_\_\_ (dürfen) \_\_\_\_\_\_\_\_\_\_\_\_ (dem Bräutigam) \_\_\_\_\_\_\_\_\_\_\_\_ (einen Ring) \_\_\_\_\_\_\_\_\_\_\_\_ (schenken).
\item In Deutschland \_\_\_\_\_\_\_\_\_\_\_\_ (müssen) \_\_\_\_\_\_\_\_\_\_\_\_ (der Mann) \_\_\_\_\_\_\_\_\_\_\_\_ (keine Mitgift) \_\_\_\_\_\_\_\_\_\_\_\_ (zahlen).
\end{enumerate}
\end{exercice}

\begin{corrige}[Exercice Entraînement]
\begin{enumerate}
\item Junior \all{muss} \all{eine große Mitgift} \all{geben}. (Accusatif : \all{eine große Mitgift} — féminin, indéfini).
\item \all{Können} \all{die Gäste} \all{der Familie} \all{helfen}? (Sujet \all{die Gäste} nominatif pluriel ; \all{der Familie} datif singulier féminin après \all{helfen}).
\item Wir \all{wollen} \all{zuerst} \all{standesamtlich} \all{heiraten}. (Position 2 : \all{wollen} ; infinitif \all{heiraten} à la fin).
\item \all{Sollen} \all{wir} \all{viele Geschenke} \all{mitbringen}? (Question fermée : modal position 1 ; \all{viele Geschenke} accusatif pluriel).
\item Die Braut \all{darf} \all{dem Bräutigam} \all{einen Ring} \all{schenken}. (\all{dem Bräutigam} datif masc. ; \all{einen Ring} accusatif masc. indéfini).
\item In Deutschland \all{muss} \all{der Mann} \all{keine Mitgift} \all{zahlen}. (\all{keine Mitgift} accusatif féminin singulier ; \all{der Mann} sujet nominatif).
\end{enumerate}
\end{corrige}

\courssub{Proposition de production et corrigé commenté}

\begin{exoresolu}[Interview et article modèles]
Énoncé : produisez l'interview (5 questions-réponses) puis l'article « Eine Hochzeit in Kamerun » (8 à 10 lignes) en respectant les contraintes : 4 modaux, 2 datifs, 5 mots du lexique.
\tcblower
\textbf{1. Interview modèle}\par
\medskip
\all{Thomas: Guten Tag, Junior! Müssen Sie bei einer Hochzeit in Kamerun eine Mitgift zahlen?}\par
\all{Junior: Ja, ich muss der Familie meiner Verlobten eine Mitgift geben, zum Beispiel Geld, Stoffe und Getränke.}\par
\all{Thomas: Wollen Sie standesamtlich oder kirchlich heiraten?}\par
\all{Junior: Wir wollen beides machen: zuerst heiraten wir standesamtlich, dann kirchlich in der Kirche.}\par
\all{Thomas: Können die Gäste der Familie bei der Vorbereitung helfen?}\par
\all{Junior: Natürlich! Die Gäste helfen der Familie und bringen oft Geschenke mit.}\par
\all{Thomas: Dürfen die Kinder auch zur Feier kommen?}\par
\all{Junior: Ja, alle Kinder dürfen kommen. Eine Hochzeit ist ein Fest für das ganze Dorf.}\par
\all{Thomas: Sollen die Deutschen auch eine Mitgift zahlen?}\par
\all{Junior: Nein, in Deutschland muss der Bräutigam keine Mitgift zahlen. Das ist ein großer Unterschied.}\par
\medskip
\textbf{2. Article modèle}\par
\medskip
\all{Eine Hochzeit in Kamerun}\par
\all{In Kamerun ist eine Hochzeit ein großes Fest für die ganze Familie. Vor der Hochzeit gibt es die Verlobung: Der Bräutigam muss der Familie der Braut eine Mitgift geben. Er kann Geld, Stoffe und Getränke bringen. Viele Paare wollen zweimal heiraten: zuerst standesamtlich im Standesamt, dann kirchlich in der Kirche. Bei der Feier dürfen alle Gäste essen, trinken und tanzen. Die Freunde sollen der Familie helfen, denn eine Hochzeit ist viel Arbeit. Die Ehefrau und der Ehemann bekommen viele Geschenke. In Deutschland muss der Mann keine Mitgift zahlen — das ist anders als in Kamerun!}\par
\medskip
\textbf{Commentaire des choix de langue :}
\begin{itemize}
\item \textbf{Modaux utilisés (5) :} \all{muss} (2×), \all{kann}, \all{wollen}, \all{dürfen}, \all{sollen} — chacun conjugué en position 2 (ou 1 en question), avec l'infinitif en fin de phrase (\all{muss ... geben}, \all{kann ... bringen}, \all{wollen ... heiraten}, \all{dürfen ... kommen}, \all{sollen ... helfen}).
\item \textbf{Datifs (soulignés dans l'article) :} \all{\textbf{der Familie der Braut}} (datif féminin singulier après \all{geben}), \all{\textbf{der Familie helfen}} (verbe \all{helfen} + datif), \all{\textbf{zur Feier}} = \all{zu der Feier} (préposition \all{zu} + datif), \all{\textbf{dem ganzen Dorf}} (datif neutre après \all{für} ? Non, \all{für} + accusatif : \all{für das ganze Dorf}. Correction : \all{für das ganze Dorf} est accusatif. Les datifs clairs sont \all{der Familie der Braut}, \all{der Familie}, \all{zur Feier}).
\item \textbf{Accusatifs :} \all{eine Mitgift geben}, \all{Geld, Stoffe und Getränke bringen}, \all{viele Geschenke bekommen}, \all{keine Mitgift zahlen}.
\item \textbf{Déclinaison adjective :} \all{ein \textbf{großes} Fest} (neutre accusatif indéfini), \all{die \textbf{ganze} Familie} (féminin accusatif défini), \all{\textbf{keine} Mitgift} (féminin accusatif négatif).
\item \textbf{Lexique de la leçon (9 mots) :} \all{die Hochzeit} (3×), \all{die Verlobung}, \all{die Mitgift} (3×), \all{der Bräutigam}, \all{die Braut}, \all{standesamtlich}, \all{kirchlich}, \all{die Ehefrau}, \all{der Ehemann}, \all{die Feier}, \all{die Gäste}, \all{die Freunde}, \all{die Familie}.
\item \textbf{Comparaison finale :} la dernière phrase oppose le Cameroun et l'Allemagne avec la négation \all{keine Mitgift} (accusatif), ce qui répond à la consigne de conclusion comparative.
\end{itemize}
\end{exoresolu}

\courssub{Bilan de l'activité}

\begin{aretenir}[Ce que j'ai appris à faire]
\begin{itemize}
\item Je sais poser des questions avec un verbe de modalité : modal en position 2 (ou 1 pour question fermée), sujet, infinitif à la fin : \all{Müssen Sie eine Mitgift zahlen?}
\item Je sais distinguer l'accusatif (COD : \all{den Ring kaufen}, \all{eine Mitgift geben}) et le datif (COI / bénéficiaire : \all{der Familie helfen}, \all{dem Bräutigam schenken}, \all{zur Feier kommen}).
\item Je connais la déclinaison des articles et adjectifs aux deux cas (masc. \all{den/einen} vs \all{dem/einem} ; fém. \all{die/eine} vs \all{der/einer} ; pluriel \all{die/—} vs \all{den/—}).
\item Je connais le vocabulaire essentiel du mariage (\all{Hochzeit, Verlobung, Mitgift, Ehemann, Ehefrau, Bräutigam, Braut, standesamtlich, kirchlich}) et je peux comparer les traditions du Cameroun et de l'Allemagne.
\item Je sais rédiger un court article journalistique structuré : titre, présentation, détails culturels, conclusion comparative.
\end{itemize}
\end{aretenir}

\begin{experience}[Pour aller plus loin : le mur des articles]
Après la lecture des meilleurs articles, la classe vote pour les trois productions les plus riches (lexique varié, cas corrects, modaux diversifiés, orthographe soignée). Les articles gagnants sont affichés au « mur des Hochzeiten » de la salle d'allemand. 
\textbf{Variante orale :} Deux binômes rejouent l'interview en changeant les rôles — le journaliste devient fiancé(e) — et en ajoutant une question sur les mariages en Allemagne : \all{Wie heiratet man in Deutschland?} (Réponse attendue : \all{Man heiratet meistens nur standesamtlich, manchmal auch kirchlich. Es gibt keine Mitgift.})
\end{experience}

\newpage

\coursec{Méthodes et exercices résolus}

\courssub{Méthodes et exercices résolus}

\begin{methode}[Comprendre un texte sur les types de mariages]
Pour réussir les questions de compréhension au Bac, suis toujours la même démarche :
\begin{enumerate}
\item \textbf{Lis le texte une première fois} sans dictionnaire : repère le thème général (ici : \cle{die Hochzeit}, les formes de mariage en Allemagne et au Cameroun).
\item \textbf{Lis les questions avant la deuxième lecture} : elles te guident vers les passages utiles.
\item \textbf{Repère les mots-clés} de chaque question et souligne dans le texte la phrase qui y répond.
\item \textbf{Réponds par une phrase complète} en allemand, avec tes propres mots : ne recopie jamais une longue phrase du texte.
\item \textbf{Vérifie la grammaire} de ta réponse : verbe conjugué en deuxième position, cas corrects après les verbes et les prépositions.
\end{enumerate}
\cle{Réflexe} : pour une question \all{Warum...?}, commence ta réponse par \all{Weil...} (verbe à la fin) ou \all{..., weil...}. Pour une question \all{Was...?}, reprends la structure du texte en la reformulant.
\end{methode}

\begin{exoresolu}[Questions de compréhension sur le texte]
Texte (extrait) : \all{In Deutschland heiraten viele Paare zuerst standesamtlich. Die standesamtliche Hochzeit findet im Rathaus statt, und nur sie ist vor dem Gesetz gültig. Danach können die Paare auch kirchlich heiraten. In Kamerun gibt es oft drei Hochzeiten: die traditionelle Hochzeit mit der Mitgift, die standesamtliche Hochzeit und die kirchliche Hochzeit. Die Familien spielen bei der traditionellen Hochzeit eine große Rolle.}

Questions : 1. \all{Wo findet die standesamtliche Hochzeit in Deutschland statt?} 2. \all{Warum ist die standesamtliche Hochzeit wichtig?} 3. \all{Wie viele Hochzeiten gibt es oft in Kamerun?}
\tcblower
\textbf{Raisonnement} : la question 1 porte sur le lieu (\all{Wo}) : le texte dit \all{im Rathaus}. La question 2 (\all{Warum}) exige une cause : \all{nur sie ist vor dem Gesetz gültig}. La question 3 demande un nombre : \all{drei Hochzeiten}.

\textbf{Réponses rédigées} :
\begin{enumerate}
\item \all{Die standesamtliche Hochzeit findet im Rathaus statt.} « Le mariage civil a lieu à la mairie. » (\all{im} = \all{in dem}, datif après \all{in} pour le lieu.)
\item \all{Sie ist wichtig, weil nur sie vor dem Gesetz gültig ist.} « Il est important parce que c'est le seul qui soit valable devant la loi. » (Subordonnée avec \all{weil} : verbe conjugué \all{ist} à la fin.)
\item \all{In Kamerun gibt es oft drei Hochzeiten: die traditionelle, die standesamtliche und die kirchliche Hochzeit.} « Au Cameroun, il y a souvent trois mariages. » (\all{es gibt} + accusatif : \all{drei Hochzeiten}.)
\end{enumerate}
\textbf{Conclusion} : chaque réponse est une phrase complète, reformulée, avec le verbe à la bonne place.
\end{exoresolu}

\begin{methode}[Employer le vocabulaire du mariage]
\begin{enumerate}
\item \textbf{Apprends chaque nom avec son article et son pluriel} : \cle{die Hochzeit, -en} ; \cle{der Ehemann, -männer} ; \cle{die Ehefrau, -en} ; \cle{die Verlobung, -en} ; \cle{die Mitgift, -en} ; \cle{das Brautpaar, -e}.
\item \textbf{Apprends les verbes avec leur construction} : \all{heiraten} + accusatif (\all{Sie heiratet ihren Freund.}) ; \all{sich verloben mit} + datif ; \all{sich verheiraten mit} + datif.
\item \textbf{Utilise les familles de mots} : \all{die Ehe} $\rightarrow$ \all{der Ehemann / die Ehefrau / ehelich} ; \all{heiraten} $\rightarrow$ \all{die Heirat / die Hochzeit}.
\item \textbf{Emploie les adjectifs utiles} : \cle{standesamtlich} (civil), \cle{kirchlich} (religieux), \cle{traditionell} (traditionnel).
\item \textbf{Dans une production}, varie : ne répète pas \all{Hochzeit} cinq fois, utilise \all{die Feier}, \all{das Fest}, \all{die Zeremonie}.
\end{enumerate}
\end{methode}

\begin{exoresolu}[Texte à trous : le vocabulaire de la Ehe]
Complète avec les mots suivants : \all{Verlobung, Mitgift, Ehemann, Hochzeit, standesamtlich}.

\all{Am Samstag ist die große (1) \dots\ von Marie und Paul. Vor einem Jahr war schon die (2) \dots\ mit den Familien. In Kamerun bringt der Bräutigam oft eine (3) \dots\ für die Familie der Braut. Am Montag gehen die beiden zum Rathaus, denn die (4) \dots\ Trauung ist vor dem Gesetz gültig. Nach der Feier ist Paul endlich der (5) \dots\ von Marie.}
\tcblower
\textbf{Raisonnement} : chaque trou se résout par le sens et par la grammaire (article, cas).
\begin{enumerate}
\item \all{Hochzeit} : \all{die große Hochzeit}, nominatif féminin après \all{ist}. « Le grand mariage. »
\item \all{Verlobung} : « il y a un an, ce fut déjà les fiançailles. »
\item \all{Mitgift} : accusatif féminin après \all{bringt} : \all{eine Mitgift}. « Le fiancé apporte souvent une dot. »
\item \all{standesamtlich} : adjectif devant \all{Trauung} (féminin, nominatif) avec article défini : \all{die standesamtliche Trauung}. « La cérémonie civile. »
\item \all{Ehemann} : nominatif masculin après \all{ist} : \all{der Ehemann von Marie}. « Le mari de Marie. »
\end{enumerate}
\textbf{Texte complété} : \all{Am Samstag ist die große Hochzeit von Marie und Paul. Vor einem Jahr war schon die Verlobung mit den Familien. In Kamerun bringt der Bräutigam oft eine Mitgift für die Familie der Braut. Am Montag gehen die beiden zum Rathaus, denn die standesamtliche Trauung ist vor dem Gesetz gültig. Nach der Feier ist Paul endlich der Ehemann von Marie.}
\end{exoresolu}

\begin{methode}[Choisir entre accusatif et datif]
\begin{enumerate}
\item \textbf{Pose la question} : \all{wen? / was?} (qui ? quoi ?) $\rightarrow$ \cle{accusatif} ; \all{wem?} (à qui ?) $\rightarrow$ \cle{datif}.
\item \textbf{Apprends les verbes à datif} par cœur : \all{danken, helfen, gratulieren, gefallen, gehören, antworten}. Exemple : \all{Wir gratulieren dem Brautpaar.} « Nous félicitons les mariés. »
\item \textbf{Prépositions + accusatif} : \all{durch, für, gegen, ohne, um}. \textbf{Prépositions + datif} : \all{aus, bei, mit, nach, seit, von, zu}.
\item \textbf{Prépositions à deux cas} (\all{in, an, auf, neben, über, unter, vor, hinter, zwischen}) : mouvement $\rightarrow$ accusatif (\all{Wohin?}) ; position $\rightarrow$ datif (\all{Wo?}). \all{Sie geht in die Kirche} (elle va à l'église, mouvement) / \all{Die Hochzeit ist in der Kirche} (position).
\item \textbf{Souviens-toi des formes} : der $\rightarrow$ den (acc.) / dem (dat.) ; die $\rightarrow$ die / der ; das $\rightarrow$ das / dem ; die (pl.) $\rightarrow$ die / den + \textbf{-n} au nom pluriel (\all{den Kindern}).
\end{enumerate}
\end{methode}

\begin{exoresolu}[Thème : phrases à traduire avec accusatif et datif]
Traduis en allemand : 1. « Nous invitons la famille à la mairie. » 2. « Le marié remercie ses parents. » 3. « Après la fête, les invités vont dans le restaurant. »
\tcblower
\textbf{Raisonnement et solutions} :
\begin{enumerate}
\item « Inviter quelqu'un » = \all{jemanden einladen} + accusatif : \all{die Familie} reste \all{die Familie}. « À la mairie » = mouvement avec \all{zu} + datif : \all{zum Standesamt}. Solution : \all{Wir laden die Familie zum Standesamt ein.} Attention : \all{einladen} est un verbe à particule séparable : \all{ein} va à la fin.
\item \all{danken} exige le datif : \all{seine Eltern} (pluriel) $\rightarrow$ \all{seinen Eltern} avec -n au nom. Solution : \all{Der Bräutigam dankt seinen Eltern.}
\item « Après la fête » = \all{nach} + datif : \all{nach der Feier}. « Vont dans le restaurant » = mouvement $\rightarrow$ accusatif : \all{in das Restaurant} = \all{ins Restaurant}. Solution : \all{Nach der Feier gehen die Gäste ins Restaurant.} Le verbe \all{gehen} reste en deuxième position malgré le complément en tête.
\end{enumerate}
\textbf{Conclusion} : toujours identifier la fonction (COD ou COI, mouvement ou position) avant de choisir le cas.
\end{exoresolu}

\begin{methode}[Conjuguer et employer les verbes de modalité au présent]
\begin{enumerate}
\item \textbf{Les six modaux} : \cle{dürfen} (permission), \cle{können} (capacité, possibilité), \cle{mögen/möchte} (aimer, vouloir poliment), \cle{müssen} (obligation), \cle{sollen} (devoir conseillé), \cle{wollen} (volonté).
\item \textbf{Conjugaison au présent} : la 1\iere{} et la 3\ieme{} personne du singulier sont identiques, sans -t : \all{ich kann, er kann} ; \all{ich muss, er muss}. Singulier souvent irrégulier : \all{ich darf, kann, mag, muss, soll, will}.
\item \textbf{Position dans la phrase} : le modal conjugué est en deuxième position, l'infinitif va \textbf{à la fin} : \all{Wir müssen die Familie einladen.} « Nous devons inviter la famille. » C'est la \cle{Satzklammer} (parenthèse de phrase).
\item \textbf{Dans une subordonnée} avec \all{weil/dass} : le modal conjugué va à la fin : \all{..., weil er die Mitgift bringen muss.}
\item \textbf{Nuances de sens} : \all{müssen} = obligation forte ; \all{sollen} = ce que les autres attendent de toi ; \all{dürfen} = autorisation ; \all{wollen} = projet personnel.
\item \textbf{Forme polie \all{möchte}} (Konjunktiv II, utilisé comme présent) : \all{ich möchte, du möchtest, er möchte, wir möchten, ihr möchtet, sie möchten}. Ex. : \all{Ich möchte heiraten.}
\end{enumerate}
\begin{center}
\begin{tabular}{|l|l|l|l|l|l|l|}
\hline
\rowcolor{popblueL} & dürfen & können & mögen & müssen & sollen & wollen \\
\hline
ich & darf & kann & mag & muss & soll & will \\
\hline
du & darfst & kannst & magst & musst & sollst & willst \\
\hline
er/sie/es & darf & kann & mag & muss & soll & will \\
\hline
wir & dürfen & können & mögen & müssen & sollen & wollen \\
\hline
ihr & dürft & könnt & mögt & müsst & sollt & wollt \\
\hline
sie/Sie & dürfen & können & mögen & müssen & sollen & wollen \\
\hline
\end{tabular}
\end{center}
\end{methode}

\begin{exoresolu}[Transformation et thème avec les modaux]
A. Mets la phrase à la forme demandée : \all{Die Gäste bringen Geschenke.} (avec \all{sollen}, à la 3\ieme{} personne du pluriel).

B. Traduis : 1. « En Allemagne, on peut se marier à la mairie. » 2. « Le fiancé doit parler avec la famille de la fiancée. » 3. « Nous voulons célébrer une grande fête. »
\tcblower
\textbf{A. Raisonnement} : \all{sollen} à la 3\ieme{} personne du pluriel = \all{sollen} (comme l'infinitif). Le verbe \all{bringen} passe à l'infinitif en fin de phrase. Solution : \all{Die Gäste sollen Geschenke bringen.} « Les invités sont censés apporter des cadeaux. »

\textbf{B. Raisonnement et solutions} :
\begin{enumerate}
\item « On » = \all{man} ; \all{können} avec \all{man} $\rightarrow$ \all{kann} ; « se marier » = \all{heiraten} (ici sans objet). Solution : \all{In Deutschland kann man standesamtlich heiraten.} ou \all{... im Standesamt heiraten.} Infinitif \all{heiraten} à la fin.
\item « Devoir » = \all{müssen} ; \all{der Bräutigam} (3\ieme{} sg.) $\rightarrow$ \all{muss} ; « parler avec » = \all{sprechen mit} + datif : \all{mit der Familie der Braut}. Solution : \all{Der Bräutigam muss mit der Familie der Braut sprechen.}
\item « Vouloir » = \all{wollen} ; \all{wir} $\rightarrow$ \all{wollen} ; « célébrer » = \all{feiern} ; « une grande fête » = accusatif : \all{ein großes Fest}. Solution : \all{Wir wollen ein großes Fest feiern.}
\end{enumerate}
\textbf{Conclusion} : dans chaque phrase, le modal est conjugué en position 2 et l'infinitif est renvoyé à la fin : c'est la structure cadre (\all{Satzklammer}).
\end{exoresolu}

\begin{methode}[Rédiger un court paragraphe : le mariage chez moi]
Pour une production guidée de 6 à 8 phrases (expression écrite ou orale) :
\begin{enumerate}
\item \textbf{Planifie} : introduction (le thème), 3 ou 4 idées (fiançailles, dot, cérémonie civile, fête), conclusion (ton opinion).
\item \textbf{Utilise des connecteurs} : \cle{zuerst} (d'abord), \cle{dann} (ensuite), \cle{danach} (après), \cle{schließlich} (finalement), \cle{meiner Meinung nach} (à mon avis).
\item \textbf{Intègre au moins deux modaux} et un datif correct pour montrer tes connaissances.
\item \textbf{Varie les débuts de phrase} mais garde le verbe en deuxième position : \all{Nach der Verlobung bringt der Bräutigam...}
\item \textbf{Relis} : majuscules aux noms, cas après prépositions, infinitif en fin avec les modaux.
\end{enumerate}
\end{methode}

\begin{exoresolu}[Production écrite guidée]
Sujet : \all{Beschreiben Sie eine Hochzeit in Ihrem Land (6--8 Sätze).} « Décrivez un mariage dans votre pays (6 à 8 phrases). »
\tcblower
\textbf{Raisonnement} : on suit le plan de la méthode : introduction, étapes avec connecteurs, deux modaux, opinion finale.

\textbf{Production modèle} :

\all{In Kamerun ist eine Hochzeit ein sehr großes Fest. Zuerst findet die Verlobung mit den beiden Familien statt. Dann muss der Bräutigam eine Mitgift für die Familie der Braut bringen. Danach gehen die jungen Leute zum Standesamt, denn die standesamtliche Trauung ist vor dem Gesetz gültig. Viele Paare wollen auch kirchlich heiraten. Bei der Feier essen und tanzen alle Gäste, und sie gratulieren dem Brautpaar. Meiner Meinung nach ist die traditionelle Hochzeit sehr schön, weil die ganze Familie zusammenkommt.}

« Au Cameroun, un mariage est une très grande fête. D'abord ont lieu les fiançailles avec les deux familles. Ensuite, le fiancé doit apporter une dot pour la famille de la fiancée. Après, les jeunes gens vont à la mairie, car la cérémonie civile est valable devant la loi. Beaucoup de couples veulent aussi se marier à l'église. À la fête, tous les invités mangent et dansent, et ils félicitent les mariés. À mon avis, le mariage traditionnel est très beau, parce que toute la famille se réunit. »

\textbf{Justifications grammaticales} : \all{muss ... bringen} et \all{wollen ... heiraten} (modaux + infinitif final) ; \all{dem Brautpaar} (datif après \all{gratulieren}) ; \all{zum Standesamt} (datif après \all{zu}) ; \all{..., weil die ganze Familie zusammenkommt} (verbe à la fin dans la subordonnée, particule \all{zusammen} rejointe au verbe en fin).
\end{exoresolu}

\begin{attention}[Les 5 erreurs qui coûtent des points au Bac]
\begin{enumerate}
\item \textbf{Oublier l'infinitif final avec un modal} : ne dis jamais \all{Er muss bringen eine Mitgift} ; dis \all{Er muss eine Mitgift bringen.} Le modal conjugué en position 2, l'infinitif à la toute fin.
\item \textbf{Confondre accusatif et datif} : \all{Ich danke die Eltern} est faux ; \all{danken} exige le datif : \all{Ich danke den Eltern.} Et au pluriel datif, n'oublie pas le -n : \all{den Gästen}, \all{mit den Kindern}.
\item \textbf{Traduire « se marier avec » par \all{heiraten mit}} : faux ami de construction ! \all{heiraten} prend l'accusatif direct : \all{Sie heiratet ihren Freund.} « Avec » n'apparaît qu'avec \all{sich verheiraten mit} ou \all{verheiratet sein mit}.
\item \textbf{Oublier la majuscule et l'article des noms} : \all{hochzeit}, \all{ehemann} sans majuscule sont des fautes. Écris toujours \all{die Hochzeit}, \all{der Ehemann}, \all{die Mitgift}.
\item \textbf{Mettre le verbe en troisième position} après un complément en tête : \all{Nach der Feier die Gäste gehen...} est faux ; dis \all{Nach der Feier gehen die Gäste ins Restaurant.} Le verbe conjugué reste toujours en deuxième position.
\end{enumerate}
\end{attention}

\newpage

\coursec{Exercices}

\courssub{Je vérifie mes connaissances}

\begin{exercice}[QCM : le mariage en Allemagne]
Entoure la bonne réponse.

\begin{enumerate}
\item En Allemagne, le mariage civil a lieu...
\begin{enumerate}
\item in der Kirche
\item im Standesamt
\item auf dem Markt
\end{enumerate}
\item \all{die Mitgift} signifie...
\begin{enumerate}
\item la dot
\item la bague
\item le témoin
\end{enumerate}
\item \all{die Verlobung} signifie...
\begin{enumerate}
\item le divorce
\item les fiançailles
\item la lune de miel
\end{enumerate}
\item \all{der Bräutigam} est...
\begin{enumerate}
\item le père de la mariée
\item le témoin
\item le futur marié
\end{enumerate}
\end{enumerate}
\end{exercice}

\begin{exercice}[Vrai ou faux ? Justifie ta réponse]
Réponds par \emph{vrai} ou \emph{faux} et justifie chaque réponse par une phrase complète en allemand.

\begin{enumerate}
\item \all{In Deutschland kann man nur in der Kirche heiraten.}
\item \all{In Kamerun gibt es oft eine traditionelle Hochzeit und eine standesamtliche Hochzeit.}
\item \all{Der Ehemann ist die Frau in der Ehe.}
\item \all{Die Mitgift ist ein Geschenk an die Familie der Braut.}
\end{enumerate}
\end{exercice}

\begin{exercice}[Vocabulaire : la famille des mots]
Complète le tableau avec le mot allemand correct (avec son article) : \all{die Hochzeit} -- \all{heiraten} -- \all{die Ehe} -- \all{der Ehemann} -- \all{die Ehefrau}.

\begin{center}
\begin{tabularx}{\linewidth}{|X|X|}
\hline
\rowcolor{popblueL} Français & Deutsch \\
\hline
le mariage (la fête) & \ldots\ldots\ldots\ldots \\
\hline
se marier (verbe) & \ldots\ldots\ldots\ldots \\
\hline
le mariage (l'union) & \ldots\ldots\ldots\ldots \\
\hline
le mari & \ldots\ldots\ldots\ldots \\
\hline
la femme (épouse) & \ldots\ldots\ldots\ldots \\
\hline
\end{tabularx}
\end{center}
\end{exercice}

\begin{exercice}[Conjugaison : les modaux au présent]
Complète avec la forme correcte du verbe modal au présent.

\begin{enumerate}
\item Ich \ldots\ldots\ldots\ldots (wollen) nächstes Jahr heiraten.
\item Du \ldots\ldots\ldots\ldots (können) die Familie der Braut besuchen.
\item Er \ldots\ldots\ldots\ldots (müssen) zum Standesamt gehen.
\item Wir \ldots\ldots\ldots\ldots (dürfen) nicht zu spät kommen.
\item Ihr \ldots\ldots\ldots\ldots (sollen) ein Geschenk kaufen.
\item Sie (pluriel) \ldots\ldots\ldots\ldots (mögen) die traditionelle Hochzeit.
\end{enumerate}
\end{exercice}

\courssub{Je m'entraîne}

\begin{exercice}[Thème : phrases à traduire]
Traduis les phrases suivantes en allemand.

\begin{enumerate}
\item Le marié offre une bague à la mariée.
\item Nous voulons nous marier en juin.
\item La famille peut donner la dot aux parents de la fiancée.
\item En Allemagne, on doit se marier à l'état civil.
\end{enumerate}
\end{exercice}

\begin{exercice}[Réécriture avec un modal imposé]
Transforme chaque phrase en y intégrant le verbe modal entre parenthèses. Attention à la place du verbe conjugué et de l'infinitif.

\begin{enumerate}
\item \all{Das Paar heiratet im Juni.} (wollen)
\item \all{Der Bräutigam kauft einen Ring.} (müssen)
\item \all{Die Gäste kommen zur Kirche.} (dürfen)
\item \all{Wir besuchen die Familie in Douala.} (können)
\item \all{Ihr lernt die Traditionen kennen.} (sollen)
\end{enumerate}
\end{exercice}

\begin{exercice}[Texte à trous : accusatif ou datif ?]
Complète les terminaisons des articles (accusatif ou datif).

\begin{textebox}[Die Hochzeit von Anna und Paul]
Anna und Paul heiraten am Samstag. Der Bräutigam schenkt d\ldots\ldots Braut ein\ldots\ldots Ring. Die Eltern geben d\ldots\ldots jungen Paar ein\ldots\ldots Geschenk. Anna dankt ihr\ldots\ldots Mutter für d\ldots\ldots Hilfe. Am Abend tanzen die Gäste mit d\ldots\ldots Ehemann und d\ldots\ldots Ehefrau.
\end{textebox}
\end{exercice}

\begin{exercice}[Appariements : mots et traductions]
Associe chaque mot allemand à sa traduction française, puis réemploie cinq de ces mots dans des phrases personnelles complètes.

\begin{center}
\begin{tabularx}{\linewidth}{|X|X|}
\hline
\rowcolor{popblueL} Deutsch & Français \\
\hline
1. die Verlobung & a. la cérémonie de mariage \\
\hline
2. die Mitgift & b. le futur marié \\
\hline
3. das Standesamt & c. les fiançailles \\
\hline
4. die Trauung & d. l'état civil \\
\hline
5. der Bräutigam & e. la dot \\
\hline
\end{tabularx}
\end{center}
\end{exercice}

\courssub{Je me prépare au Bac}

\begin{exercice}[Type Bac — Compréhension de l'écrit]
Lis le texte suivant, puis réponds aux questions \textbf{en allemand, par des phrases complètes}.

\begin{textebox}[Heiraten in der Schweiz]
In der Schweiz muss jedes Paar zuerst standesamtlich heiraten. Das ist das Gesetz: Nur die Hochzeit im Zivilstandsamt ist offiziell. Danach können viele Paare auch in der Kirche heiraten, wenn sie das wollen.

Die Schweizer Hochzeit ist oft klein und familiär. Das Paar lädt die Familie und die besten Freunde ein. Nach der Trauung gibt es ein Festessen, und die Gäste schenken dem Paar Geld oder Geschenke für das neue Haus.

In der Schweiz gibt es keine Mitgift wie in manchen afrikanischen Ländern. Aber die Familien helfen dem jungen Paar oft: Die Eltern können Geld geben, und die Großeltern schenken manchmal Möbel. Viele junge Paare wollen zuerst arbeiten und Geld sparen, bevor sie heiraten. Deshalb heiraten die Schweizer oft spät, mit 30 Jahren oder später.
\end{textebox}

\textbf{Fragen zum Text :}

\begin{enumerate}
\item Wo muss jedes Paar in der Schweiz heiraten ?
\item Kann man in der Schweiz auch in der Kirche heiraten ?
\item Wen lädt das Paar zur Hochzeit ein ?
\item Was schenken die Gäste dem Paar ?
\item Gibt es in der Schweiz eine Mitgift ?
\item Warum heiraten viele Schweizer spät ?
\end{enumerate}

\textbf{Fragen zur Sprache :}

\begin{enumerate}
\item Relève dans le texte deux verbes de modalité et donne leur infinitif.
\item Relève dans le texte un complément au datif et un complément à l'accusatif.
\item Conjugue \all{müssen} à la première personne du singulier et à la troisième personne du pluriel.
\end{enumerate}
\end{exercice}

\begin{exercice}[Type Bac — Thème et version]
\textbf{A. Version.} Traduis le texte suivant en français.

\begin{textebox}[Eine Hochzeit in Yaoundé]
Meine Schwester heiratet nächsten Samstag in Yaoundé. Zuerst gibt es die traditionelle Hochzeit: Der Bräutigam muss die Mitgift an die Familie bringen. Dann heiratet das Paar standesamtlich, und am Abend gibt es ein großes Fest mit Musik und Tanz. Alle Gäste wollen dem jungen Paar gratulieren.
\end{textebox}

\textbf{B. Thème.} Traduis les phrases suivantes en allemand.

\begin{enumerate}
\item Le marié doit offrir une bague à la fiancée. Il peut aussi offrir un cadeau à la famille.
\item Nous voulons inviter nos amis à la fête.
\item Au Cameroun, beaucoup de couples célèbrent deux mariages : le mariage traditionnel et le mariage civil.
\end{enumerate}
\end{exercice}

\begin{exercice}[Type Bac — Production écrite]
Rédige en allemand un paragraphe de \textbf{8 à 10 lignes} (environ 80 mots) sur le sujet suivant :

\begin{center}
\all{„Meine Meinung über die Mitgift“}
\end{center}

\textbf{Consignes obligatoires :}
\begin{itemize}
\item utilise \textbf{au moins trois verbes de modalité} différents au présent ;
\item utilise les connecteurs \all{zuerst}, \all{dann}, \all{schließlich} ;
\item utilise au moins une fois le datif et une fois l'accusatif avec un article ;
\item donne ton opinion personnelle (pour ou contre la dot) avec un argument.
\end{itemize}

\textbf{Aide au vocabulaire :} \all{die Meinung} (l'opinion), \all{wichtig} (important), \all{die Tradition} (la tradition), \all{respektieren} (respecter), \all{das Geld} (l'argent), \all{helfen} + datif (aider).
\end{exercice}

\newpage

\coursec{Corrigés des exercices}

\begin{corrige}[Exercice 1 — QCM : le mariage en Allemagne]
\begin{enumerate}
\item \textbf{b) im Standesamt} — En Allemagne, seul le mariage civil (à l'état civil) a une valeur légale ; la cérémonie religieuse est facultative.
\item \textbf{a) la dot} — \all{die Mitgift} est le don (argent, biens) versé selon la tradition ; attention au faux ami : ce n'est pas un « cadeau » (\all{das Geschenk}).
\item \textbf{b) les fiançailles} — \all{die Verlobung} ; famille de mots : \all{sich verloben} (se fiancer), \all{der Verlobte} (le fiancé).
\item \textbf{c) le futur marié} — \all{der Bräutigam} ; ne pas confondre avec \all{die Braut} (la mariée) et \all{der Ehemann} (le mari, déjà marié).
\end{enumerate}
\end{corrige}

\begin{corrige}[Exercice 2 — Vrai ou faux ?]
\begin{enumerate}
\item \textbf{Faux.} \all{In Deutschland kann man auch standesamtlich heiraten; nur das Standesamt ist offiziell.} « En Allemagne, on peut aussi se marier à l'état civil ; seul l'état civil est officiel. »
\item \textbf{Vrai.} \all{In Kamerun gibt es oft zwei Hochzeiten: eine traditionelle und eine standesamtliche.} « Au Cameroun, il y a souvent deux mariages : un traditionnel et un civil. »
\item \textbf{Faux.} \all{Der Ehemann ist der Mann in der Ehe; die Ehefrau ist die Frau.} « Le mari est l'homme dans le mariage ; l'épouse est la femme. »
\item \textbf{Vrai.} \all{Die Mitgift ist ein Geschenk des Bräutigams an die Familie der Braut.} « La dot est un don du marié à la famille de la mariée. »
\end{enumerate}
\emph{Méthode :} la justification doit reprendre un élément du texte support ; une phrase complète = sujet + verbe conjugué en 2\textsuperscript{e} position.
\end{corrige}

\begin{corrige}[Exercice 3 — Vocabulaire : la famille des mots]
\begin{center}
\begin{tabularx}{\linewidth}{|X|X|}
\hline
\rowcolor{popblueL} Français & Deutsch \\
\hline
le mariage (la fête) & \all{die Hochzeit} \\
\hline
se marier (verbe) & \all{heiraten} \\
\hline
le mariage (l'union) & \all{die Ehe} \\
\hline
le mari & \all{der Ehemann} \\
\hline
la femme (épouse) & \all{die Ehefrau} \\
\hline
\end{tabularx}
\end{center}
\emph{Distinction importante :} \all{die Hochzeit} désigne la fête, la cérémonie ; \all{die Ehe} désigne l'union, l'état matrimonial. Ne pas confondre les deux en production écrite.
\end{corrige}

\begin{corrige}[Exercice 4 — Conjugaison : les modaux au présent]
\begin{enumerate}
\item Ich \textbf{will} nächstes Jahr heiraten. (wollen : ich will)
\item Du \textbf{kannst} die Familie der Braut besuchen. (können : du kannst)
\item Er \textbf{muss} zum Standesamt gehen. (müssen : er muss, sans Umlaut)
\item Wir \textbf{dürfen} nicht zu spät kommen. (dürfen : wir dürfen, forme de l'infinitif)
\item Ihr \textbf{sollt} ein Geschenk kaufen. (sollen : ihr sollt)
\item Sie \textbf{mögen} die traditionelle Hochzeit. (mögen : sie mögen)
\end{enumerate}
\emph{Rappel :} aux 1\textsuperscript{re} et 3\textsuperscript{e} personnes du singulier, le modal n'a pas de terminaison (\all{ich muss, er muss}) ; les Umlauts de l'infinitif disparaissent au singulier (\all{müssen} $\rightarrow$ \all{muss}, \all{können} $\rightarrow$ \all{kann}).
\end{corrige}

\begin{corrige}[Exercice 5 — Thème : phrases à traduire]
\begin{enumerate}
\item \all{Der Bräutigam schenkt der Braut einen Ring.} — \all{schenken} + datif de la personne (\all{der Braut}) + accusatif de la chose (\all{einen Ring}).
\item \all{Wir wollen im Juni heiraten.} — modal conjugué en 2\textsuperscript{e} position, infinitif en fin de phrase ; « en juin » = \all{im Juni}.
\item \all{Die Familie kann den Eltern der Braut die Mitgift geben.} — \all{geben} + datif (\all{den Eltern}) + accusatif (\all{die Mitgift}).
\item \all{In Deutschland muss man standesamtlich heiraten.} — « on » = \all{man} ; \all{muss} en 2\textsuperscript{e} position, \all{heiraten} à la fin.
\end{enumerate}
\emph{Points de vigilance :} majuscule à tous les noms ; place de l'infinitif en fin de phrase avec un modal ; ordre datif avant accusatif quand les deux sont des noms.
\end{corrige}

\begin{corrige}[Exercice 6 — Réécriture avec un modal imposé]
\begin{enumerate}
\item \all{Das Paar \textbf{will} im Juni \textbf{heiraten}.}
\item \all{Der Bräutigam \textbf{muss} einen Ring \textbf{kaufen}.}
\item \all{Die Gäste \textbf{dürfen} zur Kirche \textbf{kommen}.}
\item \all{Wir \textbf{können} die Familie in Douala \textbf{besuchen}.}
\item \all{Ihr \textbf{sollt} die Traditionen kennen\textbf{lernen}.}
\end{enumerate}
\emph{Règle appliquée :} le modal prend la place du verbe conjugué (2\textsuperscript{e} position) et l'ancien verbe conjugué passe à l'infinitif, en fin de phrase. Le sens change légèrement : \all{will heiraten} = « veut se marier ».
\emph{Attention :} \all{kennenlernen} est un verbe à particule séparable ; à l'infinitif avec modal, il reste en un seul mot à la fin.
\end{corrige}

\begin{corrige}[Exercice 7 — Texte à trous : accusatif ou datif ?]
Texte complété :

\all{Anna und Paul heiraten am Samstag. Der Bräutigam schenkt d\textbf{er} Braut ein\textbf{en} Ring. Die Eltern geben d\textbf{em} jungen Paar ein Geschenk. Anna dankt ihr\textbf{er} Mutter für d\textbf{ie} Hilfe. Am Abend tanzen die Gäste mit d\textbf{em} Ehemann und d\textbf{er} Ehefrau.}

Justifications :
\begin{itemize}
\item \all{der Braut} : datif féminin, personne qui reçoit (\all{schenken} + dat.) ; \all{einen Ring} : accusatif masculin, chose donnée.
\item \all{dem jungen Paar} : datif neutre (\all{das Paar} $\rightarrow$ \all{dem Paar}) ; \all{ein Geschenk} : accusatif neutre (identique au nominatif).
\item \all{ihrer Mutter} : datif féminin après \all{danken} (verbe toujours suivi du datif) ; \all{für die Hilfe} : accusatif après la préposition \all{für}.
\item \all{mit dem Ehemann} (masculin) et \all{mit der Ehefrau} (féminin) : datif après la préposition \all{mit}.
\end{itemize}
\end{corrige}

\begin{corrige}[Exercice 8 — Appariements]
Appariements : \textbf{1-c, 2-e, 3-d, 4-a, 5-b}.

Exemples de phrases personnelles (modèles) :
\begin{enumerate}
\item \all{Die Verlobung von meinem Bruder ist im Dezember.} « Les fiançailles de mon frère sont en décembre. »
\item \all{Der Bräutigam bringt die Mitgift zur Familie der Braut.} « Le marié apporte la dot à la famille de la mariée. »
\item \all{Wir gehen morgen zum Standesamt.} « Nous allons demain à l'état civil. »
\item \all{Die Trauung beginnt um zehn Uhr.} « La cérémonie commence à dix heures. »
\item \all{Der Bräutigam trägt einen schönen Anzug.} « Le marié porte un beau costume. »
\end{enumerate}
\emph{Vérifier :} article correct, verbe en 2\textsuperscript{e} position, majuscule aux noms.
\end{corrige}

\begin{corrige}[Exercice 9 — Type Bac : Compréhension de l'écrit]
\textbf{Fragen zum Text} (réponses modèles) :
\begin{enumerate}
\item \all{Jedes Paar muss in der Schweiz zuerst im Zivilstandsamt heiraten.}
\item \all{Ja, danach können viele Paare auch in der Kirche heiraten, wenn sie das wollen.}
\item \all{Das Paar lädt die Familie und die besten Freunde ein.}
\item \all{Die Gäste schenken dem Paar Geld oder Geschenke für das neue Haus.}
\item \all{Nein, in der Schweiz gibt es keine Mitgift.}
\item \all{Viele Schweizer heiraten spät, weil sie zuerst arbeiten und Geld sparen wollen.}
\end{enumerate}
\emph{Méthode :} reprendre les mots de la question dans la réponse, ne jamais répondre par un seul mot.

\textbf{Fragen zur Sprache :}
\begin{enumerate}
\item Modaux du texte : \all{muss} (infinitif \all{müssen}), \all{können} (infinitif \all{können}) ; aussi \all{wollen} (\all{wollen}), \all{wollen} dans \all{wollen zuerst arbeiten}.
\item Datif : \all{dem Paar} / \all{dem jungen Paar} ; accusatif : \all{die Familie} (après \all{einlädt}) ou \all{Geld} / \all{Möbel} (compléments de \all{geben}, \all{schenken}).
\item \all{müssen} : ich \textbf{muss} ; sie (pluriel) \textbf{müssen}.
\end{enumerate}
\emph{Barème indicatif (sur 20) :} questions sur le texte 6 $\times$ 2 = 12 pts ; questions de langue 3 $\times$ 2 = 6 pts ; qualité de la langue 2 pts.
\emph{Points de vigilance :} phrase complète obligatoire ; verbe en 2\textsuperscript{e} position ; \all{weil} envoie le verbe à la fin (\all{weil sie ... sparen wollen}).
\end{corrige}

\begin{corrige}[Exercice 10 — Type Bac : Thème et version]
\textbf{A. Version (traduction modèle) :}

« Ma sœur se marie samedi prochain à Yaoundé. D'abord, il y a le mariage traditionnel : le marié doit apporter la dot à la famille. Ensuite, le couple se marie à l'état civil, et le soir il y a une grande fête avec de la musique et de la danse. Tous les invités veulent féliciter le jeune couple. »

\emph{Points de vigilance :} \all{gratulieren} + datif = « féliciter » ; \all{an die Familie bringen} = « apporter à la famille » ; ne pas traduire \all{das Fest} par « le festival ».

\textbf{B. Thème (traductions modèles) :}
\begin{enumerate}
\item \all{Der Bräutigam muss der Braut einen Ring schenken. Er kann der Familie auch ein Geschenk schenken.}
\item \all{Wir wollen unsere Freunde zum Fest einladen.}
\item \all{In Kamerun feiern viele Paare zwei Hochzeiten: die traditionelle Hochzeit und die standesamtliche Hochzeit.}
\end{enumerate}
\emph{Barème indicatif (sur 20) :} version 10 pts ; thème 3 $\times$ 3 = 9 pts ; correction globale 1 pt.
\emph{Points de vigilance :} \all{einladen} est séparable (\all{lädt ... ein}) ; datif de la personne avant accusatif de la chose ; \all{zum Fest} = \all{zu dem Fest}.
\end{corrige}

\begin{corrige}[Exercice 11 — Type Bac : Production écrite]
\textbf{Proposition de rédaction modèle :}

\begin{textebox}[Meine Meinung über die Mitgift]
Die Mitgift ist eine alte Tradition in Kamerun. Zuerst muss ich sagen: Ich finde diese Tradition wichtig, denn sie respektiert die Familie der Braut. Dann kann die Mitgift dem jungen Paar auch helfen: Die Familie gibt oft Geld oder Geschenke für das neue Leben. Aber die Mitgift darf nicht zu teuer sein. Manche Männer können nicht heiraten, weil sie kein Geld haben, und das ist ein Problem. Schließlich ist meine Meinung: Wir sollen die Tradition respektieren, aber die Familien müssen auch an die Liebe denken. Die Mitgift soll ein Symbol sein, kein Geschäft.
\end{textebox}

(Environ 90 mots.)

\textbf{Vérification des consignes :}
\begin{itemize}
\item Trois modaux différents : \all{muss}, \all{kann}, \all{darf}, \all{können}, \all{sollen}, \all{müssen}, \all{soll}.
\item Connecteurs : \all{zuerst}, \all{dann}, \all{schließlich}.
\item Datif : \all{dem jungen Paar} ; accusatif : \all{die Familie der Braut}, \all{die Mitgift}.
\item Opinion personnelle avec argument : pour la tradition, mais contre les abus.
\end{itemize}
\textbf{Barème indicatif (sur 20) :} respect des consignes 6 pts ; grammaire (modaux, cas, place du verbe) 6 pts ; vocabulaire thématique 4 pts ; cohérence et connecteurs 4 pts.
\textbf{Points de vigilance :} après \all{denn} le verbe reste en 2\textsuperscript{e} position ; après \all{weil} il va à la fin ; \all{helfen} se construit avec le datif (\all{dem Paar helfen}) ; ne pas écrire \all{die Mitgift ist ein Geschenk für die Braut} — c'est un don à la \emph{famille} de la mariée.
\end{corrige}

\newpage

\coursec{Fiche bilan}

\begin{popfigure}
\begin{tikzpicture}[mindmap, grow cyclic, every node/.style={concept, text=white, font=\small}, concept color=popblue,
level 1/.append style={level distance=3.6cm, sibling angle=51.4},
level 2/.append style={level distance=2.2cm, sibling angle=40}]
\node[concept, font=\small\bfseries] {Les types de mariages}
[clockwise from=90]
child[concept color=poporange] { node {Texte : Heiraten in Deutschland und in Kamerun}
}
child[concept color=popgreen] { node {Lexique : die Hochzeit, die Ehe}
  child { node[concept color=popgreen!70] {die Verlobung, die Mitgift} }
  child { node[concept color=popgreen!70] {standesamtlich, kirchlich, traditionell} }
}
child[concept color=poppurple] { node {Akkusativ}
  child { node[concept color=poppurple!70] {den, die, das, die} }
}
child[concept color=popblue!80] { node {Dativ}
  child { node[concept color=popblue!50] {dem, der, dem, den + n} }
}
child[concept color=popgold] { node {Verbes de modalité}
  child { node[concept color=popgold!70] {können, müssen, dürfen, wollen, sollen, mögen} }
}
child[concept color=poppink] { node {Méthode : comprendre, résumer, produire}
};
\end{tikzpicture}
\legende{Carte mentale de la leçon : les types de mariages}
\end{popfigure}

\begin{bilan}
\textbf{1. Lexique clé (à connaître avec article et pluriel)}
\begin{itemize}
\item \all{die Hochzeit, -en} : le mariage (la fête) ; \all{die Ehe, -n} : le mariage (l'union) ; \all{heiraten (+ Akk.)} : épouser
\item \all{der Ehemann, -männer} : le mari ; \all{die Ehefrau, -en} : la femme ; \all{das Ehepaar, -e} : le couple marié
\item \all{die Verlobung, -en} : les fiançailles ; \all{sich verloben (mit + Dat.)} : se fiancer ; \all{der Verlobte, -n / die Verlobte, -n} : le fiancé / la fiancée
\item \all{die Mitgift, -en} : la dot ; \all{der Bräutigam, -e} : le marié ; \all{die Braut, -äute} : la mariée
\item \all{die standesamtliche Hochzeit} : le mariage civil ; \all{die kirchliche Hochzeit} : le mariage religieux ; \all{die traditionelle Hochzeit} : le mariage traditionnel
\item \all{der Standesbeamte, -n} : l'officier d'état civil ; \all{der Trauring, -e} : l'alliance ; \all{die Einladung, -en} : l'invitation
\end{itemize}

\textbf{2. Grammaire à connaître par cœur}
\begin{center}
\begin{tabularx}{\linewidth}{|l|X|X|X|}
\hline
\rowcolor{popblueL} Cas & Masculin & Féminin / Neutre & Pluriel \\
\hline
Nominatif & der / ein & die / eine ; das / ein & die / --- \\
\hline
Akkusativ & \textbf{den / einen} & die / eine ; das / ein & die / --- \\
\hline
Dativ & \textbf{dem / einem} & \textbf{der / einer} ; dem / einem & \textbf{den + n / --- + n} \\
\hline
\end{tabularx}
\end{center}
\begin{itemize}
\item \cle{Akkusativ} : complément d'objet direct et après \all{durch, für, gegen, ohne, um}. Seul le \textbf{masculin singulier} change de forme. \all{Er heiratet die Tochter des Lehrers.} « Il épouse la fille du professeur. »
\item \cle{Dativ} : complément d'objet indirect et après \all{mit, zu, von, bei, nach, aus}. Au pluriel, le nom prend un \textbf{-n} final (sauf s'il se termine déjà par -n ou -s). \all{Sie spricht mit dem Bräutigam.} « Elle parle avec le marié. »
\item \cle{Verbes de modalité au présent} : le modal est conjugué en 2\up{e} position, l'infinitif va \textbf{à la fin} de la phrase.
\end{itemize}
\begin{center}
\begin{tabularx}{\linewidth}{|l|X|X|X|}
\hline
\rowcolor{popgreenL} Modal & ich / er & wir / sie & Sens \\
\hline
können & kann & können & pouvoir \\
\hline
müssen & muss & müssen & devoir (obligation) \\
\hline
dürfen & darf & dürfen & avoir le droit de \\
\hline
wollen & will & wollen & vouloir \\
\hline
sollen & soll & sollen & devoir (conseil, ordre) \\
\hline
mögen (möchte) & mag (möchte) & mögen (möchten) & aimer (vouloir poliment) \\
\hline
\end{tabularx}
\end{center}
\all{In Deutschland muss man standesamtlich heiraten.} « En Allemagne, on doit se marier à l'état civil. »

\textbf{3. Méthodes express}
\begin{itemize}
\item \cle{Comprendre un texte} : lire le titre et la source, repérer les mots connus, lecture globale puis détaillée, souligner les mots-clés du thème.
\item \cle{Répondre aux questions} : reprendre les mots de la question, répondre par une phrase complète, verbe en 2\up{e} position.
\item \cle{Résumer} : 3 à 5 phrases, ses propres mots, uniquement les idées principales.
\end{itemize}
\end{bilan}

\begin{attention}[Pièges fréquents des francophones]
\begin{itemize}
\item \all{die Mitgift} signifie « la dot » (versée lors du mariage traditionnel), et non « le cadeau » : le cadeau se dit \all{das Geschenk, -e}.
\item \all{heiraten} se construit \textbf{sans préposition} : \all{Er heiratet seine Freundin.} « Il épouse sa petite amie. » Ne dites jamais \all{heiraten mit}.
\item En revanche : \all{sich verloben mit + Dat.} « se fiancer avec » : \all{Sie verlobt sich mit einem Lehrer.} « Elle se fiance avec un professeur. »
\item \all{möchten} est la forme polie de \all{mögen} (Konjunktiv II) : \all{Ich möchte heiraten.} « Je voudrais me marier. »
\item Après un modal, l'infinitif se place à la fin : \all{Man muss die Mitgift bezahlen.} « On doit payer la dot. »
\item Majuscule obligatoire à tous les noms : \all{die Hochzeit, der Ehemann, die Verlobung}.
\end{itemize}
\end{attention}

\begin{exemplebox}[Les trois types de mariages en une phrase chacun]
\begin{itemize}
\item \all{Das Paar heiratet zuerst standesamtlich.} « Le couple se marie d'abord à l'état civil. »
\item \all{Am Samstag feiert die Familie die kirchliche Hochzeit in der Kirche.} « Samedi, la famille fête le mariage religieux à l'église. »
\item \all{In Kamerun bezahlt der Bräutigam oft die Mitgift an die Familie der Braut.} « Au Cameroun, le marié paie souvent la dot à la famille de la mariée. »
\end{itemize}
\end{exemplebox}

\begin{aretenir}[Checklist avant le Bac]
\begin{itemize}
\item[$\square$] Je connais le lexique du mariage avec l'article et le pluriel de chaque nom.
\item[$\square$] Je distingue \all{die Hochzeit} (la fête) et \all{die Ehe} (l'union).
\item[$\square$] Je connais les trois types de mariages : \all{standesamtlich, kirchlich, traditionell}.
\item[$\square$] Je sais que seul le masculin change à l'accusatif : \all{der} devient \all{den}.
\item[$\square$] Je connais les formes du datif : \all{dem, der, dem, den + n}.
\item[$\square$] Je cite les prépositions suivies de l'accusatif (\all{für, ohne, um...}) et du datif (\all{mit, zu, von...}).
\item[$\square$] Je conjugue sans erreur les six verbes de modalité au présent.
\item[$\square$] Je place le modal en 2\up{e} position et l'infinitif à la fin de la phrase.
\item[$\square$] Je mets une majuscule à tous les noms allemands (\all{die Hochzeit, der Ehemann}).
\item[$\square$] Je sais répondre à une question de compréhension par une phrase complète.
\item[$\square$] Je sais rédiger un court résumé (3 à 5 phrases) sur les types de mariages.
\item[$\square$] Je fais attention aux faux amis : \all{die Mitgift} = la dot, pas « le cadeau ».
\end{itemize}
\end{aretenir}

\findecours

\end{document}
